\documentclass[12pt]{article}
%last edited by wdj 2023-03-26
\usepackage{float}
\usepackage{amsthm}
\usepackage{amsmath}
\usepackage{graphicx}
\usepackage{amsfonts}
\usepackage{url}
\usepackage{fancyvrb}
\usepackage{makeidx}
\usepackage{cprotect}

\newtheorem{theorem}{Theorem}
\newtheorem{corollary}[theorem]{Corollary}
\newtheorem{conjecture}[theorem]{Conjecture}
\newtheorem{lemma}[theorem]{Lemma}
\newtheorem{proposition}[theorem]{Proposition}
\newtheorem{definition}[theorem]{Definition}
\newtheorem{example}[theorem]{Example}
\newtheorem{axiom}{Axiom}
\newtheorem{remark}{Remark}
\newtheorem{exercise}{Exercise}[section]

\def\ppp{{\mathbb{P}}}
\def\hhh{{\mathbb{H}}}
\def\aaa{{\mathbb{A}}}
\def\fff{{\mathbb{F}}}
\def\qqq{\mathbb{Q}}
\def\rrr{\mathbb{R}}
\def\ccc{\mathbb{C}}
\def\zzz{\mathbb{Z}}
\def\nnn{\mathbb{N}}
\def\pf{{\bf proof}:\ }
\def\qed{$\Box$}

\makeindex

\begin{document}


\author{David Joyner and Caroline Melles}
\title{Harmonic quotient morphisms of regular graphs: examples}
\date{2023-03-26}
\maketitle

In this paper all graphs are simple and connected.


This paper has examples of
harmonic morphisms $\phi:\Gamma_2\to \Gamma_1$,
where $\Gamma_2$ is a $k$-regular graph ($k=2,3,4,5$)
with a ``small'' number of vertices,
and $\Gamma_1=\Gamma_2/G$ is quotient graph by a subgroup
$G\subset Aut(\Gamma_2)$.

\vskip .2in

\noindent
{\bf Question}: Which regular graphs $\Gamma_2$ have no (or very few)
harmonic quotients other than $P_2$?

For example, look at $K_5$ and {\tt d5reg10a} below.

\vskip .2in

\tableofcontents

\vskip .4in

\section{Definitions}

We will always assume we have fixed an indexing of the vertices $V$ and
edges $E$ of a graph $\Gamma = (V,E)$.
Typically, the vertices are represented as $V = \{0, 1, \dots , n - 1\}$
and the edges are represented by pairs $(u, v)$, where $u < v$
are the vertices in $V$ incident to that edge.
In this situation we say $\Gamma$ is in {\it standard form}.
\index{standard form}

Recall a graph $\Gamma=(V,E)$ with automorphism group $G$ is
called {\it vertex transitive} if $G$ acts transitively on $V$.
It is called {\it edge transitive} if $G$ acts transitively on $E$.
We distinguish between (unoriented) edges and arcs, which are
oriented with a head and a tail. A graph is {\it arc transitive},
or {\it symmetric}, if $G$ acts transitively on the set of arcs
of $\Gamma$.
\index{vertex transitive graph}
\index{edge transitive graph}
\index{arc transitive graph}
Recall a
graph is called {\it Hamiltonian} if there is a simple closed path
containing each vertex of the graph.
\index{Hamiltonian graph}

A {\it circulant graph} with parameters $(n,L)$ is a
graph $\Gamma=(V,E)$ in standard form\footnote{We identify the
vertices $V=\{0,1,\dots,n-1\}$ with the cyclic group
$\zzz/n\zzz$ and $L\subset V$.} with the property
that each vertex $i$ is adjacent to
the vertices $i+j$ and $i-j$ for each $j \in L$.
\index{circulant graph}

Generally, the barbell graph $Barbell(n_1, n_2)$
is a graph with $2n_1 + n_2$ vertices ($n_1>2$, $n_2\geq 0$),
that consists of a path graph $P_{n_2+2}$ connecting two complete graphs
$K_{n_1}$.
\index{barbell graph}


\section{Several (new?) regular graph constructions}


These are, to the best of our knowledge, original. In any case, they are
mentioned, as an aside, more for amusement of the interested reader than anything else.


\begin{itemize}
\item[(A)]
Let $\Gamma_1$ be a $(k-1)$-regular graph with $n$ vertices. (Here $k$ is
``small'' relative to $n$.) Let $v$ be another vertex (not in $\Gamma_1$),
which we connect to $k$ vertices of $\Gamma_1$ by an edge. Call this
$\Gamma_2$. Make a copy of $\Gamma_2$. Connect each vertex of degree $k-1$
on one copy to its corresponding vertex on the other copy.
The resulting graph $\Gamma_3$ is
$k$-regular, with $2n+2$ vertices.

\item[(B)] Let $\Gamma_1$ be as before. Assume $n/k=d$ is an integer.
Let $v_1, \dots, v_d$ be new vertices (not in $\Gamma_1$).
Connect each $v_i$ to $k$ different vertices of
$\Gamma_1$ (assuming these sets are disjoint from each other). The result is
a $k$-regular graph with $n+d$ vertices.

\item[(C)] Let $\Gamma_1$ be as before. Assume $(n-r)/k=d$ is an integer
for some small $0<r<k$.
Let $v_1, \dots, v_d$ be new vertices (not in $\Gamma_1$). Connect each to $k$ different
vertices of $\Gamma_1$ (as in construction B). By hypothesis, there are $r$
remaining vertices without a new edge. Call this $\Gamma_2$.
Copy $\Gamma_2$ and connect these $r$ vertices in the first copy of $\Gamma_2$
to their mirror image in the $2^{nd}$ copy of $\Gamma_2$ (as in construction
A). The result is a $k$-regular graph with $2n+2d$ vertices. 
  
\end{itemize}


\section{Cubic cases}
\label{sec:cubic-examples}

Most of these are Hamiltonian.
Indeed, in 1880 P.G. Tait conjectured
(roughly speaking) that every $2$-connected cubic graph is Hamiltonian.
While it's known this conjecture is false, the counterexamples are exceptional
(e.g., the Petersen graph).
\index{Petersen graph}
\index{Tait's conjecture}

\subsection{$6$ vertices}
\label{sec:cubic-6-vertices}

There are two $3$-regular graphs with $6$ vertices, both of which are
vertex transitive but only one is also edge transitive.
All these are circulant, Hamiltonian graphs.


\begin{center}
\begin{tabular}{c|c|c|c|c}
  graph label & LCF                 & size aut. gp. & genus & size crit. gp. \\ \hline
  d3reg6a    & (6, [2, 3, -2], 2)   & 12            & 4    & $75=3\cdot 5^2$ \\
  d3reg6b    & (6, [3], 6)         & 72             & 4    & $81=3^4$ \\
\end{tabular}
\end{center}  

\begin{enumerate}
\item
  {\tt d3reg6a}\newline
\index{prism graph}  
Also known as the prism graph, denoted $Y_3$,
this has characteristic polynomial  
\[
x^6 - 9x^4 - 4x^3 + 12x^2,
\]
diameter $2$,
girth $3$,
circumference $6$.
This is vertex transitive but not an edge transitive graph. It is Hamiltonian and circulant
with parameters $(6, [2, 3, 4])$.

It has $360$ non-trivial harmonic morphisms to $St_4$, $36$ non-trivial harmonic
morphisms to $P_3$, $102$ to $P_4$, $12$ to $P_5$,
$6$ non-trivial harmonic morphisms\footnote{Exactly one of these is a
harmonic quotient morphism to $C_3$. That morphism is a
quasicover with constant vertical multiplicity $1$.} to $C_3$,
There are no non-trivial harmonic morphisms to $P_6$, to the
diamond graph (with $4$ vertices), nor are there any to $K_4$.

\item
  {\tt d3reg6b}\newline
\index{complete bipartite graph}  
Also known as the complete bipartite graph $K_{3,3}$,
this has characteristic polynomial  
\[
x^6 - 9x^4,
\]
diameter $2$,
girth $4$,
circumference $6$.
This is both a vertex transitive and an edge transitive graph.
Moreover, it is Hamiltonian and circulant
with parameters $(6, [1, 3, 5])$.


It has $444$ non-trivial harmonic morphisms\footnote{At least one of these is a
harmonic quotient morphism to $St_4$. That morphism is a homomorphism.}
to $St_4$, $48$ non-trivial harmonic
morphisms to $P_3$, $72$ to $P_4$, $72$ to $P_5$,
$36$ non-trivial harmonic morphisms\footnote{At least one of these is a
harmonic quotient morphism to $C_3$. That morphism is a
quasicover with constant vertical multiplicity $1$.} to $C_3$,
There are no non-trivial harmonic morphisms to $P_6$, to the
diamond graph, nor are there any to $K_4$.
\end{enumerate}

\subsection{$8$ vertices}
\label{sec:cubic-8-vertices}

There are five $3$-regular graphs with $8$ vertices, two of which are
vertex transitive but none are edge transitive.
All these are Hamiltonian graphs.



\begin{center}
\begin{tabular}{c|c|c|c|c}
  graph label & LCF & size aut. gp. & genus & size crit. gp. \\ \hline
  d3reg8a & (8, [4,-2,4,2], 2) & 16 & 5 & $256=2^8$ \\
  d3reg8b & (8, [2,2,-2,-2], 2) & 4 & 5 & $336=2^43\cdot 7$ \\
  d3reg8c & (8, [2,4,-2,3,3,4,-3,-3], 1) & 12 & 5 & $363=3\cdot 11^2$ \\
  d3reg8d & (8, [-3,3], 4) & 48 & 5 & $384=2^73$ \\ 
  d3reg8e & (8, [4,-3,3,4], 2) & 16 & 5 & $392=2^37^2$ \\
\end{tabular}
\end{center}  

\begin{enumerate}
\item
  {\tt d3reg8a}\newline
  The one with characteristic polynomial  
\[
x^8 - 12x^6 - 8x^5 + 38x^4 + 48x^3 - 12x^2 - 40x - 15.
\]
Diameter $3$,
girth $3$,
circumference $8$.
This is not a transitive graph but it is Hamiltonian.

It has no non-trivial harmonic morphisms to $St_4$, $P_4$
(and therefore none to $P_k$, $k>4$),% , by Lemma \ref{lemma:families}),
$C_3$, or to
$C_4$, but has $2$ non-trivial harmonic morphisms to $P_3$,
$8$ non-trivial harmonic morphisms to the
diamond graph (with $4$ vertices), and
has\footnote{According to SageMath, all these $48$ harmonic
  morphisms to $K_4$ are covers.}
$48$ non-trivial harmonic morphisms to $K_4$.

SageMath command:

\verb+Gamma = graphs.LCFGraph(8, [2, 2, -2, -2], 2)+

This is a matchstick graph.
% (see Example \ref{example:cubic-matchstick}).

\item
  {\tt  d3reg8b}\newline
The one with 
characteristic polynomial

\[
x^8 - 12x^6 - 4x^5 + 38x^4 + 16x^3 - 36x^2 - 12x + 9
\]
Diameter $3$,
girth $3$,
circumference $8$.
This is a Hamiltonian graph.

It has $2$ non-trivial harmonic morphisms to $P_3$ and
$4$ non-trivial harmonic morphisms to the diamond graph.
However, there are
no non-trivial harmonic morphisms to $St_4$, $P_4$
(and therefore $P_k$, $k>4$), $C_3$,
$C_4$, or to $K_4$.

SageMath command:

\verb@Gamma = graphs.LCFGraph(8, [4, -2, 4, 2], 2)@


\item
  {\tt d3reg8c}\newline
  The one with characteristic polynomial  
\[
x^8 - 12x^6 - 2x^5 + 36x^4 - 31x^2 + 12x
\]
Diameter $2$,
girth $3$,
circumference $8$.
This is a Hamiltonian graph.

It has $6$ non-trivial harmonic morphisms to $St_4$ and
$6$ non-trivial harmonic morphisms to $P_3$. However, there are
no non-trivial harmonic morphisms to the diamond graph, $P_4$,
$C_3$, $C_4$, or to $K_4$.

SageMath command:

\verb@Gamma = graphs.LCFGraph(8, [2, 4, -2, 3, 3, 4, -3, -3], 1)@


\item
  {\tt d3reg8d}\newline
The one (the cube graph, $Cube_3$) with characteristic polynomial
\index{cube graph}

\[
x^8 - 12x^6 + 30x^4 - 28x^2 + 9
\]
Diameter $3$,
girth $4$,
circumference $8$.
This is a vertex transitive, edge transitive, arc transitive,
distance regular, Hamiltonian graph.
\index{vertex transitive graph}
\index{edge transitive graph}
\index{arc transitive graph}

It has $24$ non-trivial harmonic morphisms to the diamond graph,
$24$ non-trivial harmonic morphisms to the $C_4$,
$12$ non-trivial harmonic morphisms to the $P_3$,
and\footnote{According to SageMath, all these $24$ harmonic
  morphisms to $K_4$ are covers.}
$24$ non-trivial harmonic morphisms to $K_4$.
However, there are
no non-trivial harmonic morphisms to $P_4$, $C_3$, or to $St_4$.

SageMath command:

\verb@Gamma = graphs.LCFGraph(8, [-3,3], 4)@

\item
  {\tt d3reg8e}\newline
The one (the {\it Wagner graph,} {\tt Wag})
with automorphism group the dihedral group $D_{16}$, and
characteristic polynomial
\index{Wagner graph}

\[
x^8 - 12x^6 + 34x^4 - 16x^3 - 20x^2 + 16x - 3
\]
Diameter $2$,
girth $4$,
circumference $8$.
This graph is vertex transitive, but not edge transitive,
and is a Hamiltonian graph.
\index{vertex transitive graph}

It has $8$ non-trivial harmonic morphisms to the $C_4$
and $4$ non-trivial harmonic morphisms to the $P_3$.
However, it has no non-trivial
harmonic morphisms to $C_3$, $St_4$, $P_4$, or to $K_4$.

SageMath command:

\verb@Gamma = graphs.LCFGraph(8, [4, -3, 3, 4], 2)@


\end{enumerate}


\subsubsection{Examples}

\begin{example}% d3reg8a
   {\tt d3reg8a}\newline 
{\rm  
Let $\Gamma_2$ have vertices

\[
V_2=\{0,1,\dots,7\}
\]
and edges

\[
E_2 = \{(0, 1), (0, 2), (0, 7), (1, 2), (1, 3), (2, 3), (3, 4),
\]
\[
(4, 5), (4, 6), (5, 6), (5, 7), (6, 7)
\}.
\]
In this case, there are $13$ non-constant harmonic quotient
morphisms\footnote{These morphisms are distinct, unless indicated otherwise.}
$\pi: \Gamma_2\to \Gamma_2/G$, for some subgroup
$G$ of $Aut(\Gamma_2)$. Of these, $2$ are isomorphic to $K_4$,
$2$ are isomorphic to the diamond graph,
$3$ are\footnote{All these are the same
morphism, even though the group $G$ is different, so there is
really only $1$ distinct harmonic quotient morphism to $P_3$.} isomorphic
to $P_3$, and $6$ are\footnote{All these are also the same
morphism, even though the group $G$ is different, so there is
really only $1$ distinct harmonic quotient morphism to $P_2$.}
isomorphic to $P_2$.

When $G=\langle (0,4)(1,6)(2,5)(3,7)\rangle$, then
$\Gamma_1=\Gamma_2/G \cong K_4$ is isomorphic to the complete graph
on $4$ letters. Note that both $\Gamma_2$ and $\Gamma_1$ are
cubic graphs.
The harmonic quotient map $\phi:\Gamma_2\to \Gamma_1$ is tabulated below
and  depicted in Figure \ref{fig:harmonic-quotient_d3reg8a-K4_30213201}.

\begin{figure}[H]
\begin{minipage}{\textwidth}
\begin{center}
%\vspace{8.0 cm}
\includegraphics[height=5cm,width=5cm]{graphics/harmonic-quotient_d3reg8a-K4_30213201.png}
\end{center}
\begin{center}
%\vspace{8.0 cm}
\includegraphics[height=4cm,width=4cm]{graphics/K4-0123b.png}
\end{center}
\end{minipage}
\caption{The cubic graph $\Gamma_2$, a degree two cover of $\Gamma_1=K_4$.}
\label{fig:harmonic-quotient_d3reg8a-K4_30213201}
\end{figure}

\begin{center}
\begin{tabular}{c|c|c|c}  
  $x$ & $\phi(x)$ & $m_\phi(x)$ & $\nu_\phi(x)$ \\ \hline
  0   &    3   &        1         &  0   \\ 
  1   &    0   &        1         &  0   \\ 
  2   &    2    &       1         &  0   \\ 
  3   &    1   &        1         &  0   \\ 
  4   &    3   &        1         &  0   \\ 
  5   &    2   &        1         &  0    \\ 
  6   &    0   &        1         &  0   \\ 
  7   &    1   &        1         &  0   \\ 
\end{tabular}
\end{center}
  
\[
  {\rm genus}(\Gamma_2) =5,\ \ \ \ \
  {\rm genus}(\Gamma_1) =3,
\]

\[
deg(\phi) =2.
\]
Note that the vertex matrix map for the map $\phi:\Gamma_2\to \Gamma_1\cong K_4$ is given by

\[
\Phi =
\left(\begin{array}{rrrr}
0 & 0 & 0 & 1 \\
1 & 0 & 0 & 0 \\
0 & 0 & 1 & 0 \\
0 & 1 & 0 & 0 \\
0 & 0 & 0 & 1 \\
0 & 0 & 1 & 0 \\
1 & 0 & 0 & 0 \\
0 & 1 & 0 & 0
\end{array}\right).
\]
This matrix is invariant under the quotient group
$G = \langle (0,4)(1,6)(2,5)(3,7)\rangle$, regarded as
a row permutation of the matrix.

When $G=\langle (0,7)(1,5)(2,6)(3,4) \rangle$, then
$\Gamma_1=\Gamma_2/G$ is isomorphic to the diamond graph.
The harmonic quotient morphism is depicted in
Figure \ref{fig:harmonic-quotient_d3reg8a-diamond_02311230}
and tabulated below.

\begin{center}
\begin{tabular}{c|c|c|c}  
  $x$ & $\phi(x)$ & $m_\phi(x)$ & $\nu_\phi(x)$ \\ \hline
  0   &    0   &        1         &  1   \\ 
  1   &    2   &        1         &  0   \\ 
  2   &    3    &       1         &  0   \\ 
  3   &    1   &        1         &  1   \\ 
  4   &    1   &        1         &  1   \\ 
  5   &    2   &        1         &  0    \\ 
  6   &    3   &        1         &  0   \\ 
  7   &    0   &        1         &  1   \\ 
\end{tabular}
\end{center}

\[
  {\rm genus}(\Gamma_2) =5,\ \ \ \ \
  {\rm genus}(\Gamma_1) =2,
\]

\[
deg(\phi) =2.
\]

\begin{figure}[H]
\begin{minipage}{\textwidth}
\begin{center}
%\vspace{8.0 cm}
\includegraphics[height=5cm,width=5cm]{graphics/harmonic-quotient_d3reg8a-diamond_02311230.png}
\end{center}
\begin{center}
%\vspace{8.0 cm}
\includegraphics[height=5cm,width=3cm]{graphics/diamond_0123c.png}
\end{center}
\end{minipage}
\caption{The cubic graph  $\Gamma_2$ ({\tt d3reg8a}), a degree two quasicover of the diamond graph.}
\label{fig:harmonic-quotient_d3reg8a-diamond_02311230}
\end{figure}


Note all three of these graphs (namely, $\Gamma_2$, $K_4$, and the diamond graph)
are matchstick graphs.
}
\end{example}



\begin{example}% d3reg8b
    {\tt d3reg8b}\newline
{\rm  
Let $\Gamma_2$ have vertices

\[
V_2=\{0,1,\dots,7\}
\]
and edges

\[
E_2 = \{(0, 1), (0, 4), (0, 7), (1, 2), (1, 7), (2, 3), (2, 6),
\]
\[
(3, 4), (3, 5), (4, 5), (5, 6), (6, 7)\}.
\]
In this case, there is only $1$ harmonic quotient morphism
$\pi: \Gamma_2\to \Gamma_2/G$, for some subgroup
$G$ of $Aut(\Gamma_2)$.
When $G=\langle (0,4)(1,5)(2,6)(3,7)\rangle$, then
$\Gamma_1=\Gamma_2/G$ is isomorphic to the diamond graph.
The harmonic quotient morphism (a quasicover) is depicted in
Figure \ref{fig:harmonic-quotient_d3reg8b-diamond_21302130}
and tabulated below.

\begin{center}
\begin{tabular}{c|c|c|c}  
  $x$ & $\phi(x)$ & $m_\phi(x)$ & $\nu_\phi(x)$ \\ \hline
  0   &    2   &        1         &  0   \\ 
  1   &    1   &        1         &  1   \\ 
  2   &    3    &       1         &  0   \\ 
  3   &    0   &        1         &  1   \\ 
  4   &    2   &        1         &  0   \\ 
  5   &    1   &        1         &  1    \\ 
  6   &    3   &        1         &  0   \\ 
  7   &    0   &        1         &  1   \\ 
\end{tabular}
\end{center}

\[
  {\rm genus}(\Gamma_2) =5,\ \ \ \ \
  {\rm genus}(\Gamma_1) =2,
\]

\[
deg(\phi) =2.
\]

\begin{figure}[H]
\begin{minipage}{\textwidth}
\begin{center}
%\vspace{8.0 cm}
\includegraphics[height=5cm,width=5cm]{graphics/harmonic-quotient_d3reg8b-diamond_21302130.png}
\end{center}
\begin{center}
%\vspace{8.0 cm}
\includegraphics[height=5cm,width=4cm]{graphics/diamond-0123b.png}
\end{center}
\end{minipage}
\caption{The cubic graph  $\Gamma_2$ ({\tt d3reg8b}), a degree two quasicover of the diamond graph.}
\label{fig:harmonic-quotient_d3reg8b-diamond_21302130}
\end{figure}

}
\end{example}


\begin{example}% d3reg8c
    {\tt d3reg8c}\newline
{\rm  
Let $\Gamma_2$ have vertices

\[
V_2=\{0,1,\dots,7\}
\]
and edges

\[
E_2 = \{
(0, 1), (0, 2), (0, 7), (1, 2), (1, 5), (2, 3), (3, 4),
\]
\[
(3, 6), (4, 5), (4, 7), (5, 6), (6, 7)
\}.
\]
In this case, there are only $2$ distinct
subgroups $G$ of $Aut(\Gamma_2)$ with
harmonic quotient morphisms $\pi: \Gamma_2\to \Gamma_2/G$.
However, in both cases, $\Gamma_1=\Gamma_2/G$
is the same graph (isomorphic to the star graph $St_4$ on $4$ vertices)
and the morphisms are also the same.
\index{star graph}

When $G=\langle (0,2,1)(3,5,7) \rangle$, 
the harmonic quotient morphism is depicted in
Figure \ref{fig:harmonic-quotient_d3reg8c-star4_33302010}
and tabulated below.

\begin{center}
\begin{tabular}{c|c|c|c}  
  $x$ & $\phi(x)$ & $m_\phi(x)$ & $\nu_\phi(x)$ \\ \hline
  0   &    3   &        1         &  2   \\ 
  1   &    3   &        1         &  2   \\ 
  2   &    3    &       1         &  2   \\ 
  3   &    0   &        1         &  0   \\ 
  4   &    2   &        3         &  0   \\ 
  5   &    0   &        1         &  0    \\ 
  6   &    1   &        3         &  0   \\ 
  7   &    0   &        1         &  0   \\ 
\end{tabular}
\end{center}

\[
  {\rm genus}(\Gamma_2) =5,\ \ \ \ \
  {\rm genus}(\Gamma_1) =0,
\]

\[
deg(\phi) =3.
\]

\begin{figure}[H]
\begin{minipage}{\textwidth}
\begin{center}
%\vspace{8.0 cm}
\includegraphics[height=5cm,width=5cm]{graphics/harmonic-quotient_d3reg8c-star4_33302010.png}
\end{center}
\begin{center}
%\vspace{8.0 cm}
\includegraphics[height=5cm,width=4cm]{graphics/star4-0123.png}
\end{center}
\end{minipage}
\caption{The cubic graph  $\Gamma_2$ ({\tt d3reg8c}), a degree three harmonic quotient to the star graph, $St_4$.}
\label{fig:harmonic-quotient_d3reg8c-star4_33302010}
\end{figure}

When $G=\langle (1,2)(3,5), (0,2,1)(3,5,7) \rangle$
(a subgroup of order $6$),
the harmonic quotient morphism is the same as in
Figure \ref{fig:harmonic-quotient_d3reg8c-star4_33302010},
tabulated above.

}
\end{example}



\begin{example}% d3reg8d
    {\tt d3reg8d}\newline
{\rm  
Let $\Gamma_2$ have vertices

\[
V_2=\{0,1,\dots,7\}
\]
and edges

\[
E_2 = \{(0, 1), (0, 5), (0, 7), (1, 2), (1, 4), (2, 3), (2, 7),
\]
\[
(3, 4), (3, 6), (4, 5), (5, 6), (6, 7)
\}.
\]
This graph has $55$ non-constant harmonic quotient
morphisms\footnote{These morphisms are distinct, unless indicated otherwise.}
$\pi: \Gamma_2\to \Gamma_2/G$, for some subgroup
$G$ of $Aut(\Gamma_2)$. Of these, $3$ of these quotient graphs
are isomorphic to $C_4$, $6$ are isomorphic to the diamond graph,
$1$ is isomorphic to $K_4$,
$6$ are isomorphic to $P_3$, and the
remaining\footnote{Of these $39$, many are the same
morphism, even though the group $G$ is different, so there are
really only $17$ distinct harmonic quotient morphisms to $P_2$.}
$39$ are $P_2$. 

When $G=\langle (0,3)(1,6)(2,5)(4,7) \rangle$, 
the harmonic quotient morphism is depicted in
Figure \ref{fig:harmonic-quotient_d3reg8d-K4_30132102}
and tabulated below.

\begin{center}
\begin{tabular}{c|c|c|c}  
  $x$ & $\phi(x)$ & $m_\phi(x)$ & $\nu_\phi(x)$ \\ \hline
  0   &    3   &        1         &  0   \\ 
  1   &    0   &        1         &  0   \\ 
  2   &    1    &       1         &  0   \\ 
  3   &    3   &        1         &  0   \\ 
  4   &    2   &        1         &  0   \\ 
  5   &    1   &        1         &  0    \\ 
  6   &    0   &        1         &  0   \\ 
  7   &    2   &        1         &  0   \\ 
\end{tabular}
\end{center}

\[
  {\rm genus}(\Gamma_2) =5,\ \ \ \ \
  {\rm genus}(\Gamma_1) =3,
\]

\[
deg(\phi) =2.
\]

\begin{figure}[H]
\begin{minipage}{\textwidth}
\begin{center}
%\vspace{8.0 cm}
\includegraphics[height=5cm,width=5cm]{graphics/harmonic-quotient_d3reg8d-K4_30132102.png}
\end{center}
\begin{center}
%\vspace{8.0 cm}
\includegraphics[height=5cm,width=4cm]{graphics/K4-0123.png}
\end{center}
\end{minipage}
\caption{The $3$-dimensional cube graph  $\Gamma_2$ ({\tt d3reg8d}), with a degree two harmonic quotient morphism (a covering) to the complete graph, $K_4$.}
\label{fig:harmonic-quotient_d3reg8d-K4_30132102}
\end{figure}

When $G=\langle (0,1)(2,7)(3,6)(4,5) \rangle$,
the harmonic quotient morphism (a quasicover) is depicted in
Figure \ref{fig:harmonic-quotient_d3reg8d-C4_11320023} and
tabulated above.

\begin{center}
\begin{tabular}{c|c|c|c}  
  $x$ & $\phi(x)$ & $m_\phi(x)$ & $\nu_\phi(x)$ \\ \hline
  0   &    1   &        1         &  1   \\ 
  1   &    1   &        1         &  1   \\ 
  2   &    3   &        1         &  1   \\ 
  3   &    2   &        1         &  1   \\ 
  4   &    0   &        1         &  1   \\ 
  5   &    0   &        1         &  1    \\ 
  6   &    2   &        1         &  1   \\ 
  7   &    3   &        1         &  1   \\ 
\end{tabular}
\end{center}

\[
  {\rm genus}(\Gamma_2) =5,\ \ \ \ \
  {\rm genus}(\Gamma_1) =1,
\]

\[
deg(\phi) =2.
\]


\begin{figure}[H]
\begin{minipage}{\textwidth}
\begin{center}
%\vspace{8.0 cm}
\includegraphics[height=5cm,width=5cm]{graphics/harmonic-quotient_d3reg8d-C4_11320023.png}
\end{center}
\begin{center}
%\vspace{8.0 cm}
\includegraphics[height=5cm,width=4cm]{graphics/cycle4-0123b.png}
\end{center}
\end{minipage}
\caption{The Wagner graph  $\Gamma_2$ ({\tt d3reg8d}), with a degree two harmonic quotient morphism (a quasicovering) to the cycle graph, $C_4$.}
\label{fig:harmonic-quotient_d3reg8d-C4_11320023}
\end{figure}

}
\end{example}



\begin{example}% d3reg8e
    {\tt d3reg8e}\newline
{\rm  
Let $\Gamma_2$ have vertices

\[
V_2=\{0,1,\dots,7\}
\]
and edges

\[
E_2 = \{(0, 1), (0, 4), (0, 7), (1, 2), (1, 6), (2, 3),
\]
\[
(2, 5), (3, 4), (3, 7), (4, 5), (5, 6), (6, 7)
\}.
\]
This Wagner graph has $7$ non-constant harmonic quotient
morphisms\footnote{These morphisms are distinct, unless indicated otherwise.}
$\pi: \Gamma_2\to \Gamma_2/G$, for some subgroup
$G$ of $Aut(\Gamma_2)$. Of these quotient graphs,
$1$ is isomorphic to $C_4$, $2$ are isomorphic to $P_3$, and
$4$ are\footnote{Of these $4$, two are the same
morphism, even though the group $G$ is different, so there are
really only $3$ distinct harmonic quotient morphisms to $P_2$.}
isomorphic to $P_2$.

When $G=\langle (0,7)(1,6)(2,5)(3,4) \rangle$,
the harmonic quotient morphism (a quasicover) is depicted in
Figure \ref{fig:harmonic-quotient_d3reg8e-C4_03211230} and
tabulated above.

\begin{center}
\begin{tabular}{c|c|c|c}  
  $x$ & $\phi(x)$ & $m_\phi(x)$ & $\nu_\phi(x)$ \\ \hline
  0   &    1   &        1         &  1   \\ 
  1   &    1   &        1         &  1   \\ 
  2   &    3   &        1         &  1   \\ 
  3   &    2   &        1         &  1   \\ 
  4   &    0   &        1         &  1   \\ 
  5   &    0   &        1         &  1    \\ 
  6   &    2   &        1         &  1   \\ 
  7   &    3   &        1         &  1   \\ 
\end{tabular}
\end{center}

\[
  {\rm genus}(\Gamma_2) =5,\ \ \ \ \
  {\rm genus}(\Gamma_1) =1,
\]

\[
deg(\phi) =2.
\]


\begin{figure}[H]
\begin{minipage}{\textwidth}
\begin{center}
%\vspace{8.0 cm}
\includegraphics[height=5cm,width=5cm]{graphics/harmonic-quotient_d3reg8e-C4_03211230.png}
\end{center}
\begin{center}
%\vspace{8.0 cm}
\includegraphics[height=5cm,width=4cm]{graphics/Cycle4-0123.png}
\end{center}
\end{minipage}
\caption{The cubic graph  $\Gamma_2$ ({\tt d3reg8e}), with a degree two harmonic quotient morphism (a quasicovering) to the cycle graph, $C_4$.}
\label{fig:harmonic-quotient_d3reg8e-C4_03211230}
\end{figure}

}
\end{example}






\subsection{$10$ vertices}
%%% file-of-3regular10-graphs.txt
\label{subsec:d3reg10}

There are 19 $3$-regular graphs with $10$ vertices.
Three of these cubic graphs are vertex transitive and one of these
three (namely, the Petersen graph) is also edge transitive.
All but two of these graphs are Hamiltonian (circumference $|V|=10$).
One of these two has a bridge (the first one listed below)
and circumference $5$.
The other one of these two non-Hamiltonian graphs,
the Petersen graph, has circumference $9$.

{\small{
\begin{center}
\begin{tabular}{c|c|c|c|c}
  graph label & LCF & $|Aut(\Gamma)|$ & genus & $|K(\Gamma)|$ \\ \hline
  d3reg10a & -- & 32 & 6 & $576=2^63^2$ \\
  d3reg10b & (10, [4,2,3,-2,-4,-3,2,2,-2,-2], 1) & 4 & 6 & $1160=2^35\cdot 29$ \\
  d3reg10c & (10, [2,-3,-2,2,2,-2,-2,2,3,-2], 1) & 8 & 6 & $1080=2^33^35$ \\
  d3reg10d & (10, [-2,-2,3,3,3,-3,-3,-3,2,2], 1) & 16 & 6 & $1224=2^33^217$ \\ 
  d3reg10e & (10, [2,2,-2,-2,5], 2) & 16 & 6 & $960=2^63\cdot 5$ \\
  d3reg10f & (10, [2,3,-2,5,-3], 2) & 2 & 6 & $1488=2^43\cdot 31$ \\  
  d3reg10g & (10, [2,-4,-2,5,2,4,-2,4,5,-4], 1) & 12 & 6 & $1600=2^65^2$ \\  
  d3reg10h & (10, [5,3,5,-4,-3,5,2,5,-2,4], 1) & 2 & 6 & $1717=17\cdot 101$ \\  
  d3reg10i & (10, [-4,3,3,5,-3,-3,4,2,5,-2], 1) & 4 & 6 & $1599=3\cdot 13\cdot 41$ \\  
  d3reg10j & (10, [3,-3,5,-3,2,4,-2,5,3,-4], 1) & 6 & 6 & $1682 = 2\cdot 29^2$ \\  
  d3reg10k & (10, [-4, 4, 2, 5, -2], 2) & 4 & 6 & $1463=7\cdot 11\cdot 19$ \\  
  d3reg10l & (10, [5,-2,2,4,-2,5,2,-4,-2,2], 1) & 6 & 6 & $1445= 5\cdot 17^2$ \\  
  d3reg10m & (10, [2,5,-2,5,5], 2) & 8 & 6 & $1573=11^213$ \\
  d3reg10n & (10, [5,-3,-3,3,3], 2) & 48 & 6 & $1728 = 2^6 3^3$ \\    
  d3reg10o & (10, [5,-4,4,-4,4], 2) & 8  & 6 & $1920=2^7 3\cdot 5$ \\    
  d3reg10p & (10, [5,-4,4,5,5], 2) & 4  & 6 & $1881=3^2 11\cdot 19$ \\    
  d3reg10q & (10, [5,5,-3,5,3], 2)  & 20  & 6 & $1815=3\cdot 5\cdot 11^2$ \\    
  d3reg10r & (10,[-4, 4, -3, 5, 3],2) & 20  & 6 & $1805=5\cdot 19^2$ \\    
  d3reg10s & --  & 120  & 6 & $2000=2^45^3$ \\    
\end{tabular}
\end{center}  
}}

\begin{enumerate}
\item
  {\tt d3reg10a}\newline
The one with vertices and edges:

\[
V=\{ 0, 1, 2, 3, 4, 5, 6, 7, 8, 9\},
\]
\[
E = \{ (0, 1), (0, 6), (0, 9), (1, 2), (1, 5), (2, 3), (2, 4), (3, 4),
\]
\[
(3, 5), (4, 5), (6, 7), (6, 8), (7, 8), (7, 9), (8, 9)\}.
\]
It has %automorphism group of order $32$,
characteristic polynomial  
\[
x^{10} - 15x^8 - 8x^7 + 63x^6 + 64x^5 - 37x^4 - 56x^3 - 12x^2,
\]
Diameter $5$,
girth $3$,
circumference $5$.
This is not Hamiltonian (indeed, it has a bridge).
%It has critical group of cardinality $576=2^63^2$.

It has $8$ non-trivial harmonic morphisms to $P_3$, but
no non-trivial harmonic morphisms to $P_4$, $St_4$, $C_3$, $C_4$, or to $K_4$.


\item
  {\tt d3reg10b}\newline
The one with characteristic polynomial  
\[
x^{10} - 15x^8 - 6x^7 + 69x^6 + 48x^5 - 96x^4 - 76x^3 + 30x^2 + 26x + 3
\]
Diameter $4$,
girth $3$,
circumference $10$.
This is a Hamiltonian graph.
%It has critical group of cardinality $1160=2^35\cdot 29$.

This particular graph is very unusual in that it lacks non-trivial
harmonic morphisms to all ``very small'' graphs.
It has no non-trivial harmonic morphisms to the diamond graph,
$St_4$, $P_3$, (hence none to $P_4$, $P_5$, \dots),
%by Lemma \ref{lemma:families}),
$C_3$, $C_4$, or to $K_4$.

SageMath command:

\verb@graphs.LCFGraph(10, [4, 2, 3, -2, -4, -3, 2, 2, -2, -2], 1)@


\item
  {\tt d3reg10c}\newline
The one with %automorphism group of order $8$ and having
characteristic polynomial  
\[
x^{10} - 15x^8 - 8x^7 + 71x^6 + 68x^5 - 93x^4 - 132x^3 - 36x^2 
\]
Diameter $3$,
girth $3$,
circumference $10$.
This is a Hamiltonian graph.% (with circumference $10$).
%It has critical group of cardinality $1080=2^33^35$.

It has $8$ non-trivial harmonic morphisms to $P_3$,
$6$ non-trivial harmonic morphisms to $C_3$,
but no non-trivial harmonic morphisms to the diamond graph,
$St_4$, $P_4$, $C_4$, or to $K_4$.

SageMath command:

\verb@graphs.LCFGraph(10, [2, -3, -2, 2, 2, -2, -2, 2, 3, -2], 1)@


\item
  {\tt d3reg10d}\newline
The one with %automorphism group of order $16$ and having
characteristic polynomial  
\[
x^{10} - 15x^8 - 4x^7 + 63x^6 + 63x^5 - 61x^4 - 56x^3 - 12x^2 
\]
Diameter $3$,
girth $3$,
circumference $10$.
This is a Hamiltonian graph.
%It has critical group of cardinality $1224=2^33^217$.

It has $12$ non-trivial harmonic morphisms to $St_4$,
$12$ non-trivial harmonic morphisms to $P_3$,
$12$ non-trivial harmonic morphisms to $C_3$,
but no non-trivial harmonic morphisms to the diamond graph,
$P_4$, $C_4$, or to $K_4$.

SageMath command:

\verb@graphs.LCFGraph(10, [-2, -2, 3, 3, 3, -3, -3, -3, 2, 2], 1)@


\item
  {\tt d3reg10e}\newline
The one with %automorphism group of order $16$ and having
characteristic polynomial  

\[
x^{10} - 15x^8 - 8x^7 + 71x^6 + 64x^5 - 101x^4 - 104x^3 + 44x^2 + 48x 
\]
Diameter $4$,
girth $3$,
circumference $10$.
This is a Hamiltonian graph.% (with circumference $10$).
%It has critical group of cardinality $960=2^63\cdot 5$.

It has $4$ non-trivial harmonic morphisms to $P_3$,
but no non-trivial harmonic morphisms to $St_4$, $P_4$, $C_3$, $C_4$,
or to $K_4$.

SageMath command:

\verb@graphs.LCFGraph(10,[2,2,-2,-2,5],2)@


\item
  {\tt d3reg10f}\newline
The one with %automorphism group of order $2$ and having
characteristic polynomial  

\[
x^{10} - 15x^8 - 4x^7 + 71x^6 + 28x^5 - 121x^4 - 48x^3 + 64x^2 + 24x 
\]
Diameter $3$,
girth $3$,
circumference $10$.
This is a Hamiltonian graph.
%It has critical group of cardinality $1488=2^43\cdot 31$.

It has $2$ non-trivial harmonic morphisms to $P_3$,
but no non-trivial harmonic morphisms to $St_4$, $P_4$, $C_3$, $C_4$, or to
$K_4$.

SageMath command:

\verb@graphs.LCFGraph(10,[2,3,-2,5,-3],2)@


\item
  {\tt d3reg10g}\newline
The one with %automorphism group of order $12$ and having
characteristic polynomial  

\[
x^{10} - 15x^8 - 4x^7 + 75x^6 + 24x^5 - 157x^4 - 36x^3 + 144x^2 + 16x - 48
\]
Diameter $3$,
girth $3$,
circumference $10$.
This is a Hamiltonian graph.
%It has critical group of cardinality $1600=2^65^2$.

It has $6$ non-trivial harmonic morphisms to $St_4$,
$6$ non-trivial harmonic morphisms to $P_3$, but
no non-trivial harmonic morphisms to $P_4$, $C_3$, $C_4$, or to $K_4$.

SageMath command:

\verb@graphs.LCFGraph(10, [2, -4, -2, 5, 2, 4, -2, 4, 5, -4], 1)@


\item
  {\tt d3reg10h}\newline
The one with %automorphism group of order $2$ and having
characteristic polynomial  

\[
x^{10} - 15x^8 - 2x^7 + 71x^6 + 8x^5 - 132x^4 - 2x^3 + 91x^2 - 8x - 12
\]
Diameter $3$,
girth $3$,
circumference $10$.
This is a Hamiltonian graph.
%This has critical group of order $1717=17\cdot 101$.

It has $2$ non-trivial harmonic morphisms to $P_3$, but
no non-trivial harmonic morphisms to $St_4$, $P_4$,
$C_3$, $C_4$, or to $K_4$.

SageMath command:

\verb@graphs.LCFGraph(10,[5, 3, 5, -4, -3, 5, 2, 5, -2, 4], 1)@


\item
  {\tt d3reg10i}\newline
The one with %automorphism group of order $4$ and having
characteristic polynomial  

\[
x^{10} - 15x^8 - 2x^7 + 67x^6 + 12x^5 - 96x^4 - 22x^3 + 35x^2 + 12x 
\]
Diameter $3$,
girth $3$,
circumference $10$.
This is a Hamiltonian graph.% (with circumference $10$).
%This has critical group of order $1599=3\cdot 13\cdot 41$.

It has $6$ non-trivial harmonic morphisms to $St_4$,
$6$ non-trivial harmonic morphisms to $P_3$,
but
no non-trivial harmonic morphisms to $P_4$, $C_3$, $C_4$, or to $K_4$.

SageMath command:

\verb@graphs.LCFGraph(10,[-4, 3, 3, 5, -3, -3, 4, 2, 5, -2],1)@


\item
  {\tt d3reg10j}\newline
The one with %automorphism group of order $6$ and having
characteristic polynomial  

\[
x^{10} - 15x^8 - 2x^7 + 69x^6 + 12x^5 - 116x^4 - 24x^3 + 54x^2 + 26x + 3 
\]
Diameter $3$,
girth $3$,
circumference $10$.
%has girth $3$ and diameter $3$.
This is a Hamiltonian graph.% (with circumference $10$).
%This has critical group of order $1682 = 2\cdot 29^2$.

It has $6$ non-trivial harmonic morphisms to $P_3$, but
no non-trivial harmonic morphisms to
$St_4$, $P_4$, $C_3$, $C_4$, or to $K_4$.

SageMath command:

\verb@graphs.LCFGraph(10,[3, -3, 5, -3, 2, 4, -2, 5, 3, -4], 1)@


\item
  {\tt d3reg10k}\newline
The one with %automorphism group of order $4$ and having
characteristic polynomial  

\[
x^{10} - 15x^8 - 4x^7 + 69x^6 + 32x^5 - 105x^4 - 64x^3 + 23x^2 + 20x + 3 
\]
Diameter $3$,
girth $3$,
circumference $10$.
This is a Hamiltonian graph.% (with circumference $10$).
This has critical group of order $1463=7\cdot 11\cdot 19$.

It has $6$ non-trivial harmonic morphisms to $P_3$, but
no non-trivial harmonic morphisms to $St_4$,
$P_4$, $C_3$, $C_4$, or to $K_4$.

SageMath command:

\verb@graphs.LCFGraph(10,[-4, 4, 2, 5, -2], 2)@


\item
  {\tt d3reg10l}\newline
The one with %automorphism group of order $6$ and having
characteristic polynomial  

\[
x^{10} - 15x^8 - 6x^7 + 75x^6 + 48x^5 - 144x^4 - 114x^3 + 75x^2 + 68x + 12 
\]
Diameter $3$,
girth $3$,
circumference $10$.
This is a Hamiltonian graph.
%This has critical group of order $1445= 5\cdot 17^2$.

It has no non-trivial harmonic morphisms to $St_4$, $P_4$, $C_3$, $C_4$,
or to $K_4$. However, it has $6$ non-trivial harmonic morphisms to $P_3$.

SageMath command:

\verb@graphs.LCFGraph(10, [5, -2, 2, 4, -2, 5, 2, -4, -2, 2], 1)@


\item
  {\tt d3reg10m}\newline
The one with %automorphism group of order $8$ and having
characteristic polynomial  

\[
x^{10} - 15x^8 - 4x^7 + 73x^6 + 28x^5 - 141x^4 - 52x^3 + 99x^2 + 16x - 21 
\]
Diameter $3$,
girth $3$,
circumference $10$.
This is a Hamiltonian graph.% (with circumference $10$).
%This has critical group of order $1573=11^213$.

It has $2$ non-trivial harmonic morphisms to $P_3$, but
no non-trivial harmonic morphisms to $St_4$, $P_4$, $C_4$, or to $K_4$.

SageMath command:

\verb@graphs.LCFGraph(10,[2, 5, -2, 5, 5],2)@


\item
  {\tt d3reg10n}\newline
The one with %automorphism group of order $48$ and having
characteristic polynomial  

\[
x^{10} - 15x^8 + 63x^6 - 85x^4 + 36x^2  
\]
Diameter $3$,
girth $4$,
circumference $10$.
This is an integral\footnote{Its adjacency matrix has
  eigenvalues $3, 2, 1, 1, 0, 0, -1, -1, -2, -3$, while
  (since the graph is cubic) its Laplacian has eigenvalues
  $6, 5, 4, 4, 3, 3, 2, 2, 1, 0$.},
Hamiltonian graph.
%This has critical group of cardinality $1728 = 2^6 3^3$.

It has $18$ non-trivial harmonic morphisms to $St_4$,
$22$ non-trivial harmonic morphisms to $P_3$,
$8$ non-trivial harmonic morphisms to $P_4$,
but no non-trivial harmonic morphisms to $C_3$, $C_4$ or $K_4$.
By a computer check using SageMath , we know there
are no non-trivial harmonic morphisms to $P_5$.
%More generally, using Theorem \ref{thrm:diameter-bound}
Therefore, we know there
are no on-trivial harmonic morphisms to $P_k$, for $k\geq 5$.
\index{integral graph}

SageMath command:

\verb@graphs.LCFGraph(10, [5, -3, -3, 3, 3], 2)@


\item
  {\tt d3reg10o}\newline
The one with %automorphism group of order $8$ and having
characteristic polynomial  

\[
x^{10} - 15x^8 + 71x^6 - 15x^5 - 133x^4 + 64x^3 + 76x^2 - 48x  
\]
Diameter $3$,
girth $4$,
circumference $10$.
This is a Hamiltonian graph.
%The critical group has cardinality $1920=2^7 3\cdot 5$.

It has $6$ non-trivial harmonic morphisms to $St_4$,
$8$ non-trivial harmonic morphisms to $P_3$,
but no non-trivial harmonic morphisms to $P_4$, $C_3$, $C_4$, or to $K_4$.

SageMath command:

\verb@graphs.LCFGraph(10,[5, -4, 4, -4, 4],2)@


\item
  {\tt d3reg10p}\newline
The one with %automorphism group of order $4$ and having
characteristic polynomial  

\[
x^{10} - 15x^8 + 69x^6 -12x^5 - 117x^4 +36x^3 + 59x^2 -12x - 9  
\]
Diameter $3$,
girth $4$,
circumference $10$.
This is a Hamiltonian graph.% (with circumference $10$).
%This has critical group of order $1881=3^2 11\cdot 19$.

It has $2$ non-trivial harmonic morphisms to $P_3$,
but no non-trivial harmonic morphisms to $St_4$, $P_4$,
$C_3$, $C_4$, or to $K_4$.


SageMath command:

\verb@graphs.LCFGraph(10,[5, -4, 4, 5, 5],2)@


\item
  {\tt d3reg10q}\newline
The one with %automorphism group of order $20$ (acting transitively on the vertices of the graph) and having
characteristic polynomial  

\[
x^{10} - 15x^8 + 65x^6 - 105x^4 + 55x^2 - 9  
\]
Diameter $3$,
girth $4$,
circumference $10$.
This graph is vertex transitive, not edge transitive, but it is
Hamiltonian.% (with circumference $10$).
This has critical group of order $1815=3\cdot 5\cdot 11^2$.

It has $10$ non-trivial harmonic morphisms to $P_3$,
but no non-trivial harmonic morphisms to $St_4$, $P_4$, $C_3$,
$C_4$, or to $K_4$.

SageMath command:

\verb@graphs.LCFGraph(10,[5, 5, -3, 5, 3],2)@


\item
  {\tt d3reg10r}\newline
The one with %automorphism group of order $20$ (acting transitively on the vertices of the graph) and having
characteristic polynomial  

\[
x^{10} - 15x^8 + 65x^6 - 4x^5 - 85x^4 - 20x^3 + 35x^2 + 20x + 3
\]
Diameter $3$,
girth $4$,
circumference $10$.
This graph is vertex transitive, not edge transitive,
but it is Hamiltonian.% (with circumference $10$).
%This has critical group of order $1805=5\cdot 19^2$.

It has $10$ non-trivial harmonic morphisms to $P_3$,
but no non-trivial harmonic morphisms to $St_4$, $P_4$, $C_3$, $C_4$,
or to $K_4$.


SageMath command:

\verb@graphs.LCFGraph(10,[-4, 4, -3, 5, 3],2)@


\item
  {\tt d3reg10s}
  The Petersen graph with vertices and edges:

\[
V=\{ 0,1,2,3,4,5,6,7,8,9\},
\]
\[
E = \{ (0, 1), (0, 4), (0, 5), (1, 2), (1, 6), (2, 3), (2, 7), (3, 4), (3, 8),
\]
\[
(4, 9), (5, 7), (5, 8), (6, 8), (6, 9), (7, 9) \}.
\]
The automorphism group acts transitively
on the vertices {\it and} the edges of the graph. It has
characteristic polynomial  
\index{Petersen graph}

\[
x^{10} - 15x^8 + 75x^6 - 24x^5 - 165x^4 + 120x^3 + 120x^2 - 160x +48,
\]
Diameter $2$,
girth $5$,
circumference $9$.
It is symmetric (arc transitive),
distance regular, integral, and is strongly regular with
parameters $(10,3,0,1)$.
This is {\it not} a Hamiltonian graph, but has circumference $9$.
%This has critical group of order $2000=2^45^3$.
\index{integral graph}

It has $30$ non-trivial harmonic morphisms to $St_4$ and
$30$ non-trivial harmonic morphisms to $P_3$.
However, it has
no non-trivial harmonic morphisms to $P_4$, $C_3$, $C_4$, or to $K_4$.

\end{enumerate}  



\subsubsection{Examples}
\label{sec:d3reg10-examples}

\begin{example}% d3reg10a
  {\tt d3reg10a}\newline
{\rm  
Let $\Gamma_2$ have vertices

\[
V_2=\{0,1,\dots,9\}
\]
and edges

\[
E_2 = \{
(0, 1), (0, 6), (0, 9), (1, 2), (1, 5), (2, 3), (2, 4), (3, 4),
\]
\[
(3, 5), (4, 5), (6, 7), (6, 8), (7, 8), (7, 9), (8, 9)
\}.
\]
This graph has $4$ distinct non-constant harmonic quotient
morphisms $\pi: \Gamma_2\to \Gamma_2/G$, for some subgroup
$G$ of $Aut(\Gamma_2)$. All quotients are isomorphic to the
graph\footnote{This graph is basically the wheel graph
with $4$ ``rim'' vertices, but remove one of the ``spoke'' edges.
}
$\Gamma_1$ have vertices

\[
V_1=\{0,1,\dots,4\}
\]
and edges

\[
E_1 = \{
(0, 1), (0, 3), (0, 4), (1, 2), (1, 4), (2, 3), (3, 4)
\}.
\]
One of these quotient morphisms is tabulated below.
\begin{center}
\begin{tabular}{c|c|c|c}  
  $x$ & $\phi(x)$ & $m_\phi(x)$ & $\nu_\phi(x)$ \\ \hline
  0   &    2   &        1         &  1   \\ 
  1   &    2   &        1         &  1   \\ 
  2   &    3   &        1         &  0   \\ 
  3   &    0   &        1         &  0   \\ 
  4   &    4   &        1         &  0   \\ 
  5   &    1   &        1         &  0    \\ 
  6   &    1   &        1         &  0   \\ 
  7   &    4   &        1         &  0   \\
  8   &    0   &        1         &  0   \\
  9   &    3   &        1         &  0   \\ 
\end{tabular}
\end{center}

\[
  {\rm genus}(\Gamma_2) =6,\ \ \ \ \
  {\rm genus}(\Gamma_1) =3,
\]

\[
deg(\phi) =2.
\]

}
\end{example}


\begin{example}% d3reg10b
  {\tt d3reg10b}\newline
{\rm  
Let $\Gamma_2$ have vertices

\[
V_2=\{0,1,\dots,9\}
\]
and edges

\[
E_2 = \{
(0, 1), (0, 4), (0, 9), (1, 2), (1, 3), (2, 3), (2, 5), (3, 4),
\]
\[
(4, 5), (5, 6), (6, 7), (6, 8), (7, 8), (7, 9), (8, 9)
\}.
\]
This graph has no non-constant harmonic quotient
morphisms $\pi: \Gamma_2\to \Gamma_2/G$, for some subgroup
$G$ of $Aut(\Gamma_2)$.

}
\end{example}


\begin{example}% d3reg10c
  {\tt d3reg10c}\newline
{\rm  
Let $\Gamma_2$ have vertices

\[
V_2=\{0,1,\dots,9\}
\]
and edges

\[
E_2 = \{
(0, 1), (0, 2), (0, 9), (1, 2), (1, 8), (2, 3), (3, 4), (3, 5),
\]
\[
(4, 5), (4, 6), (5, 6), (6, 7), (7, 8), (7, 9), (8, 9)
\}.
\]
This graph has no non-constant harmonic quotient
morphisms $\pi: \Gamma_2\to \Gamma_2/G$, for some subgroup
$G$ of $Aut(\Gamma_2)$.

}
\end{example}

\begin{example}% d3reg10d
  {\tt d3reg10d}\newline
{\rm  
Let $\Gamma_2$ have vertices

\[
V_2=\{0,1,\dots,9\}
\]
and edges

\[
E_2 = \{
(0, 1), (0, 8), (0, 9), (1, 2), (1, 9), (2, 3), (2, 5), (3, 4),
\]
\[
(3, 6), (4, 5), (4, 7), (5, 6), (6, 7), (7, 8), (8, 9)
\}.
\]
This graph has no non-constant harmonic quotient
morphisms $\pi: \Gamma_2\to \Gamma_2/G$, for some subgroup
$G$ of $Aut(\Gamma_2)$.

}
\end{example}

\begin{example}% d3reg10e
  {\tt d3reg10e}\newline
{\rm  
Let $\Gamma_2$ have vertices

\[
V_2=\{0,1,\dots,9\}
\]
and edges

\[
E_2 = \{
(0, 1), (0, 2), (0, 9), (1, 2), (1, 3), (2, 3), (3, 4), (4, 5),
\]
\[
(4, 9), (5, 6), (5, 7), (6, 7), (6, 8), (7, 8), (8, 9)
\}.
\]
This graph has $2$ distinct non-constant harmonic quotient
morphisms $\pi: \Gamma_2\to \Gamma_2/G$, for some subgroup
$G$ of $Aut(\Gamma_2)$. All quotients are isomorphic to the
graph $\Gamma_1$ having vertices

\[
V_1=\{0,1,\dots,4\}
\]
and edges

\[
E_1 = \{
(0, 1), (0, 2), (0, 3), (1, 2), (1, 4), (2, 4), (3, 4)
\}.
\]
One of these quotient morphisms is tabulated below.
\begin{center}
\begin{tabular}{c|c|c|c}  
  $x$ & $\phi(x)$ & $m_\phi(x)$ & $\nu_\phi(x)$ \\ \hline
  0   &    4   &        1         &  0   \\ 
  1   &    1   &        1         &  0   \\ 
  2   &    2   &        1         &  0   \\ 
  3   &    0   &        1         &  0   \\ 
  4   &    3   &        1         &  0   \\ 
  5   &    4   &        1         &  0    \\ 
  6   &    2   &        1         &  0   \\ 
  7   &    1   &        1         &  0   \\
  8   &    0   &        1         &  0   \\
  9   &    3   &        1         &  1   \\ 
\end{tabular}
\end{center}

\[
  {\rm genus}(\Gamma_2) =6,\ \ \ \ \
  {\rm genus}(\Gamma_1) =3,
\]

\[
deg(\phi) =2.
\]
The uncolored graphs are depicted in Figure \ref{fig:harmonic-quotient_d3reg10e-Gamma1}.

\begin{figure}[H]
\begin{minipage}{\textwidth}
\begin{center}
%\vspace{8.0 cm}
\includegraphics[height=5cm,width=5cm]{graphics/graph_d3reg10e.png}
\end{center}
\begin{center}
%\vspace{8.0 cm}
\includegraphics[height=4cm,width=4cm]{graphics/C4-C3-C3.png}
\end{center}
\end{minipage}
\caption{The cubic graph  $\Gamma_2$ ({\tt d3reg10e}), with a degree two harmonic quotient morphism to and quasicover of, the graph, $\Gamma_1$.}
\label{fig:harmonic-quotient_d3reg10e-Gamma1}
\end{figure}

}
\end{example}


\begin{example}% d3reg10f
  {\tt d3reg10f}\newline
{\rm  
Let $\Gamma_2$ have vertices

\[
V_2=\{0,1,\dots,9\}
\]
and edges

\[
E_2 = \{
(0, 1), (0, 2), (0, 9), (1, 2), (1, 4), (2, 3), (3, 4), (3, 8), 
\]
\[
(4, 5), (5, 6), (5, 7), (6, 7), (6, 9), (7, 8), (8, 9)
\}.
\]
This graph has no non-constant harmonic quotient
morphisms $\pi: \Gamma_2\to \Gamma_2/G$, for some subgroup
$G$ of $Aut(\Gamma_2)$.

}
\end{example}

\begin{example}% d3reg10g
  {\tt d3reg10g}\newline
{\rm  
Let $\Gamma_2$ have vertices

\[
V_2=\{0,1,\dots,9\}
\]
and edges

\[
E_2 = \{
(0, 1), (0, 2), (0, 9), (1, 2), (1, 7), (2, 3), (3, 4), (3, 8),
\]
\[
(4, 5), (4, 6), (5, 6), (5, 9), (6, 7), (7, 8), (8, 9)
\}.
\]
This graph has only one\footnote{Technically, there are
two groups for which there is a harmonic quotient, but the
resulting quotient graphs and quotient morphisms are the same.}
non-constant harmonic quotient
morphisms $\pi: \Gamma_2\to \Gamma_2/G$, for some subgroup
$G$ of $Aut(\Gamma_2)$.

The quotient morphism associated to the subgroup
$G=\langle (0,1,2)(3,9,7)(4,5,6)\rangle$
is tabulated below and depicted in Figure \ref{fig:harmonic-quotient_d3reg10g-star4_3330111020}.

\begin{center}
\begin{tabular}{c|c|c|c}  
  $x$ & $\phi(x)$ & $m_\phi(x)$ & $\nu_\phi(x)$ \\ \hline
  0   &    3   &        1         &  2   \\ 
  1   &    3   &        1         &  2   \\ 
  2   &    3   &        1         &  2   \\ 
  3   &    0   &        1         &  0   \\ 
  4   &    1   &        1         &  2   \\ 
  5   &    1   &        1         &  2    \\ 
  6   &    1   &        1         &  2   \\ 
  7   &    0   &        1         &  0   \\
  8   &    2   &        3         &  0   \\
  9   &    0   &        1         &  0   \\ 
\end{tabular}
\end{center}

\[
  {\rm genus}(\Gamma_2) =6,\ \ \ \ \
  {\rm genus}(\Gamma_1) =0,
\]

\[
deg(\phi) =3.
\]
\begin{figure}[H]
\begin{minipage}{\textwidth}
\begin{center}
%\vspace{8.0 cm}
\includegraphics[height=5cm,width=5cm]{graphics/harmonic-quotient_d3reg10g-star4_3330111020.png}
\end{center}
\begin{center}
%\vspace{8.0 cm}
\includegraphics[height=5cm,width=4cm]{graphics/star4-0123.png}
\end{center}
\end{minipage}
\caption{The cubic graph  $\Gamma_2$ ({\tt d3reg10g}), a degree three harmonic quotient to the star graph, $St_4$.}
\label{fig:harmonic-quotient_d3reg10g-star4_3330111020}
\end{figure}

}
\end{example}


\begin{example}% d3reg10h
  {\tt d3reg10h}\newline
{\rm  
Let $\Gamma_2$ have vertices

\[
V_2=\{0,1,\dots,9\}
\]
and edges

\[
E_2 = \{
(0, 1), (0, 5), (0, 9), (1, 2), (1, 4), (2, 3), (2, 7), (3, 4), 
\]
\[
(3, 9), (4, 5), (5, 6), (6, 7), (6, 8), (7, 8), (8, 9)
\}.
\]
This graph has no non-constant harmonic quotient
morphisms $\pi: \Gamma_2\to \Gamma_2/G$, for some subgroup
$G$ of $Aut(\Gamma_2)$.

}
\end{example}


\begin{example}% d3reg10i
  {\tt d3reg10i}\newline
{\rm  
Let $\Gamma_2$ have vertices

\[
V_2=\{0,1,\dots,9\}
\]
and edges

\[
E_2 = \{
(0, 1), (0, 6), (0, 9), (1, 2), (1, 4), (2, 3), (2, 5), (3, 4),
\]
\[
(3, 8), (4, 5), (5, 6), (6, 7), (7, 8), (7, 9), (8, 9)
\}.
\]
This graph has no non-constant harmonic quotient
morphisms $\pi: \Gamma_2\to \Gamma_2/G$, for some subgroup
$G$ of $Aut(\Gamma_2)$.

}
\end{example}



\begin{example}% d3reg10j
  {\tt d3reg10j}\newline
{\rm  
Let $\Gamma_2$ have vertices

\[
V_2=\{0,1,\dots,9\}
\]
and edges

\[
E_2 = \{
(0, 1), (0, 3), (0, 9), (1, 2), (1, 8), (2, 3), (2, 7), (3, 4),
\]
\[
(4, 5), (4, 6), (5, 6), (5, 9), (6, 7), (7, 8), (8, 9)
\}.
\]
This graph has no non-constant harmonic quotient
morphisms $\pi: \Gamma_2\to \Gamma_2/G$, for some subgroup
$G$ of $Aut(\Gamma_2)$.

}
\end{example}


\begin{example}% d3reg10k
  {\tt d3reg10k}\newline
{\rm  
Let $\Gamma_2$ have vertices

\[
V_2=\{0,1,\dots,9\}
\]
and edges

\[
E_2 = \{
(0, 1), (0, 6), (0, 9), (1, 2), (1, 5), (2, 3), (2, 4), (3, 4),
\]
\[
(3, 8), (4, 5), (5, 6), (6, 7), (7, 8), (7, 9), (8, 9)
\}.
\]
This graph has only $1$ non-constant harmonic quotient
morphisms $\pi: \Gamma_2\to \Gamma_2/G$, for some subgroup
$G$ of $Aut(\Gamma_2)$. In this case,
$G = \langle (0,1)(2,9)(3,8)(4,7)(5,6) \rangle$ and
$\Gamma_1 =  \Gamma_2/G$ is the house graph, with vertices and edges:

\[
V_1 = \{0, \dots, 4\},
\]
\[
E_1 = \{(0, 3), (0, 4), (1, 2), (1, 4), (2, 3), (3, 4)\}.
\]
Note that vertices $3$ and $4$ are those with degree $3$.

The quotient morphism 
associated to the subgroup
$G=\langle (0,1,2)(3,9,7)(4,5,6)\rangle$
is tabulated below. As we can see, it's a quasicover
$\phi:\Gamma_2\to \Gamma_1$.

\begin{center}
\begin{tabular}{c|c|c|c}  
  $x$ & $\phi(x)$ & $m_\phi(x)$ & $\nu_\phi(x)$ \\ \hline
  0   &    2   &        1         &  1   \\ 
  1   &    2   &        1         &  1   \\ 
  2   &    3   &        1         &  0   \\ 
  3   &    0   &        1         &  1   \\ 
  4   &    4   &        1         &  0   \\ 
  5   &    1   &        1         &  1    \\ 
  6   &    1   &        1         &  1   \\ 
  7   &    4   &        1         &  0   \\
  8   &    0   &        1         &  1   \\
  9   &    3   &        1         &  0   \\ 
\end{tabular}
\end{center}

\[
  {\rm genus}(\Gamma_2) =6,\ \ \ \ \
  {\rm genus}(\Gamma_1) =2,
\]

\[
deg(\phi) =2.
\]
The uncolored graphs are depicted in Figure \ref{fig:harmonic-quotient_d3reg10k+house}.

\begin{figure}[H]
\begin{minipage}{\textwidth}
\begin{center}
%\vspace{8.0 cm}
\includegraphics[height=5cm,width=5cm]{graphics/graph_d3reg10k.png}
\end{center}
\begin{center}
%\vspace{8.0 cm}
\includegraphics[height=5cm,width=4cm]{graphics/house-graph.png}
\end{center}
\end{minipage}
\caption{The cubic graph  $\Gamma_2$ ({\tt d3reg10k}), with a degree two harmonic quotient morphism to and quasicover of, the house graph, $\Gamma_1$.}
\label{fig:harmonic-quotient_d3reg10k+house}
\end{figure}

}
\end{example}



\begin{example}% d3reg10l
  {\tt d3reg10l}\newline
{\rm  
Let $\Gamma_2$ have vertices

\[
V_2=\{0,1,\dots,9\}
\]
and edges

\[
E_2 = \{
(0, 1), (0, 5), (0, 9), (1, 2), (1, 9), (2, 3), (2, 4), (3, 4),
\]
\[
(3, 7), (4, 5), (5, 6), (6, 7), (6, 8), (7, 8), (8, 9)
\}.
\]
This graph has no non-constant harmonic quotient
morphisms $\pi: \Gamma_2\to \Gamma_2/G$, for some subgroup
$G$ of $Aut(\Gamma_2)$.

}
\end{example}



\begin{example}% d3reg10m
  {\tt d3reg10m}\newline
{\rm  
Let $\Gamma_2$ have vertices

\[
V_2=\{0,1,\dots,9\}
\]
and edges

\[
E_2 = \{
(0, 1), (0, 2), (0, 9), (1, 2), (1, 6), (2, 3), (3, 4), (3, 8), 
\]
\[
(4, 5), (4, 9), (5, 6), (5, 7), (6, 7), (7, 8), (8, 9)
\}.
\]
This graph has only $4$ non-constant harmonic quotient
morphisms $\pi: \Gamma_2\to \Gamma_2/G$, for some subgroup
$G$ of $Aut(\Gamma_2)$.
For $2$ subgroups, the quotient $ \Gamma_2/G$ is isomorphic to the house graph
and the quotient morphisms are distinct. For the remaining two subgroups,
the quotient $ \Gamma_2/G$ is
isomorphic\footnote{The vertices are $0,1,2$, but their ordering is different that usual:
  $0$ and $1$ are the vertices of degree $1$, while vertex $2$ has degree $2$.}
  to $P_3$ and both quotient morphisms are the same.

In case $G = \langle (0,5)(1,6)(2,7)(3,8)(4,9) \rangle$,
$\Gamma_1 =  \Gamma_2/G$ is the house graph, with vertices and edges:

\[
V_1 = \{0, \dots, 4\},
\]
\[
E_1 = \{(0, 1), (0, 2), (1, 4), (2, 3), (2, 4), (3, 4)\}.
\]
Note that vertices $2$ and $4$ are those with degree $3$.
The quotient morphism 
associated to this subgroup $G$ is tabulated below. As we can see, it's a quasicover
$\phi:\Gamma_2\to \Gamma_1$.

\begin{center}
\begin{tabular}{c|c|c|c}  
  $x$ & $\phi(x)$ & $m_\phi(x)$ & $\nu_\phi(x)$ \\ \hline
  0   &    4   &        1         &  0   \\ 
  1   &    3   &        1         &  1   \\ 
  2   &    2   &        1         &  0   \\ 
  3   &    0   &        1         &  1   \\ 
  4   &    1   &        1         &  1   \\ 
  5   &    4   &        1         &  0    \\ 
  6   &    3   &        1         &  1   \\ 
  7   &    2   &        1         &  0   \\
  8   &    0   &        1         &  1   \\
  9   &    1   &        1         &  1   \\ 
\end{tabular}
\end{center}

\[
  {\rm genus}(\Gamma_2) =6,\ \ \ \ \
  {\rm genus}(\Gamma_1) =2,
\]

\[
deg(\phi) =2.
\]

The quotient morphism $\phi:\Gamma_2\to \Gamma_1=P_3$ 
associated to the subgroup $G = \langle (0,2)(3,9)(4,8)(5,7), (0,5)(1,6)(2,7)(3,8)(4,9) \rangle$
is tabulated below. 
$\phi:\Gamma_2\to \Gamma_1$.

\begin{center}
\begin{tabular}{c|c|c|c}  
  $x$ & $\phi(x)$ & $m_\phi(x)$ & $\nu_\phi(x)$ \\ \hline
  0   &    2   &        1         &  1   \\ 
  1   &    0   &        2         &  1   \\ 
  2   &    2   &        1         &  1   \\ 
  3   &    1   &        1         &  2   \\ 
  4   &    1   &        1         &  2   \\ 
  5   &    2   &        1         &  1    \\ 
  6   &    0   &        2         &  1   \\ 
  7   &    2   &        1         &  1   \\
  8   &    1   &        1         &  2   \\
  9   &    1   &        1         &  2   \\ 
\end{tabular}
\end{center}

\[
  {\rm genus}(\Gamma_2) =6,\ \ \ \ \
  {\rm genus}(\Gamma_1) =0,
\]

\[
deg(\phi) =4.
\]

}
\end{example}




\begin{example}% d3reg10n
  {\tt d3reg10n}\newline
{\rm  
Let $\Gamma_2$ have vertices

\[
V_2=\{0,1,\dots,9\}
\]
and edges

\[
E_2 = \{
(0, 1), (0, 5), (0, 9), (1, 2), (1, 8), (2, 3), (2, 9), (3, 4),
\]
\[
(3, 6), (4, 5), (4, 7), (5, 6), (6, 7), (7, 8), (8, 9)
\}.
\]
This graph has $31$ non-constant harmonic quotient
morphisms $\pi: \Gamma_2\to \Gamma_2/G$, for some subgroup
$G$ of $Aut(\Gamma_2)$.

Two of these groups yield the same harmonic quotient morphism to
the ``$H$-tree'' - a genus $0$ graph $\Gamma_1 = (V_1,E_1)$
with vertices and edges,

\[
V_1=\{0,1,\dots, 5\},
\]
\[
E_1=\{
(0, 1), (0, 3), (0, 5), (1, 2), (1, 4)
\}.
\]
For example, the group $G=\langle (0,8,2)(3,5,7) \rangle$
yields this quotient graph $\Gamma_1=\Gamma_2/G$.

Eight of these groups yield the same harmonic quotient morphism to
$K_{2,3}$ - a genus $2$ graph $\Gamma_1 = (V_1,E_1)$
with vertices and edges,

\[
V_1=\{0,1,\dots, 4\},
\]
\[
E_1=\{
(0, 2), (0, 3), (0, 4), (1, 2), (1, 3), (1, 4)
\}.
\]
For example, the group $G=\langle (0,5)(1,4)(2,3)(6,9)(7,8) \rangle$
yields this quotient graph $\Gamma_1=\Gamma_2/G$.
One of these is tabulated below and the uncolored graphs is
depicted in Figure \ref{fig:harmonic-quotient_d3reg10n+K23}.
\begin{center}
\begin{tabular}{c|c|c|c}  
  $x$ & $\phi(x)$ & $m_\phi(x)$ & $\nu_\phi(x)$ \\ \hline
  0   &    3   &        1         &  1   \\ 
  1   &    1   &        1         &  0   \\ 
  2   &    2   &        1         &  1   \\ 
  3   &    2   &        1         &  1   \\ 
  4   &    1   &        1         &  0   \\ 
  5   &    3   &        1         &  1    \\ 
  6   &    0   &        1         &  0   \\ 
  7   &    4   &        1         &  1   \\
  8   &    4   &        1         &  1   \\
  9   &    0   &        1         &  0   \\ 
\end{tabular}
\end{center}

\[
  {\rm genus}(\Gamma_2) =6,\ \ \ \ \
  {\rm genus}(\Gamma_1) =2,
\]

\[
deg(\phi) =2.
\]

\begin{figure}[H]
\begin{minipage}{\textwidth}
\begin{center}
%\vspace{8.0 cm}
\includegraphics[height=5cm,width=5cm]{graphics/graph_d3reg10n.png}
\end{center}
\begin{center}
%\vspace{8.0 cm}
\includegraphics[height=5cm,width=4cm]{graphics/complete-bipartite-K-2-3.png}
\end{center}
\end{minipage}
\caption{The cubic graph  $\Gamma_2$ ({\tt d3reg10n}), with a degree two harmonic quotient morphism to and quasicover of, the complete bipartite graph, $\Gamma_1=K_{2,3}$.}
\label{fig:harmonic-quotient_d3reg10n+K23}
\end{figure}

Three of these groups yield the same harmonic quotient morphism to
$St_4$, the star graph on $4$ vertices.
\index{star graph}

The remaining harmonic quotient morphisms are either to $P_3$ or to
$P_2$.

}
\end{example}





\begin{example}% d3reg10o
  {\tt d3reg10o}\newline
{\rm  
Let $\Gamma_2$ have vertices

\[
V_2=\{0,1,\dots,9\}
\]
and edges

\[
E_2 = \{
(0, 1), (0, 5), (0, 9), (1, 2), (1, 7), (2, 3), (2, 6), (3, 4),
\]
\[
(3, 9), (4, 5), (4, 8), (5, 6), (6, 7), (7, 8), (8, 9)
\}.
\]
This graph has $3$ non-constant harmonic quotient
morphisms $\pi: \Gamma_2\to \Gamma_2/G$, for some subgroup
$G$ of $Aut(\Gamma_2)$.

The group $G = \langle (0,5)(1,4)(2,3)(6,9)(7,8) \rangle$ yields the
harmonic quotient morphism to the genus $2$ graph
$\Gamma_1 = (V_1,E_1) = K_{2,3}$
with vertices and edges,

\[
V_1=\{0,1,\dots, 4\},
\]
\[
E_1=\{
(0, 2), (0, 3), (0, 4), (1, 2), (1, 3), (1, 4)
\}.
\]
This is tabulated below (and agrees with a table in a previous example).

\begin{center}
\begin{tabular}{c|c|c|c}  
  $x$ & $\phi(x)$ & $m_\phi(x)$ & $\nu_\phi(x)$ \\ \hline
  0   &    3   &        1         &  1   \\ 
  1   &    1   &        1         &  0   \\ 
  2   &    2   &        1         &  1   \\ 
  3   &    2   &        1         &  1   \\ 
  4   &    1   &        1         &  0   \\ 
  5   &    3   &        1         &  1    \\ 
  6   &    0   &        1         &  0   \\ 
  7   &    4   &        1         &  1   \\
  8   &    4   &        1         &  1   \\
  9   &    0   &        1         &  0   \\ 
\end{tabular}
\end{center}

\[
  {\rm genus}(\Gamma_2) =6,\ \ \ \ \
  {\rm genus}(\Gamma_1) =2,
\]

\[
deg(\phi) =2.
\]

The group $G = \langle (0,5)(1,4)(2,8)(3,7)(6,9) \rangle$ yields the
harmonic quotient morphism to the genus $3$ graph
$\Gamma_1 = (V_1,E_1)$
with vertices and edges,

\[
V_1=\{0,1,\dots, 4\},
\]
\[
E_1=\{
(0, 1), (0, 2), (0, 3), (1, 3), (1, 4), (2, 3), (2, 4)
\}.
\]
This is tabulated below (and agrees with a table in a previous example).

\begin{center}
\begin{tabular}{c|c|c|c}  
  $x$ & $\phi(x)$ & $m_\phi(x)$ & $\nu_\phi(x)$ \\ \hline
  0   &    4   &        1         &  1   \\ 
  1   &    2   &        1         &  0   \\ 
  2   &    0   &        1         &  0   \\ 
  3   &    3   &        1         &  0   \\ 
  4   &    2   &        1         &  0   \\ 
  5   &    4   &        1         &  0    \\ 
  6   &    1   &        1         &  0   \\ 
  7   &    3   &        1         &  0   \\
  8   &    0   &        1         &  0   \\
  9   &    1   &        1         &  0   \\ 
\end{tabular}
\end{center}

\[
  {\rm genus}(\Gamma_2) =6,\ \ \ \ \
  {\rm genus}(\Gamma_1) =3,
\]

\[
deg(\phi) =2.
\]

Finally, the group
$G = \langle (1,9)(2,3)(4,6)(7,8), (0,5)(1,6)(4,9) \rangle$ yields the
harmonic quotient morphism to the genus $0$ graph
$\Gamma_1 = (V_1,E_1) = St_4$, with

\[
V_1=\{0, 1, 2, 3\},
\]
and $0$ is the degree $3$ (``center'') vertex.
This is tabulated below.

\begin{center}
\begin{tabular}{c|c|c|c}  
  $x$ & $\phi(x)$ & $m_\phi(x)$ & $\nu_\phi(x)$ \\ \hline
  0   &    3   &        2         &  1   \\ 
  1   &    0   &        1         &  0   \\ 
  2   &    1   &        2         &  1   \\ 
  3   &    1   &        2         &  1   \\ 
  4   &    0   &        1         &  0   \\ 
  5   &    3   &        2         &  1    \\ 
  6   &    0   &        1         &  0   \\ 
  7   &    2   &        2         &  1   \\
  8   &    2   &        2         &  1   \\
  9   &    0   &        1         &  0   \\ 
\end{tabular}
\end{center}

\[
  {\rm genus}(\Gamma_2) =6,\ \ \ \ \
  {\rm genus}(\Gamma_1) =0,
\]

\[
deg(\phi) =4.
\]


}
\end{example}



\begin{example}% d3reg10p
  {\tt d3reg10p}\newline
{\rm  
Let $\Gamma_2$ have vertices

\[
V_2=\{0,1,\dots,9\}
\]
and edges

\[
E_2 = \{
(0, 1), (0, 5), (0, 9), (1, 2), (1, 7), (2, 3), (2, 6), (3, 4), 
\]
\[
 (3, 8), (4, 5), (4, 9), (5, 6), (6, 7), (7, 8), (8, 9)
\}.
\]
This graph has only $1$ non-constant harmonic quotient
morphisms $\pi: \Gamma_2\to \Gamma_2/G$, for some subgroup
$G$ of $Aut(\Gamma_2)$.


The group $G = \langle (0,3)(1,2)(4,9)(5,8)(6,7) \rangle$ yields the
harmonic quotient morphism to the house graph
$\Gamma_1 = (V_1,E_1) = K_{2,3}$
with vertices and edges,

\[
V_1=\{0,1,\dots, 4\},
\]
\[
E_1=\{
(0, 1), (0, 3), (1, 2), (1, 4), (2, 4), (3, 4)
\}.
\]
This is tabulated below.

\begin{center}
\begin{tabular}{c|c|c|c}  
  $x$ & $\phi(x)$ & $m_\phi(x)$ & $\nu_\phi(x)$ \\ \hline
  0   &    1   &        1         &  0   \\ 
  1   &    0   &        1         &  1   \\ 
  2   &    0   &        1         &  1   \\ 
  3   &    1   &        1         &  0   \\ 
  4   &    2   &        1         &  1   \\ 
  5   &    4   &        1         &  0    \\ 
  6   &    3   &        1         &  1   \\ 
  7   &    3   &        1         &  1   \\
  8   &    4   &        1         &  0   \\
  9   &    2   &        1         &  1   \\ 
\end{tabular}
\end{center}

\[
  {\rm genus}(\Gamma_2) =6,\ \ \ \ \
  {\rm genus}(\Gamma_1) =2,
\]

\[
deg(\phi) =2.
\]

}
\end{example}





\begin{example}% d3reg10q
  {\tt d3reg10q}\newline
{\rm  
Let $\Gamma_2$ have vertices

\[
V_2=\{0,1,\dots,9\}
\]
and edges

\[
E_2 = \{
(0, 1), (0, 5), (0, 9), (1, 2), (1, 6), (2, 3), (2, 9), (3, 4), 
\]
\[
(3, 8), (4, 5), (4, 7), (5, 6), (6, 7), (7, 8), (8, 9)
\}.
\]
This graph has $13$ non-constant harmonic quotient
morphisms $\pi: \Gamma_2\to \Gamma_2/G$, for some subgroup
$G$ of $Aut(\Gamma_2)$.

Five of these quotient graphs are the house graph (each with a
distinct quotient morphism),
another five of these quotient graphs are the path graph $P_3$
(also with distinct quotient morphisms), 
however, the two groups that yield $P_2$ as a quotient have the
same quotient morphism.
Finally, the
 group $G = \langle (0,1)(2,9)(3,8)(4,7)(5,6) \rangle$ yields the
harmonic quotient morphism to the cycle graph
$\Gamma_1 = (V_1,E_1) = C_5$
with vertices and edges,

\[
V_1=\{0,1,\dots, 4\},
\]
\[
E_1=\{
(0, 3), (0, 4), (1, 2), (1, 4), (2, 3)
\}.
\]
This is tabulated below.

\begin{center}
\begin{tabular}{c|c|c|c}  
  $x$ & $\phi(x)$ & $m_\phi(x)$ & $\nu_\phi(x)$ \\ \hline
  0   &    2   &        1         &  1   \\ 
  1   &    2   &        1         &  1   \\ 
  2   &    3   &        1         &  1   \\ 
  3   &    0   &        1         &  1   \\ 
  4   &    4   &        1         &  1   \\ 
  5   &    1   &        1         &  1    \\ 
  6   &    1   &        1         &  1   \\ 
  7   &    4   &        1         &  1   \\
  8   &    0   &        1         &  1   \\
  9   &    3   &        1         &  1   \\ 
\end{tabular}
\end{center}

\[
  {\rm genus}(\Gamma_2) =6,\ \ \ \ \
  {\rm genus}(\Gamma_1) =1,
\]

\[
deg(\phi) =2.
\]


\begin{figure}[H]
\begin{minipage}{\textwidth}
\begin{center}
%\vspace{8.0 cm}
\includegraphics[height=5cm,width=5cm]{graphics/graph_d3reg10q.png}
\end{center}
\begin{center}
%\vspace{8.0 cm}
\includegraphics[height=5cm,width=4cm]{graphics/C5-quotient-of-d3reg10q.png}
\end{center}
\end{minipage}
\caption{The cubic graph  $\Gamma_2$ ({\tt d3reg10q}), with a degree two harmonic quotient morphism to and quasicover of, the cycle graph, $\Gamma_1=C_5$.}
\label{fig:harmonic-quotient_d3reg10q+C5}
\end{figure}

}
\end{example}




\begin{example}% d3reg10r
  {\tt d3reg10r}\newline
{\rm  
Let $\Gamma_2$ have vertices

\[
V_2=\{0,1,\dots,9\}
\]
and edges

\[
E_2 = \{
(0, 1), (0, 6), (0, 9), (1, 2), (1, 5), (2, 3), (2, 9), (3, 4), 
\]
\[
(3, 8), (4, 5), (4, 7), (5, 6), (6, 7), (7, 8), (8, 9)
\}.
\]
This graph has $8$ non-constant harmonic quotient
morphisms $\pi: \Gamma_2\to \Gamma_2/G$, for some subgroup
$G$ of $Aut(\Gamma_2)$.

Five of these quotient graphs are the path graph $P_3$
(with distinct quotient morphisms), 
however, the two groups that yield $P_2$ as a quotient have the
same quotient morphism.
Finally, the
group

\[
G = \langle (0,2)(1,9)(3,6)(4,7)(5,8), (0,9)(1,2)(3,5)(6,8) \rangle
\]
yields the
harmonic quotient morphism to the cycle graph
$\Gamma_1 = (V_1,E_1) = C_5$
with vertices and edges,

\[
V_1=\{0,1,\dots, 4\},
\]
\[
E_1=\{
(0, 3), (0, 4), (1, 2), (1, 4), (2, 3)
\}.
\]
This is tabulated below (the same table as in the previus example).

\begin{center}
\begin{tabular}{c|c|c|c}  
  $x$ & $\phi(x)$ & $m_\phi(x)$ & $\nu_\phi(x)$ \\ \hline
  0   &    2   &        1         &  1   \\ 
  1   &    2   &        1         &  1   \\ 
  2   &    3   &        1         &  1   \\ 
  3   &    0   &        1         &  1   \\ 
  4   &    4   &        1         &  1   \\ 
  5   &    1   &        1         &  1    \\ 
  6   &    1   &        1         &  1   \\ 
  7   &    4   &        1         &  1   \\
  8   &    0   &        1         &  1   \\
  9   &    3   &        1         &  1   \\ 
\end{tabular}
\end{center}

\[
  {\rm genus}(\Gamma_2) =6,\ \ \ \ \
  {\rm genus}(\Gamma_1) =1,
\]

\[
deg(\phi) =2.
\]

}
\end{example}



\begin{example}% d3reg10s
  {\tt d3reg10s}\newline
{\rm  
Let $\Gamma_2$ be the Petersen graph with vertices

\[
V_2=\{0,1,\dots,9\}
\]
and edges

\[
E_2 = \{
(0, 1), (0, 4), (0, 5), (1, 2), (1, 6), (2, 3), (2, 7), (3, 4),
\]
\[
(3, 8), (4, 9), (5, 7), (5, 8), (6, 8), (6, 9), (7, 9)
\}.
\]
This graph has $27$ non-constant harmonic quotient
morphisms $\pi: \Gamma_2\to \Gamma_2/G$, for some subgroup
$G$ of $Aut(\Gamma_2)$. Five of these quotient graphs are 
the star graph with $4$ vertices, $St_4$, while the remaining
graphs are the path graph, $P_2$.
\index{star graph}

The group
$G = \langle (0,4)(1,3)(5,9)(6,8), (0,4)(1,9)(2,7)(3,5) \rangle$ yields the
harmonic quotient morphism to the genus $0$ graph
$\Gamma_1 = (V_1,E_1) = St_4$, with

\[
V_1=\{0, 1, 2, 3\},
\]
and $3$ is the degree $3$ (``center'') vertex.
This is tabulated below.

\begin{center}
\begin{tabular}{c|c|c|c}  
  $x$ & $\phi(x)$ & $m_\phi(x)$ & $\nu_\phi(x)$ \\ \hline
  0   &    1   &        2         &  1   \\ 
  1   &    1   &        2         &  1   \\ 
  2   &    3   &        1         &  0   \\ 
  3   &    0   &        2         &  1   \\ 
  4   &    3   &        1         &  0   \\ 
  5   &    3   &        1         &  0    \\ 
  6   &    3   &        1         &  0   \\ 
  7   &    2   &        2         &  1   \\
  8   &    0   &        2         &  1   \\
  9   &    2   &        2         &  1   \\ 
\end{tabular}
\end{center}

\[
  {\rm genus}(\Gamma_2) =6,\ \ \ \ \
  {\rm genus}(\Gamma_1) =0,
\]

\[
deg(\phi) =4.
\]

}
\end{example}




\subsection{$12$ vertices}
\label{sec:cubic-12-vertices}
%%% file-of-d3reg12-graphs.txt and  file-of-3reg12-graphs.txt

There are are $85$ non-isomorphic, connected, cubic graphs on $12$ vertices,
$80$ of which are Hamiltonian and $4$ of which are vertex transitive.
However, none of the $85$ graphs
in this collection are edge transitive.

Here are several examples,
including those $4$ which are vertex transitive and
those $4$ that are not Hamiltonian.


{\small{
\begin{center}
\begin{tabular}{c|c|c|c}
  graph label & LCF & $|Aut(\Gamma)|$ & genus  \\ \hline
  d3reg12a & -- & 16 & 7   \\
  d3reg12b & -- & 16 & 7  \\
  d3reg12c & -- & 8 & 7   \\
  d3reg12d & -- & 32 & 7  \\
  d3reg12e & (12, [3, -2, -4, -3, 4, 2], 2) & 4 & 7  \\
  d3reg12f & (12, [3,-2,-4,-3,3,3,3,-3,-3,-3,4,2], 1) & 8 & 7   \\
  d3reg12g & (12, [3, 3, 3, -3, -3, -3], 2) & 4 & 7  \\
  d3reg12h & (12, [3, 3, 3, -3, -3, -3], 2) & 64 & 7  \\
  d3reg12i & (12, [2,-3,-2,3,3,3,-3,-3,-3,2,3,-2], 1) & 16 & 7   \\
  d3reg12j & (12, [2, 3, -2, 2, -3, -2], 2) & 16 & 7  \\
  d3reg12k & (12, [-2,3,6,3,-3,2,-3,-2,6,2,2,-2], 1) & 2 & 7  \\
  d3reg12l & (12, [6, -2, 2], 4) & 24 & 7  \\
  d3reg12m & (12, [-3, 3], 6) & 24 & 7  \\
  d3reg12n & (12, [5, -5], 6) & 48 & 7  \\
  
\end{tabular}
\end{center}  
}}



\begin{enumerate}
\item
  {\tt d3reg12a}\newline % or #83
  The one with vertices and edges:

\[
V=\{ 0,1,2,3,4,5,6,7,8,9,10,11\},
\]
\[
E = \{
(0, 1), (0, 2), (0, 9), (1, 2), (1, 3), (2, 3), (3, 9), (4, 5), (4, 6), (4, 7),
\]
\[
 (5, 8), (5, 11), (6, 7), (6, 10), (7, 10), (8, 9), (8, 11), (10, 11)
\},
\]
has automorphism group of order $16$ and
characteristic polynomial  
\[
x^{12} - 18x^{10} - 10x^9 + 109x^8 + 112x^7 - 223x^6 - 326x^5 + 58x^4 + 196x^3 + 9x^2 - 36x.
\]
It has girth $3$ and diameter $6$.
This is {\it not} a Hamiltonian graph, with circumference $7$.
The cardinality of its critical group is $2112=2^63\cdot 11$.

It has no non-trivial harmonic morphisms to $St_4$ or $P_4$, $C_3$, $C_4$,
or to $K_4$.
However, there are $12$ non-trivial harmonic morphisms to $P_3$.

%At this time, the number of non-trivial harmonic morphisms to $K_4$ is unknown.


\item
  {\tt d3reg12b}\newline % or #60
The one with vertices and edges:

\[
V=\{ 0,1,2,3,4,5,6,7,8,9,10,11\},
\]
\[
E = \{
(0, 1), (0, 9), (0, 10), (1, 4), (1, 5), (2, 3), (2, 6), (2, 11), (3, 6), 
\]
\[
(3, 9), (4, 5), (4, 10), (5, 10), (6, 7), (7, 8), (7, 11), (8, 9), (8, 11)
\}
\]
has automorphism group of order $16$ and
characteristic polynomial  
\[
x^{12} - 18x^{10} - 8x^9 + 109x^8 + 84x^7 - 240x^6 - 220x^5 + 172x^4 + 168x^3 .
\]
It has girth $3$ and diameter $5$.
This is {\it not} a Hamiltonian graph, with circumference $7$.
The size of the critical group of this graph is $2520=2^33^25\cdot 7$.

It has no non-trivial harmonic morphisms to $St_4$, $P_4$, $C_4$, or to $K_4$.
However, it has $14$ non-trivial harmonic morphisms to $P_3$ and
$6$ non-trivial harmonic morphisms to $C_3$.

\item
  {\tt d3reg12c}\newline % or #41
  The one with vertices and edges:

\[
V=\{ 0,1,2,3,4,5,6,7,8,9,10,11\},
\]
\[
E = \{
(0, 1), (0, 10), (0, 11), (1, 3), (1, 4), (2, 3), (2, 5), (2, 11), (3, 5),
\]
\[
(4, 5), (4, 11), (6, 7), (6, 8), (6, 10), (7, 8), (7, 9), (8, 9), (9, 10)
\}.
\]
has automorphism group of order $8$ and
characteristic polynomial  
\[
x^{12} - 18x^{10} - 6x^9 + 105x^8 + 60x^7 - 211x^6 - 122x^5 + 146x^4 + 52x^3 - 39x^2 .
\]
It has girth $3$ and diameter $6$.
This is {\it not} a Hamiltonian graph, with circumference $7$.
The size of the critical group of this graph is $2640=2^43\cdot 5\cdot 11$.

It has no non-trivial harmonic morphisms to $St_4$, $P_4$, $C_3$, $C_4$, or to $K_4$.
However, it has $6$ non-trivial harmonic morphisms to $P_3$.


\item
  {\tt d3reg12d}\newline % or 40
The one with vertices and edges:

\[
V=\{ 0,1,2,3,4,5,6,7,8,9,10,11\},
\]
\[
E = \{
(0, 1), (0, 10), (0, 11), (1, 3), (1, 4), (2, 3), (2, 4), (2, 11), (3, 5),
\]
\[
(4, 5), (5, 11), (6, 7), (6, 8), (6, 10), (7, 8), (7, 9), (8, 9), (9, 10)
\}.
\]
has automorphism group of order $32$ and
characteristic polynomial  
\[
x^{12} - 18x^{10} - 4x^9 + 101x^8 + 36x^7 - 176x^6 - 40x^5 + 84x^4 .
\]
It has girth $3$ and diameter $5$.
This is {\it not} a Hamiltonian graph, with circumference $7$.
The size of the critical group of this graph is $2808=2^33^313$.

It has $12$ non-trivial harmonic morphisms to $St_4$,
$22$ non-trivial harmonic morphisms to $P_3$, 
$12$ non-trivial harmonic morphisms to $C_3$, but
no non-trivial harmonic morphisms to $P_4$, $C_4$, or to $K_4$.


\item
  {\tt d3reg12e}\newline % or #33
 The one with automorphism group of order $4$ and
characteristic polynomial  
\[
x^{12} - 18x^{10} - 4x^9 + 109x^8 + 36x^7 - 256x^6 - 64x^5 + 228x^4 + 16x^3 - 48x^2 
\]
has girth $3$ and diameter $5$.
This is a Hamiltonian graph (with circumference $12$).
The size of the critical group of this graph is $5250=2\cdot 3\cdot 5^37$.

It has no non-trivial harmonic morphisms to $St_4$, $P_4$, $C_4$, or to $K_4$.
However, it has $4$ non-trivial harmonic morphisms to $P_3$ and
$6$ non-trivial harmonic morphisms to $C_3$.

SageMath notation:

\verb@graphs.LCFGraph(12,[3, -2, -4, -3, 4, 2], 2)@


\item
  {\tt d3reg12f}\newline % or #32
The one with automorphism group of order $8$ and
having characteristic polynomial  
\[
x^{12} - 18x^{10} - 2x^9 + 103x^8 + 18x^7 - 201x^6 - 26x^5 + 105x^4 
\]
has girth $3$ and diameter $4$.
This is a Hamiltonian graph (with circumference $12$).
The size of the critical group of this graph is $5535=3^35\cdot 41$.

It has $12$ non-trivial harmonic morphisms to $P_3$,
$36$ non-trivial harmonic morphisms to $C_3$, and
$12$ non-trivial harmonic morphisms to $St_4$.
However, there were no
non-trivial harmonic morphisms to $P_4$, $C_4$, or to $K_4$.

SageMath notation:

\verb@graphs.LCFGraph(12,[3,-2,-4,-3,3,3,3,-3,-3,-3,4,2],1)@



\item
  {\tt d3reg12g}\newline % or #31
The one with automorphism group of order $4$ and
having characteristic polynomial  
\[
x^{12} - 18x^{10} - 6x^9 + 111x^8 + 62x^7 - 265x^6 - 166x^5 + 213x^4 + 93x^3 - 60x^2 
\]
has girth $3$ and diameter $4$.
This is a Hamiltonian graph (with circumference $12$).
The size of the critical group of this graph is $4875=3\cdot 5^313$.

It has $12$ non-trivial harmonic morphisms to $St_4$,
$6$ non-trivial harmonic morphisms to $C_3$, but no
non-trivial harmonic morphisms to $P_3$, $P_4$, $C_4$, or to $K_4$.

SageMath notation:

\verb@graphs.LCFGraph(12,[4,2,3,-2,-4,-3,2,3,-2,2,-3,-2],1)@


\item
  {\tt d3reg12h}\newline % or #4
The one with automorphism group of order $64$ and
characteristic polynomial  
\[
x^{12} - 18x^{10} + 97x^8  - 144x^6 
\]
has girth $4$ and diameter $4$.
This is a bipartite, Hamiltonian graph (with circumference $12$).
The size of the critical group of this graph is $5832=2^33^6$.

It
has\footnote{Not only is $168$ some sort of record for
  the data we are looking at here, some of these are quasicovers
  (where the horizontal multiplicity, as a function of the
  vertices of $\Gamma$, is the constant $1$). One of these has a weed
  with $6$ roots.}
$168$ non-trivial harmonic morphisms to $C_3$,
$48$ non-trivial harmonic morphisms to $St_4$,
$56$ non-trivial harmonic morphisms to $P_3$, but no
non-trivial harmonic morphisms to $P_4$, $C_4$, or to $K_4$.

SageMath notation:

\verb@graphs.LCFGraph(12, [3, 3, 3, -3, -3, -3], 2)@


\item
  {\tt d3reg12i}\newline % or #3
The one with automorphism group of order $16$ and
having characteristic polynomial  

\[
x^{12} - 18x^{10} - 4x^9 + 105x^8 + 44x^7 - 216x^6 - 104x^5 + 96x^4 
\]
has girth $3$ and diameter $4$.
This is a Hamiltonian graph (with circumference $12$).
The size of the critical group of this graph is $5130=2\cdot 3^35\cdot 19$.

It has $42$ non-trivial harmonic morphisms to $C_3$,
$12$ non-trivial harmonic morphisms to $P_3$,
$12$ non-trivial harmonic morphisms to $St_4$,
but no
non-trivial harmonic morphisms to $P_4$, $C_4$, or to $K_4$.

SageMath notation:

\verb@graphs.LCFGraph(12,[2,-3,-2,3,3,3,-3,-3,-3,2,3,-2],1)@


\item
  {\tt d3reg12j}\newline % or #28
The one with automorphism group of order $16$ and
having characteristic polynomial  
\[
x^{12} - 18x^{10} - 8x^9 + 113x^8 + 88x^7 - 272x^6 - 272x^5 + 176x^4 + 192x^3
\]
has girth $3$ and diameter $4$.
This is a Hamiltonian graph (with circumference $12$).
The size of the critical group of this graph is $4500=2^23^25^3$.

It has $24$ non-trivial harmonic morphisms to $P_3$,
$24$ non-trivial harmonic morphisms to $C_3$,
$2$ non-trivial harmonic morphisms to $P_4$, but no
non-trivial harmonic morphisms to $St_4$, $C_4$, or to $K_4$.
By a computer check using SageMath we know there
are no non-trivial harmonic morphisms to $P_5$.
%Using Theorem \ref{thrm:diameter-bound}
Therefore, we know there
are no non-trivial harmonic morphisms to $P_k$, for $k\geq 6$.

SageMath notation:

\verb@graphs.LCFGraph(12, [2, 3, -2, 2, -3, -2], 2)@


\item
  {\tt d3reg12k}\newline % or #66
The one with automorphism group of order $2$ and
having characteristic polynomial  

\[
x^{12} - 18x^{10} - 6x^9 + 111x^8 + 68x^7 - 275x^6 - 220x^5 + 257x^4 + 236x^3 - 61x^2 - 54x + 9
\]
has girth $3$ and diameter $4$.
This is a Hamiltonian graph (with circumference $12$).
The size of the critical group of this graph is $5104=2^411\cdot 29$.

It has no non-trivial harmonic morphisms to $St_4$, $C_3$, $C_4$, $P_4$, or to $K_4$.
However, it has $4$ non-trivial harmonic morphisms to $P_3$.

%At this time, the number of non-trivial harmonic morphisms to $K_4$ is unknown.

SageMath notation:

\verb@graphs.LCFGraph(12,[-2,3,6,3,-3,2,-3,-2,6,2,2,-2],1)@


\item
  {\tt d3reg12l}\newline %{\tt d3reg12-17}
The truncated tetrahedron graph with automorphism group of order $24$
(a group isomorphic to $S_4$ acting transitively on the vertices)
has characteristic polynomial  
\index{truncated tetrahedron graph}
  
\[
x^{12} - 18x^{10} - 8x^9 + 117x^8 + 96x^7 - 316x^6 - 384x^5 + 240x^4 + 512x^3 - 192x^2 \, .
\]
This graph has chromatic number $3$ and is not edge transitive.
It has girth $3$, diameter $3$, and
is a Hamiltonian graph (with circumference $12$).
The size of the critical group of this graph is $6000=2^43\cdot 5^3$.

This graph has $12$ non-trivial harmonic morphisms to $P_3$,
$6$ non-trivial harmonic morphisms to $C_3$,
$6$ non-trivial harmonic morphisms to $P_4$,
$24$ non-trivial harmonic morphisms\footnote{All 24 of
  these non-trivial harmonic morphisms of $\Gamma_2\to\Gamma_1=K_4$
  are covers of degree $3$
  and are permuted amongst each other by the automorphism group.}
to $K_4$, but
no non-trivial harmonic morphisms to $C_4$ or to $St_4$.
All the non-trivial harmonic morphisms to $K_4$ are covers
and they are permuted transitively under the action of the
automorphism group of the graph.

SageMath notation:

\verb@graphs.LCFGraph(12, [6, -2, 2], 4)@


\item
  {\tt d3reg12m}\newline %{\tt d3reg12-11}
The graph, $\Gamma$ say, with automorphism group of order $24$
(acting transitively on the vertices) has characteristic polynomial  

\[
x^{12} - 18x^{10} + 105x^8 - 236x^6 + 144x^4  \, .
\]
It is an integral graph\footnote{Its adjacency spectrum
multi-set  is $\{3, 2, 2, 1, 0, 0, 0, 0, -1, -2, -2, -3\}$.},
with girth $4$, diameter $4$, and
circumference $12$ (i.e., is a Hamiltonian graph). 
It is bipartite but not edge transitive.
The size of the critical group of this graph
is\footnote{Note there is another
cubic graph, $\Gamma'$ say, on $12$ vertices, not presented here,
having this size critical group but trivial automorphism group:
$|Aut(\Gamma')|=1$. In particular, $\Gamma'$ is not a transitive graph.}
$8100=2^23^45^2$.
\index{integral graph}

There are no non-trivial harmonic morphisms to $C_4$ or to $K_4$.
However, there are $30$ non-trivial harmonic morphisms to $P_3$
$6$ non-trivial harmonic morphisms to $P_4$,
$12$ non-trivial harmonic morphisms to $St_4$, and
$36$ non-trivial harmonic morphisms\footnote{All these
  morphisms are quasicovers, with constant horizontal multiplicity
  $1$ and constant vertical multiplicity $1$.} to $C_3$.
By a computer check using SageMath we know there
are no non-trivial harmonic morphisms to $P_5$.
%Using Theorem \ref{thrm:diameter-bound}
Therefore, we know there
are no non-trivial harmonic morphisms to $P_k$, for $k\geq 6$.

SageMath notation:

\verb@graphs.LCFGraph(12, [-3, 3], 6)@

\item
 {\tt d3reg12n}\newline %  {\tt d3reg12-2}
The Franklin graph with automorphism group of order $48$
(acting transitively on the vertices but not on the edges)
has characteristic polynomial  
\index{Franklin graph}

\[
x^{12} - 18x^{10} + 111x^8 - 316x^6 + 447x^4 - 306x^2 + 81 \, .
\]
It has girth $4$, diameter $3$, and is a bipartite,
Hamiltonian graph (with circumference $12$).
The size of the critical group of this graph is $9216=2^{10}\cdot 3^2$.

There are $42$ non-trivial harmonic morphisms to $P_3$,
$18$ non-trivial harmonic morphisms to $St_4$,
but no non-trivial harmonic morphisms to $C_3$, $C_4$, $P_4$ or to $K_4$.

%At this time, the number of non-trivial harmonic morphisms to $K_4$ is unknown.

SageMath notation:

\verb@graphs.FranklinGraph()@

\noindent
or

\verb@graphs.LCFGraph(12, [5, -5], 6)@

%\item
  % #65
  
\end{enumerate}  




\subsubsection{Examples}
\label{sec:d3reg12-examples}

\begin{example}% d3reg12a
  {\tt d3reg12a}\newline
{\rm  
Let $\Gamma_2$ have vertices

\[
V_2=\{0,1,\dots,11\}
\]
and edges

\[
E_2 = \{
(0, 1), (0, 2), (0, 9), (1, 2), (1, 3), (2, 3), (3, 9), (4, 5), (4, 6), (4, 7), 
\]
\[
(5, 8), (5, 11), (6, 7), (6, 10), (7, 10), (8, 9), (8, 11), (10, 11)
\}.
\]
This graph has no non-constant harmonic quotient
morphisms $\pi: \Gamma_2\to \Gamma_2/G$, for some subgroup
$G$ of $Aut(\Gamma_2)$.

}
\end{example}


\begin{example}% d3reg12b
  {\tt d3reg12b}\newline
{\rm  
Let $\Gamma_2$ have vertices

\[
V_2=\{0,1,\dots,11\}
\]
and edges

\[
E_2 = \{
(0, 1), (0, 9), (0, 10), (1, 4), (1, 5), (2, 3), (2, 6),
(2, 11), (3, 6),
\]
\[
(3, 9), (4, 5), (4, 10), (5, 10), (6, 7), (7, 8), (7, 11), (8,
9), (8, 11)
\}.
\]
This graph has no non-constant harmonic quotient
morphisms $\pi: \Gamma_2\to \Gamma_2/G$, for some subgroup
$G$ of $Aut(\Gamma_2)$.

}
\end{example}



\begin{example}% d3reg12c
  {\tt d3reg12c}\newline
{\rm  
Let $\Gamma_2$ have vertices

\[
V_2=\{0,1,\dots,11\}
\]
and edges

\[
E_2 = \{
(0, 1), (0, 10), (0, 11), (1, 3), (1, 4), (2, 3), (2, 5),
(2, 11), (3, 5),
\]
\[
(4, 5), (4, 11), (6, 7), (6, 8), (6, 10), (7, 8), (7, 9), (8, 9), (9, 10)
\}.
\]
This graph has no non-constant harmonic quotient
morphisms $\pi: \Gamma_2\to \Gamma_2/G$, for some subgroup
$G$ of $Aut(\Gamma_2)$.

}
\end{example}



\begin{example}% d3reg12d
  {\tt d3reg12d}\newline
{\rm  
Let $\Gamma_2$ have vertices

\[
V_2=\{0,1,\dots,11\}
\]
and edges

\[
E_2 = \{
(0, 1), (0, 10), (0, 11), (1, 3), (1, 4), (2, 3), (2, 4),
(2, 11), (3, 5),
\]
\[
(4, 5), (5, 11), (6, 7), (6, 8), (6, 10), (7, 8), (7, 9), (8, 9), (9, 10)
\}.
\]
This graph has no non-constant harmonic quotient
morphisms $\pi: \Gamma_2\to \Gamma_2/G$, for some subgroup
$G$ of $Aut(\Gamma_2)$.

}
\end{example}



\begin{example}% d3reg12e
  {\tt d3reg12e}\newline
{\rm  
Let $\Gamma_2$ be the Petersen graph with vertices

\[
V_2=\{0,1,\dots,11\}
\]
and edges

\[
E_2 = \{
[(0, 1), (0, 3), (0, 11), (1, 2), (1, 11), (2, 3), (2, 10), (3, 4), (4, 5), (4, 8),
\]
\[
(5, 6), (5, 7), (6, 7), (6, 9), (7, 8), (8, 9), (9, 10), (10, 11)
\}.
\]
This graph has $2$ non-constant harmonic quotient
morphisms $\pi: \Gamma_2\to \Gamma_2/G$, for some subgroup
$G$ of $Aut(\Gamma_2)$. Both of these quotient graphs have $6$ vertices.

The group
$G = \langle (0,6)(1,7)(2,8)(3,9)(4,10)(5,11) \rangle$ yields the
harmonic quotient morphism (in fact, a cover) to the cubic prism graph
$\Gamma_1 = (V_1,E_1)$, with
\index{prism graph}

\[
V_1=\{0, 1, 2, 3, 4, 5\},
\]
\[
E_1=\{
(0, 1), (0, 2), (0, 5), (1, 3), (1, 5), (2, 3), (2, 4), (3, 4), (4, 5)
\}.
\]
This is tabulated below. The uncolored graphs are depicted in
Figure \ref{fig:harmonic-quotient_d3reg12e+prism6}.

\begin{center}
\begin{tabular}{c|c|c|c}  
  $x$ & $\phi(x)$ & $m_\phi(x)$ & $\nu_\phi(x)$ \\ \hline
  0   &    4   &        1         &  0   \\ 
  1   &    3   &        1         &  0   \\ 
  2   &    1   &        1         &  0   \\ 
  3   &    5   &        1         &  0   \\ 
  4   &    0   &        1         &  0   \\ 
  5   &    2   &        1         &  0    \\ 
  6   &    4   &        1         &  0   \\ 
  7   &    3   &        1         &  0   \\
  8   &    1   &        1         &  0   \\
  9   &    5   &        1         &  0   \\
  10   &   0   &        1         &  0   \\
  11  &    2   &        1         &  0   \\ 
\end{tabular}
\end{center}

\[
  {\rm genus}(\Gamma_2) =7,\ \ \ \ \
  {\rm genus}(\Gamma_1) =4,
\]

\[
deg(\phi) =2.
\]



\begin{figure}[H]
\begin{minipage}{\textwidth}
\begin{center}
%\vspace{8.0 cm}
\includegraphics[height=5cm,width=5cm]{graphics/graph_d3reg12e.png}
\end{center}
\begin{center}
%\vspace{8.0 cm}
\includegraphics[height=5cm,width=4cm]{graphics/prism6-quotient-of-d3reg12e.png}
\end{center}
\end{minipage}
\caption{The cubic graph  $\Gamma_2$ ({\tt d3reg12e}), with a degree two harmonic quotient morphism to and cover of, the cubic prism graph, $\Gamma_1$.}
\label{fig:harmonic-quotient_d3reg12e+prism6}
\end{figure}

The group
$G = \langle (0,5)(1,7)(2,8)(3,4)(6,11)(9,10) \rangle$ yields the
harmonic quotient morphism (in fact, a quasicover) to the graph
$\Gamma_1 = (V_1,E_1)$, with

\[
V_1=\{0, 1, 2, 3, 4, 5\},
\]
\[
E_1=\{
(0, 1), (0, 5), (1, 2), (1, 3), (2, 4), (2, 5), (3, 4), (4, 5)
\}.
\]
This is tabulated below.

\begin{center}
\begin{tabular}{c|c|c|c}  
  $x$ & $\phi(x)$ & $m_\phi(x)$ & $\nu_\phi(x)$ \\ \hline
  0   &    5   &        1         &  0   \\ 
  1   &    2   &        1         &  0   \\ 
  2   &    1   &        1         &  0   \\ 
  3   &    0   &        1         &  1   \\ 
  4   &    0   &        1         &  1   \\ 
  5   &    5   &        1         &  0    \\ 
  6   &    4   &        1         &  0   \\ 
  7   &    2   &        1         &  0   \\
  8   &    1   &        1         &  0   \\
  9   &    3   &        1         &  1   \\
  10   &   3   &        1         &  1   \\
  11  &    4   &        1         &  0   \\ 
\end{tabular}
\end{center}

\[
  {\rm genus}(\Gamma_2) =7,\ \ \ \ \
  {\rm genus}(\Gamma_1) =3,
\]

\[
deg(\phi) =2.
\]

}
\end{example}



\begin{example}% d3reg12f
  {\tt d3reg12f}\newline
{\rm  
Let $\Gamma_2$ have vertices

\[
V_2=\{0,1,\dots,11\}
\]
and edges

\[
E_2 = \{
(0, 1), (0, 3), (0, 11), (1, 2), (1, 11), (2, 3), (2, 10), (3, 4), (4, 5), (4, 7),
\]
\[
(5, 6), (5, 8), (6, 7), (6, 9), (7, 8), (8, 9), (9, 10), (10, 11)
\}.
\]
This graph has no non-constant harmonic quotient
morphisms $\pi: \Gamma_2\to \Gamma_2/G$, for some subgroup
$G$ of $Aut(\Gamma_2)$.

}
\end{example}


\begin{example}% d3reg12g
  {\tt d3reg12g}\newline
{\rm  
Let $\Gamma_2$ have vertices

\[
V_2=\{0,1,\dots,11\}
\]
and edges

\[
E_2 = \{
(0, 1), (0, 4), (0, 11), (1, 2), (1, 3), (2, 3), (2, 5), (3, 4), (4, 5), (5, 6), 
\]
\[
(6, 7), (6, 8), (7, 8), (7, 10), (8, 9), (9, 10), (9, 11), (10, 11)
\}.
\]
This graph has no non-constant harmonic quotient
morphisms $\pi: \Gamma_2\to \Gamma_2/G$, for some subgroup
$G$ of $Aut(\Gamma_2)$.

}
\end{example}



\begin{example}% d3reg12h
  {\tt d3reg12h}\newline
{\rm  
Let $\Gamma_2$ have vertices

\[
V_2=\{0,1,\dots,11\}
\]
and edges

\[
E_2 = \{
(0, 1), (0, 3), (0, 11), (1, 2), (1, 4), (2, 3), (2, 5), (3, 4), (4, 5), (5, 6),
\]
\[
 (6, 7), (6, 9), (7, 8), (7, 10), (8, 9), (8, 11), (9, 10), (10, 11)
\}.
\]
This graph has $52$ non-constant harmonic quotient
morphisms $\pi: \Gamma_2\to \Gamma_2/G$, for some subgroup
$G$ of $Aut(\Gamma_2)$.

Four of these $52$ are distinct non-constant harmonic quotient
morphisms $\pi: \Gamma_2\to \Gamma_2/G=K_{3,3}$, for some subgroup $G$.
(Indeed, these are all covers of this cubic graph.)
%\footnote{One way to visualize the complete bipartite graph $K_{3,3}$
%is to draw $C_6$ and then add ``interior'' edges connecting
%vertices on opposite sides of the circle.}.
\index{complete bipartite graph}
One of these, for the subgroup

\[
G_0=\langle (0,6)(1,7)(2,10)(3,9)(4,8)(5,11) \rangle,
\]
is tabulated below. The uncolored graphs are depicted in
Figure \ref{fig:harmonic-quotient_d3reg12h+K33}.

\begin{center}
\begin{tabular}{c|c|c|c}  
  $x$ & $\phi(x)$ & $m_\phi(x)$ & $\nu_\phi(x)$ \\ \hline
  0   &    3   &        1         &  0   \\ 
  1   &    2   &        1         &  0   \\ 
  2   &    1   &        1         &  0   \\ 
  3   &    4   &        1         &  0   \\ 
  4   &    5   &        1         &  0   \\ 
  5   &    0   &        1         &  0    \\ 
  6   &    3   &        1         &  0   \\ 
  7   &    2   &        1         &  0   \\
  8   &    5   &        1         &  0   \\
  9   &    4   &        1         &  0   \\
  10  &    1   &        1         &  0   \\
  11  &    0   &        1         &  0   \\ 
\end{tabular}
\end{center}

\[
  {\rm genus}(\Gamma_2) =7,\ \ \ \ \
  {\rm genus}(\Gamma_1) =4,
\]

\[
deg(\phi) =2.
\]


\begin{figure}[H]
\begin{minipage}{\textwidth}
\begin{center}
%\vspace{8.0 cm}
\includegraphics[height=5cm,width=5cm]{graphics/graph_d3reg12h.png}
\end{center}
\begin{center}
%\vspace{8.0 cm}
\includegraphics[height=5cm,width=4cm]{graphics/K33-quotient-of-d3reg12h.png}
\end{center}
\end{minipage}
\caption{The cubic graph  $\Gamma_2$ ({\tt d3reg12h}), with a degree two harmonic quotient morphism to and cover of, the complete bipartite graph, $K_{3,3}=\Gamma_2/G_0$.}
\label{fig:harmonic-quotient_d3reg12h+K33}
\end{figure}


Four of the remaining $48$ are distinct non-constant harmonic quotient
morphisms $\pi: \Gamma_2\to \Gamma_2/G=\Gamma_1$, for some $G$.
(Indeed, these are all quasicovers of this
graph\footnote{One way to visualize this graph is to
remove one of the ``interior'' edges of $K_{3,3}$.}
$\Gamma_1$.
One of these quotients, for the subgroup

\[
G_1=\langle (0,11)(1,10)(2,7)(3,8)(4,9)(5,6) \rangle,
\]
is tabulated below. The uncolored graphs are depicted in
Figure \ref{fig:harmonic-quotient_d3reg12h+K33m1}.

\begin{center}
\begin{tabular}{c|c|c|c}  
  $x$ & $\phi(x)$ & $m_\phi(x)$ & $\nu_\phi(x)$ \\ \hline
  0   &    2   &        1         &  1   \\ 
  1   &    0   &        1         &  0   \\ 
  2   &    4   &        1         &  0   \\ 
  3   &    5   &        1         &  0   \\ 
  4   &    3   &        1         &  0   \\ 
  5   &    1   &        1         &  1    \\ 
  6   &    1   &        1         &  1   \\ 
  7   &    4   &        1         &  0   \\
  8   &    0   &        1         &  0   \\
  9   &    3   &        1         &  0   \\
  10  &    5   &        1         &  0   \\
  11  &    2   &        1         &  1   \\ 
\end{tabular}
\end{center}

\[
  {\rm genus}(\Gamma_2) =7,\ \ \ \ \
  {\rm genus}(\Gamma_1) =3,
\]

\[
deg(\phi) =2.
\]

\begin{figure}[H]
\begin{minipage}{\textwidth}
\begin{center}
%\vspace{8.0 cm}
\includegraphics[height=5cm,width=5cm]{graphics/graph_d3reg12h.png}
\end{center}
\begin{center}
%\vspace{8.0 cm}
\includegraphics[height=4cm,width=4cm]{graphics/K33m1-quotient-of-d3reg12h.png}
\end{center}
\end{minipage}
\caption{The cubic graph  $\Gamma_2$ ({\tt d3reg12h}), with a degree two harmonic quotient morphism to and quasicover of, $\Gamma_1=\Gamma_2/G_1$.}
\label{fig:harmonic-quotient_d3reg12h+K33m1}
\end{figure}

Four of the remaining $44$ are distinct non-constant harmonic quotient
morphisms $\pi: \Gamma_2\to \Gamma_2/G$, for some $G$, of the barbell
graph, $\Gamma_1=Barbell(3,0)$.
(Indeed, these are all quasicovers of this graph $\Gamma_1$.)
One of these quotients, for the subgroup
\index{barbell graph}

\[
G_2=\langle (0,5)(1,2)(3,4)(6,11)(7,8)(9,10) \rangle,
\]
is tabulated below. The uncolored graphs are depicted in
Figure \ref{fig:harmonic-quotient_d3reg12h+barbell30}.


\begin{center}
\begin{tabular}{c|c|c|c}  
  $x$ & $\phi(x)$ & $m_\phi(x)$ & $\nu_\phi(x)$ \\ \hline
  0   &    1   &        1         &  0   \\ 
  1   &    3   &        1         &  1   \\ 
  2   &    3   &        1         &  1   \\ 
  3   &    0   &        1         &  1   \\ 
  4   &    0   &        1         &  1   \\ 
  5   &    1   &        1         &  0    \\ 
  6   &    4   &        1         &  0   \\ 
  7   &    5   &        1         &  1   \\
  8   &    5   &        1         &  1   \\
  9   &    2   &        1         &  1   \\
  10  &    2   &        1         &  1   \\
  11  &    4   &        1         &  0   \\ 
\end{tabular}
\end{center}

\[
  {\rm genus}(\Gamma_2) =7,\ \ \ \ \
  {\rm genus}(\Gamma_1) =2,
\]

\[
deg(\phi) =2.
\]

\begin{figure}[H]
\begin{minipage}{\textwidth}
\begin{center}
%\vspace{8.0 cm}
\includegraphics[height=5cm,width=5cm]{graphics/graph_d3reg12h.png}
\end{center}
\begin{center}
%\vspace{8.0 cm}
\includegraphics[height=4cm,width=4cm]{graphics/barbell30-quotient-of-d3reg12h.png}
\end{center}
\end{minipage}
\caption{The cubic graph  $\Gamma_2$ ({\tt d3reg12h}), with a degree two harmonic quotient morphism to and quasicover of, $\Gamma_1=Barbell(3,0)=\Gamma_2/G_2$.}
\label{fig:harmonic-quotient_d3reg12h+barbell30}
\end{figure}

Eight of the remaining $40$ are non-constant harmonic quotient
morphisms $\pi: \Gamma_2\to \Gamma_2/G$, for some $G$, of the
cycle graph, $C_3$.
(Indeed, these are all quasicovers of this graph $\Gamma_1$.)
One of these quotients, for the subgroup\footnote{Note the $G_2$ is a normal subgroup of $G_3$. Also note that $C_3$ is a quotient graph of the barbell
graph $Barbell(3,0)$.}

\[
G_3=\langle (0,5)(1,2)(3,4)(6,11)(7,8)(9,10), (0,6)(1,9)(2,8)(3,7)(4,10)(5,11) \rangle,
\]
is tabulated below. The colored graphs are depicted in
Figure \ref{fig:harmonic-quotient_d3reg12h+C3}.

\begin{center}
\begin{tabular}{c|c|c|c}  
  $x$ & $\phi(x)$ & $m_\phi(x)$ & $\nu_\phi(x)$ \\ \hline
  0   &    1   &        1         &  1   \\ 
  1   &    2   &        1         &  1   \\ 
  2   &    2   &        1         &  1   \\ 
  3   &    0   &        1         &  1   \\ 
  4   &    0   &        1         &  1   \\ 
  5   &    1   &        1         &  1    \\ 
  6   &    1   &        1         &  1   \\ 
  7   &    0   &        1         &  1   \\
  8   &    2   &        1         &  1   \\
  9   &    2   &        1         &  1   \\
  10  &    0   &        1         &  1   \\
  11  &    1   &        1         &  1   \\ 
\end{tabular}
\end{center}

\[
  {\rm genus}(\Gamma_2) =7,\ \ \ \ \
  {\rm genus}(\Gamma_1) =1,
\]

\[
deg(\phi) =4.
\]

\begin{figure}[H]
\begin{minipage}{\textwidth}
\begin{center}
%\vspace{8.0 cm}
\includegraphics[height=5cm,width=5cm]{graphics/harmonic-quotient_d3reg12h-C3_122001102201.png}
\end{center}
\begin{center}
%\vspace{8.0 cm}
\includegraphics[height=4cm,width=4cm]{graphics/cycle3-012.png}
\end{center}
\end{minipage}
\caption{The cubic graph  $\Gamma_2$ ({\tt d3reg12h}), with a degree two harmonic quotient morphism to and quasicover of, $\Gamma_1=C_3=\Gamma_2/G_3$.}
\label{fig:harmonic-quotient_d3reg12h+C3}
\end{figure}

}
\end{example}



\begin{example}% d3reg12i
  {\tt d3reg12i}\newline
{\rm  
Let $\Gamma_2$ have vertices

\[
V_2=\{0,1,\dots,11\}
\]
and edges

\[
E_2 = \{
(0, 1), (0, 2), (0, 11), (1, 2), (1, 10), (2, 3), (3, 4), (3, 6), (4, 5),
\]
\[
(4, 7), (5, 6), (5, 8), (6, 7), (7, 8), (8, 9), (9, 10), (9, 11), (10, 11)
\}.
\]
This graph has $2$ non-constant harmonic quotient
morphisms $\pi: \Gamma_2\to \Gamma_2/G$, for some subgroup
$G$ of $Aut(\Gamma_2)$. Both are distinct non-constant harmonic quotient
morphisms $\pi: \Gamma_2\to \Gamma_2/G$, for some $G$, of the barbell
graph, $\Gamma_1=Barbell(3,0)$.
\index{barbell graph}
(Indeed, these are all quasicovers of this graph $\Gamma_1$.)
One of these quotients, for the subgroup

\[
G_2=\langle (0,11)(1,10)(2,9)(3,8)(4,5)(6,7) \rangle,
\]
is tabulated below.
\begin{center}
\begin{tabular}{c|c|c|c}  
  $x$ & $\phi(x)$ & $m_\phi(x)$ & $\nu_\phi(x)$ \\ \hline
  0   &    1   &        1         &  1   \\ 
  1   &    4   &        1         &  1   \\ 
  2   &    5   &        1         &  0   \\ 
  3   &    0   &        1         &  0   \\ 
  4   &    2   &        1         &  1   \\ 
  5   &    2   &        1         &  1    \\ 
  6   &    3   &        1         &  1   \\ 
  7   &    3   &        1         &  1   \\
  8   &    0   &        1         &  0   \\
  9   &    5   &        1         &  0   \\
  10  &    4   &        1         &  1   \\
  11  &    1   &        1         &  1   \\ 
\end{tabular}
\end{center}

\[
  {\rm genus}(\Gamma_2) =7,\ \ \ \ \
  {\rm genus}(\Gamma_1) =2,
\]

\[
deg(\phi) =2.
\]

}
\end{example}



\begin{example}% d3reg12j
  {\tt d3reg12j}\newline
{\rm  
Let $\Gamma_2$ have vertices

\[
V_2=\{0,1,\dots,11\}
\]
and edges

\[
E_2 = \{
(0, 1), (0, 2), (0, 11), (1, 2), (1, 4), (2, 3), (3, 4), (3, 5), (4, 5),
\]
\[
(5, 6), (6, 7), (6, 8), (7, 8), (7, 10), (8, 9), (9, 10), (9, 11), (10, 11)
\}.
\]
This graph has $14$ non-constant harmonic quotient
morphisms $\pi: \Gamma_2\to \Gamma_2/G$, for some subgroup
$G$ of $Aut(\Gamma_2)$.

Two of these are distinct non-constant harmonic quotient
morphisms $\pi: \Gamma_2\to \Gamma_2/G$, for some $G$, of
the (cubic) prism graph (in fact, both a
degree $2$ covers).
\index{prism graph}
One of these quotients, for the subgroup

\[
G_2=\langle (0,6)(1,7)(2,8)(3,9)(4,10)(5,11) \rangle,
\]
is tabulated below.
\begin{center}
\begin{tabular}{c|c|c|c}  
  $x$ & $\phi(x)$ & $m_\phi(x)$ & $\nu_\phi(x)$ \\ \hline
  0   &    4   &        1         &  0   \\ 
  1   &    3   &        1         &  0   \\ 
  2   &    1   &        1         &  0   \\ 
  3   &    5   &        1         &  0   \\ 
  4   &    0   &        1         &  0   \\ 
  5   &    2   &        1         &  0    \\ 
  6   &    4   &        1         &  0   \\ 
  7   &    3   &        1         &  0   \\
  8   &    1   &        1         &  0   \\
  9   &    5   &        1         &  0   \\
  10  &    0   &        1         &  0   \\
  11  &    2   &        1         &  0   \\ 
\end{tabular}
\end{center}

\[
  {\rm genus}(\Gamma_2) =7,\ \ \ \ \
  {\rm genus}(\Gamma_1) =4,
\]

\[
deg(\phi) =2.
\]


One of the remaining $12$ is a non-constant harmonic quotient
morphism $\pi: \Gamma_2\to \Gamma_2/G$, for some $G$,
for the barbell
graph, $\Gamma_1=Barbell(3,0)$.
\index{barbell graph}
This quotient, for the subgroup

\[
G=\langle  (0,5)(1,4)(2,3)(6,11)(7,10)(8,9) \rangle,
\]
is tabulated below.
\begin{center}
\begin{tabular}{c|c|c|c}  
  $x$ & $\phi(x)$ & $m_\phi(x)$ & $\nu_\phi(x)$ \\ \hline
  0   &    5   &        1         &  0   \\ 
  1   &    0   &        1         &  1   \\ 
  2   &    1   &        1         &  1   \\ 
  3   &    1   &        1         &  1   \\ 
  4   &    0   &        1         &  1   \\ 
  5   &    5   &        1         &  0    \\ 
  6   &    4   &        1         &  0   \\ 
  7   &    2   &        1         &  1   \\
  8   &    3   &        1         &  1   \\
  9   &    3   &        1         &  1   \\
  10  &    2   &        1         &  1   \\
  11  &    4   &        1         &  0   \\ 
\end{tabular}
\end{center}

\[
  {\rm genus}(\Gamma_2) =7,\ \ \ \ \
  {\rm genus}(\Gamma_1) =2,
\]

\[
deg(\phi) =2.
\]

Two of the remaining $11$ are distinct non-constant harmonic quotient
morphisms $\pi: \Gamma_2\to \Gamma_2/G$, for some $G$, of
codomino graph.
\index{codomino graph}
One of these quotients, for the subgroup

\[
G=\langle (0,11)(1,9)(2,10)(3,7)(4,8)(5,6) \rangle,
\]
is tabulated below. The uncolored graphs are depicted in
Figure \ref{fig:harmonic-quotient_d3reg12j+codomino}.

\begin{center}
\begin{tabular}{c|c|c|c}  
  $x$ & $\phi(x)$ & $m_\phi(x)$ & $\nu_\phi(x)$ \\ \hline
  0   &    3   &        1         &  1   \\ 
  1   &    1   &        1         &  0  \\ 
  2   &    2   &        1         &  0  \\ 
  3   &    4   &        1         &  0  \\ 
  4   &    5   &        1         &  0   \\ 
  5   &    0   &        1         &  1    \\ 
  6   &    0   &        1         &  1   \\ 
  7   &    4   &        1         &  0   \\
  8   &    5   &        1         &  0   \\
  9   &    1   &        1         &  0   \\
  10  &    2   &        1         &  0   \\
  11  &    3   &        1         &  1   \\ 
\end{tabular}
\end{center}

\[
  {\rm genus}(\Gamma_2) =7,\ \ \ \ \
  {\rm genus}(\Gamma_1) =3,
\]

\[
deg(\phi) =2.
\]

\begin{figure}[H]
\begin{minipage}{\textwidth}
\begin{center}
%\vspace{8.0 cm}
\includegraphics[height=5cm,width=5cm]{graphics/graph_d3reg12j.png}
\end{center}
\begin{center}
%\vspace{8.0 cm}
\includegraphics[height=4cm,width=4cm]{graphics/codomino-quotient-of-d3reg12j.png}
\end{center}
\end{minipage}
\caption{The cubic graph  $\Gamma_2$ ({\tt d3reg12j}), with a degree two harmonic quotient morphism to and quasicover of the codomino graph, $\Gamma_1=\Gamma_2/G_2$.}
\label{fig:harmonic-quotient_d3reg12j+codomino}
\end{figure}

We tabulate the subgroups and their (harmonic) quotients
$\Gamma_1=\Gamma_2/G$. 

{\small{
\begin{center}
\begin{tabular}{c|c}
  gens of $G\subset Aut(\Gamma_2)$ & $\Gamma_2/G$ \\ \hline
  (0,6)(1,7)(2,8)(3,9)(4,10)(5,11) & prism \\ \hline
  (0,6)(1,8)(2,7)(3,10)(4,9)(5,11) & prism \\ \hline
  (0,5)(1,4)(2,3)(6,11)(7,10)(8,9) & $Barbell(3,0)$ \\ \hline
  (0,11)(1,9)(2,10)(3,7)(4,8)(5,6) & codomino \\ \hline
  (0,11)(1,10)(2,9)(3,8)(4,7)(5,6) & codomino \\ \hline
  (0,5)(1,4)(2,3)(6,11)(7,10)(8,9), & \\
  (0,6)(1,7)(2,8)(3,9)(4,10)(5,11) & $C_3$ \\ \hline
  (0,5)(1,4)(2,3)(6,11)(7,10)(8,9), & \\
  (0,6)(1,8)(2,7)(3,10)(4,9)(5,11)  & $C_3$ \\ \hline
(1,2)(3,4)(7,8)(9,10), & \\
(0,11)(1,10,2,9)(3,8,4,7)(5,6)  & $P_4$ \\ \hline
  (1,2)(3,4)(7,8)(9,10), & \\
  (0,11)(1,9)(2,10)(3,7)(4,8)(5,6)  & $P_4$ \\ \hline
  (7,8)(9,10), (1,2)(3,4), & \\
  (0,11)(1,9)(2,10)(3,7)(4,8)(5,6)  & $P_4$ \\ \hline
  (1,2)(3,4)(7,8)(9,10),  & \\
  (0,5)(1,3)(2,4)(6,11)(7,9)(8,10), & \\
  (0,6)(1,7)(2,8)(3,9)(4,10)(5,11) & $P_2$ \\ \hline
  (1,2)(3,4)(7,8)(9,10), & \\
  (0,5)(1,3)(2,4)(6,11)(7,10)(8,9), & \\
  (0,6)(1,7)(2,8)(3,9)(4,10)(5,11) & $P_2$ \\ \hline
  (1,2)(3,4)(7,8)(9,10), & \\
  (0,5)(1,3)(2,4)(6,11)(7,9)(8,10), & \\
  (0,6)(1,8,2,7)(3,10,4,9)(5,11) & $P_2$ \\ \hline
  (1,2)(3,4)(7,8)(9,10), & \\
  (0,5)(1,3)(2,4)(6,11)(7,10)(8,9), & \\
  (0,6)(1,8,2,7)(3,10,4,9)(5,11) & $P_2$ \\
\end{tabular}
\end{center}
}}
All the quotient morphisms to $P_4$ in the above table are the same,
and all those to $P_2$ are the same.

Using this table one can see that

(a) One of the quotient morphisms to the prism graph has a further
quotient to $C_3$.

(b) One of the quotient morphisms to the codomino graph has a further
quotient to $P_4$.

(c) One of the quotient morphisms to the barbell graph has a further
quotient to $C_3$.
}
\end{example}



\begin{example}% d3reg12k
  {\tt d3reg12k}\newline
{\rm  
Let $\Gamma_2$ have vertices

\[
V_2=\{0,1,\dots,11\}
\]
and edges

\[
E_2 = \{
(0, 1), (0, 10), (0, 11), (1, 2), (1, 4), (2, 3), (2, 8), (3, 4), (3, 6),
\]
\[
(4, 5), (5, 6), (5, 7), (6, 7), (7, 8), (8, 9), (9, 10), (9, 11), (10, 11)
\}.
\]
This graph has no non-constant harmonic quotient
morphisms $\pi: \Gamma_2\to \Gamma_2/G$, for some subgroup
$G$ of $Aut(\Gamma_2)$.

}
\end{example}


\begin{example}% d3reg12l
  {\tt d3reg12l}\newline
{\rm  
Let $\Gamma_2$ have vertices

\[
V_2=\{0,1,\dots,11\}
\]
and edges

\[
E_2 = \{
(0, 1), (0, 6), (0, 11), (1, 2), (1, 11), (2, 3), (2, 4), (3, 4), (3, 9),
\]
\[
(4, 5), (5, 6), (5, 7), (6, 7), (7, 8), (8, 9), (8, 10), (9, 10), (10, 11)
\}.
\]
This graph has $10$ non-constant harmonic quotient
morphisms $\pi: \Gamma_2\to \Gamma_2/G$, for some subgroup
$G$ of $Aut(\Gamma_2)$.

Three of these $10$ are distinct non-constant harmonic quotient
morphisms $\pi: \Gamma_2\to \Gamma_2/G$, for some $G$, of
codomino graph.
\index{codomino graph}
One of these quotients, for the subgroup

\[
G_1=\langle (0,3)(1,2)(4,11)(5,10)(6,9)(7,8) \rangle,
\]
is tabulated below.

\begin{center}
\begin{tabular}{c|c|c|c}  
  $x$ & $\phi(x)$ & $m_\phi(x)$ & $\nu_\phi(x)$ \\ \hline
  0   &    4   &        1         &  0   \\ 
  1   &    3   &        1         &  1  \\ 
  2   &    3   &        1         &  1  \\ 
  3   &    4   &        1         &  0  \\ 
  4   &    1   &        1         &  0   \\ 
  5   &    2   &        1         &  0    \\ 
  6   &    0   &        1         &  0   \\ 
  7   &    5   &        1         &  1   \\
  8   &    5   &        1         &  1   \\
  9   &    0   &        1         &  0   \\
  10  &    2   &        1         &  0   \\
  11  &    1   &        1         &  0   \\ 
\end{tabular}
\end{center}

\[
  {\rm genus}(\Gamma_2) =7,\ \ \ \ \
  {\rm genus}(\Gamma_1) =3,
\]

\[
deg(\phi) =2.
\]

One of the remaining $7$ is a distinct non-constant harmonic quotient
morphism $\pi: \Gamma_2\to \Gamma_2/G$, for some $G$, of
cycle graph, $C_3$.
The quotient $\phi: \Gamma_2\to \Gamma_2/G_2$ for the subgroup

\[
G_2=\langle (0,3)(1,2)(4,11)(5,10)(6,9)(7,8),
(0,6)(1,7)(2,8)(3,9)(4,10)(5,11) \rangle,
\]
is tabulated below. Note that $G_1$ is a (normal) subgroup of $G_2$.
We conclude that one of the quotient morphisms to the
codomino graph has a further quotient to $C_3$.

\begin{center}
\begin{tabular}{c|c|c|c}  
  $x$ & $\phi(x)$ & $m_\phi(x)$ & $\nu_\phi(x)$ \\ \hline
  0   &    2   &        1         &  1   \\ 
  1   &    1   &        1         &  1  \\ 
  2   &    1   &        1         &  1  \\ 
  3   &    2   &        1         &  1  \\ 
  4   &    0   &        1         &  1   \\ 
  5   &    0   &        1         &  1    \\ 
  6   &    2   &        1         &  1   \\ 
  7   &    1   &        1         &  1   \\
  8   &    1   &        1         &  1   \\
  9   &    2   &        1         &  1   \\
  10  &    0   &        1         &  1   \\
  11  &    0   &        1         &  1   \\ 
\end{tabular}
\end{center}

\[
  {\rm genus}(\Gamma_2) =7,\ \ \ \ \
  {\rm genus}(\Gamma_1) =1,
\]

\[
deg(\phi) =4.
\]

Three of the remaining $6$ are distinct non-constant harmonic quotient
morphisms $\pi: \Gamma_2\to \Gamma_2/G$, for some $G$, of
the path graph, $P_4$.
The remaining $3$ are distinct non-constant harmonic quotient
morphisms $\pi: \Gamma_2\to \Gamma_2/G$, for some $G$, of
the path graph, $P_2$.

}
\end{example}




\begin{example}% d3reg12m
  {\tt d3reg12m}\newline
{\rm  
Let $\Gamma_2$ have vertices

\[
V_2=\{0,1,\dots,11\}
\]
and edges

\[
E_2 = \{
(0, 1), (0, 9), (0, 11), (1, 2), (1, 4), (2, 3), (2, 11), (3, 4), (3, 6),
\]
\[
(4, 5), (5, 6), (5, 8), (6, 7), (7, 8), (7, 10), (8, 9), (9, 10), (10, 11)
\}.
\]
This graph has $31$ non-constant harmonic quotient
morphisms $\pi: \Gamma_2\to \Gamma_2/G$, for some subgroup
$G$ of $Aut(\Gamma_2)$.

Three of these $31$ are distinct non-constant harmonic quotient
morphisms $\pi: \Gamma_2\to \Gamma_2/G$, for some $G$, of
codomino graph.
\index{codomino graph}

Three of the remaining $28$ are distinct non-constant harmonic quotient
morphisms $\pi: \Gamma_2\to \Gamma_2/G$, for some $G$, of
grid graph, $Grid(2,1)$.
\index{grid graph}

One of the remaining $25$ is a distinct non-constant harmonic quotient
morphism $\pi: \Gamma_2\to \Gamma_2/G$, for some $G$, of
complete bipartite graph, $K_{3,3}$.

One of the remaining $24$ is a distinct non-constant harmonic quotient
morphism $\pi: \Gamma_2\to \Gamma_2/G$, for some $G$, of
(cubic) prism graph.
\index{prism graph}

One of the remaining $23$ is a distinct non-constant harmonic quotient
morphism $\pi: \Gamma_2\to \Gamma_2/G$, for some $G$, of
cycle graph, $C_6$.
The quotient $\phi: \Gamma_2\to \Gamma_2/G$ for the subgroup

\[
G=\langle (0,11)(1,2)(3,4)(5,6)(7,8)(9,10) \rangle,
\]
is tabulated below. 

\begin{center}
\begin{tabular}{c|c|c|c}  
  $x$ & $\phi(x)$ & $m_\phi(x)$ & $\nu_\phi(x)$ \\ \hline
  0   &    2   &        1         &  1   \\ 
  1   &    4   &        1         &  1  \\ 
  2   &    4   &        1         &  1  \\ 
  3   &    0   &        1         &  1  \\ 
  4   &    0   &        1         &  1   \\ 
  5   &    1   &        1         &  1    \\ 
  6   &    1   &        1         &  1   \\ 
  7   &    5   &        1         &  1   \\
  8   &    5   &        1         &  1   \\
  9   &    3   &        1         &  1   \\
  10  &    3   &        1         &  1   \\
  11  &    2   &        1         &  1   \\ 
\end{tabular}
\end{center}

\[
  {\rm genus}(\Gamma_2) =7,\ \ \ \ \
  {\rm genus}(\Gamma_1) =1,
\]

\[
deg(\phi) =2.
\]

Three of the remaining $22$ are distinct non-constant harmonic quotient
morphisms $\pi: \Gamma_2\to \Gamma_2/G$, for some $G$, of
path graph, $P_4$.

Four of the remaining $19$ are distinct non-constant harmonic quotient
morphisms $\pi: \Gamma_2\to \Gamma_2/G$, for some $G$, of
cycle graph, $C_3$.

Three of the remaining $15$ are distinct non-constant harmonic quotient
morphisms $\pi: \Gamma_2\to \Gamma_2/G$, for some $G$, of
path graph, $P_3$.

The remaining $12$ are (not necessarily distinct) non-constant harmonic quotient
morphisms $\pi: \Gamma_2\to \Gamma_2/G$, for some $G$, of
path graph, $P_2$.


}
\end{example}



\begin{example}% d3reg12n
  {\tt d3reg12n}\newline
{\rm  
Let $\Gamma_2$ be the Franklin graph, with vertices
\index{Franklin graph}

\[
V_2=\{0,1,\dots,11\}
\]
and edges

\[
E_2 = \{
(0, 1), (0, 5), (0, 6), (1, 2), (1, 7), (2, 3), (2, 8), (3, 4), (3, 9), (4, 5), 
\]
\[
(4, 10), (5, 11), (6, 7), (6, 9), (7, 10), (8, 9), (8, 11), (10, 11)
\}.
\]
This graph has $35$ non-constant harmonic quotient
morphisms $\pi: \Gamma_2\to \Gamma_2/G$, for some subgroup
$G$ of $Aut(\Gamma_2)$.

Six of these $35$ are distinct non-constant harmonic quotient
morphisms $\pi: \Gamma_2\to \Gamma_2/G$, for some $G$,  to
the {\tt EUYo}
co-twin-$C_5$ graph\footnote{We are using the naming conventions
of graph6 \cite{graph6}. In the ISGCI it is the co-twin-$C_5$ graph
\cite{ISGCI}.}.
\index{co-twin-$C_5$ graph}
One of these quotients, for the subgroup

\[
G=\langle (0,1)(2,5)(3,11)(4,8)(6,7)(9,10) \rangle,
\]
is tabulated below. The uncolored graphs are depicted in
Figure \ref{fig:harmonic-quotient_d3reg12n+cotwinC5}.

\begin{center}
\begin{tabular}{c|c|c|c}  
  $x$ & $\phi(x)$ & $m_\phi(x)$ & $\nu_\phi(x)$ \\ \hline
  0   &    1   &        1         &  1   \\ 
  1   &    1   &        1         &  1  \\ 
  2   &    4   &        1         &  0  \\ 
  3   &    3   &        1         &  0  \\ 
  4   &    5   &        1         &  0   \\ 
  5   &    4   &        1         &  0  \\ 
  6   &    2   &        1         &  1   \\ 
  7   &    2   &        1         &  1   \\
  8   &    5   &        1         &  0   \\
  9   &    0   &        1         &  0   \\
  10  &    0   &        1         &  0   \\
  11  &    3   &        1         &  0   \\ 
\end{tabular}
\end{center}

\[
  {\rm genus}(\Gamma_2) =7,\ \ \ \ \
  {\rm genus}(\Gamma_1) =3,
\]

\[
deg(\phi) =2.
\]

\begin{figure}[H]
\begin{minipage}{\textwidth}
\begin{center}
%\vspace{8.0 cm}
\includegraphics[height=5cm,width=5cm]{graphics/graph_d3reg12n.png}
\end{center}
\begin{center}
%\vspace{8.0 cm}
\includegraphics[height=4cm,width=4cm]{graphics/cotwinC5-quotient-of-d3reg12n.png}
\end{center}
\end{minipage}
\caption{The cubic graph  $\Gamma_2$ ({\tt d3reg12n}), with a degree two harmonic quotient morphism to and quasicover of the codomino graph, $\Gamma_1=\Gamma_2/G_2$.}
\label{fig:harmonic-quotient_d3reg12n+cotwinC5}
\end{figure}

One of the remaining $29$ is a distinct non-constant harmonic quotient
morphism $\pi: \Gamma_2\to \Gamma_2/G$, for some $G$, of
the $X_{37}$ graph.
\index{$X_{37}$ graph}
One of these quotients, for the subgroup

\[
G=\langle (0,1)(2,5)(3,4)(6,7)(8,11)(9,10) \rangle,
\]
is tabulated below. The uncolored graphs are depicted in
Figure \ref{fig:harmonic-quotient_d3reg12n+X37}.

\begin{center}
\begin{tabular}{c|c|c|c}  
  $x$ & $\phi(x)$ & $m_\phi(x)$ & $\nu_\phi(x)$ \\ \hline
  0   &    4   &        1         &  1   \\ 
  1   &    4   &        1         &  1  \\ 
  2   &    5   &        1         &  0  \\ 
  3   &    0   &        1         &  1  \\ 
  4   &    0   &        1         &  1   \\ 
  5   &    5   &        1         &  0  \\ 
  6   &    1   &        1         &  1   \\ 
  7   &    1   &        1         &  1   \\
  8   &    2   &        1         &  1   \\
  9   &    3   &        1         &  0   \\
  10  &    3   &        1         &  0   \\
  11  &    2   &        1         &  1   \\ 
\end{tabular}
\end{center}

\[
  {\rm genus}(\Gamma_2) =7,\ \ \ \ \
  {\rm genus}(\Gamma_1) =2,
\]

\[
deg(\phi) =2.
\]

\begin{figure}[H]
\begin{minipage}{\textwidth}
\begin{center}
%\vspace{8.0 cm}
\includegraphics[height=5cm,width=5cm]{graphics/graph_d3reg12n.png}
\end{center}
\begin{center}
%\vspace{8.0 cm}
\includegraphics[height=4cm,width=4cm]{graphics/X37-quotient-of-d3reg12n.png}
\end{center}
\end{minipage}
\caption{The cubic graph  $\Gamma_2$ ({\tt d3reg12n}), with a degree two harmonic quotient morphism to and quasicover of the $X_{37}$ graph, $\Gamma_1=\Gamma_2/G_2$.}
\label{fig:harmonic-quotient_d3reg12n+X37}
\end{figure}

Five of the remaining $28$ are distinct non-constant harmonic quotient
morphisms $\pi: \Gamma_2\to \Gamma_2/G$, for some $G$, of
twin-$C_5$ graph.
\index{twin-$C_5$ graph}
One of these quotients, for the subgroup

\[
G=\langle (0,3)(1,2)(4,5)(6,9)(7,8)(10,11) \rangle,
\]
is tabulated below. The uncolored graphs are depicted in
Figure \ref{fig:harmonic-quotient_d3reg12n+twinC5}.

\begin{center}
\begin{tabular}{c|c|c|c}  
  $x$ & $\phi(x)$ & $m_\phi(x)$ & $\nu_\phi(x)$ \\ \hline
  0   &    2   &        1         &  0   \\ 
  1   &    1   &        1         &  1  \\ 
  2   &    1   &        1         &  1  \\ 
  3   &    2   &        1         &  0  \\ 
  4   &    3   &        1         &  1   \\ 
  5   &    3   &        1         &  1  \\ 
  6   &    0   &        1         &  1   \\ 
  7   &    5   &        1         &  0   \\
  8   &    5   &        1         &  0   \\
  9   &    0   &        1         &  1   \\
  10  &    4   &        1         &  1   \\
  11  &    4   &        1         &  1   \\ 
\end{tabular}
\end{center}

\[
  {\rm genus}(\Gamma_2) =7,\ \ \ \ \
  {\rm genus}(\Gamma_1) =2,
\]

\[
deg(\phi) =2.
\]

\begin{figure}[H]
\begin{minipage}{\textwidth}
\begin{center}
%\vspace{8.0 cm}
\includegraphics[height=5cm,width=5cm]{graphics/graph_d3reg12n.png}
\end{center}
\begin{center}
%\vspace{8.0 cm}
\includegraphics[height=4cm,width=4cm]{graphics/twinC5b-quotient-of-d3reg12n.png}
\end{center}
\end{minipage}
\caption{The cubic graph  $\Gamma_2$ ({\tt d3reg12n}), with a degree two harmonic quotient morphism to and quasicover of the twin-$C_5$ graph, $\Gamma_1=\Gamma_2/G_2$.}
\label{fig:harmonic-quotient_d3reg12n+twinC5}
\end{figure}

Nine of the remaining $23$ are (not necessarily distinct)
non-constant harmonic quotient
morphism $\pi: \Gamma_2\to \Gamma_2/G$, for some $G$, of
the star graph $St_4$ having $4$ vertices.
\index{star graph}

All of the remaining $14$ are (not necessarily distinct)
non-constant harmonic quotient
morphism $\pi: \Gamma_2\to \Gamma_2/G$, for some $G$, of
the path graph $P_2$.
}
\end{example}



\section{Quartic cases}

\subsection{$5$ vertices}

The complete graph $K_5$ is the only quartic graph on $5$ vertices.

Since a non-trivial harmonic morphism is surjective on the vertices and
since it sends neighbors to neighbors, there can be no non-trivial
harmonic morphism from $K_5$ (or any complete graph) to
$St_4$, $P_4$, or to $C_4$.
%By Theorem \ref{thm:harm-morphism-complete},
By general theory, we know
the number of non-trivial harmonic morphisms to $K_3=C_3$ or to $K_4$ is $0$. 


\subsubsection{Examples}

\begin{example}% d4reg5a
  {\tt d4reg5a}\newline
{\rm  
  Let $\Gamma_2=K_5$, in standard form.

There are $75$ harmonic quotient morphisms
$\pi: \Gamma_2\to \Gamma_1=\Gamma_2/G$,
for some $G$, but most are not distinct. Every quotient is
$\Gamma_1=P_2$, and there are no others.
An example for
\[
G = \langle (1,3)(2,4), (1,4)(2,3)\rangle
\]
is tabulated below.

\begin{center}
\begin{tabular}{c|c|c|c}  
  $x$ & $\phi(x)$ & $m_\phi(x)$ & $\nu_\phi(x)$ \\ \hline
  0   &    1   &        4         &  0   \\ 
  1   &    0   &        1         &  3  \\ 
  2   &    0   &        1         &  3  \\ 
  3   &    0   &        1         &  3  \\ 
  4   &    0   &        1         &  3   \\ 
\end{tabular}
\end{center}

\[
  {\rm genus}(\Gamma_2) =6,\ \ \ \ \
  {\rm genus}(\Gamma_1) =0,
\]

\[
deg(\phi) =4.
\]
}
\end{example}


\subsection{$6$ vertices}

There is (up to isomorphism) exactly one $4$-regular connected graph
on $6$ vertices, the octahedral graph (also known as the graph of the
Fano diagram).
\index{octahedral graph}

Define $\Gamma=(V,E)$ to have vertices $V=$
% d4reg6a

\[
V = \{0,1,2,3,4,5\},
\]
and edges $E=$
\[
\{(0, 1), (0, 2), (0, 4), (0, 5), (1, 2), (1, 3), (1, 4), (2, 3), (2, 5), (3, 4), (3, 5), (4, 5)\}.
\]
It is an integral graph\footnote{The multiset of eigenvalues of
  the adjacency matrix is $\{4, 0, 0, 0, -2, -2\}$.} with girth $3$,
diameter $2$, and circumference $6$ (i.e., is a
Hamiltonian graph).
\index{integral graph}
Its automorphism group is order $48$ and it's vertex transitive,
edge transitive, strongly regular, and distance regular.
The size of the critical group of this graph is $384=2^73$.
\index{vertex transitive graph}
\index{edge transitive graph}
\index{distance regular graph}
\index{strongly regular graph}

There are no non-trivial harmonic morphisms to $St_4$, $P_4$, $C_4$,
or to $K_4$.
There are exactly $6$ non-trivial harmonic morphisms $\Gamma\to P_3$ and
$6$ non-trivial harmonic
morphisms\footnote{All $6$ of these morphisms to $C_3$ are
  $2$-multicovers of degree $4$.} $\Gamma\to C_3$.



\subsubsection{Examples}

\begin{example}% d4reg6a
 {\tt d4reg6a}\newline
{\rm  
  Let $\Gamma_2$ denote the octahedral graph as given above.
There are $59$ harmonic quotient morphisms
$\pi: \Gamma_2\to \Gamma_1=\Gamma_2/G$,
for some $G$, but most are not distinct.

Six of these $59$ are distinct harmonic quotient morphisms
to the diamond graph.

Of the remaining $53$, $14$ are (not necessarily distinct)
harmonic quotient morphisms to the cycle graph $C_3$. As we
know from the paragraph above, all of these morphisms to $C_3$ are
$2$-multicovers.

The remaining $39$ are (not necessarily distinct)
harmonic quotient morphisms to the path graph $P_2$.
}
\end{example}



\subsection{$7$ vertices}

There are (up to isomorphism) exactly two $4$-regular connected graphs on $7$ vertices.

\begin{enumerate}
\item
  {\tt d4reg7a}
Define $\Gamma=(V,E)$ to have vertex vertex set $\{0,1,2,3,4,5,6\}$ and edge set:

\[
\{(0, 1), (0, 3), (0, 5), (0, 6), (1, 2), (1, 3), (1, 5), (2, 3), (2, 4),
\]
\[
(2, 6), (3, 4), (4, 5), (4, 6), (5, 6)\}.
\]
It has girth $3$, diameter $2$, and is a
Hamiltonian graph (with circumference $7$).
Its automorphism group is order $14$ and acts transitively on the vertices
(but it's not edge transitive).
The size of the critical group of this graph is $1183=7\cdot 13^2$.

There are no non-trivial harmonic morphisms to $St_4$, $P_3$, $P_4$, $C_3$, $C_4$ or to $K_4$.

\item
{\tt d4reg7b}
Define $\Gamma=(V,E)$ to have vertex vertex set $\{0,1,2,3,4,5,6\}$ and edge set:

\[
\{
(0, 1), (0, 3), (0, 4), (0, 6), (1, 2), (1, 5), (1, 6), (2, 3), (2, 4),
\]
\[
(2, 6), (3, 4), (3, 5), (4, 5), (5, 6)
\}.
\]
It has girth $3$, diameter $2$, and is a
Hamiltonian graph (with circumference $7$).
Its automorphism group is order $48$ and it does not act transitively on the vertices.
The size of the critical group of this graph is $1200=2^43\cdot 5^2$.

There are no non-trivial harmonic morphisms to $C_3$, $C_4$, $P_4$ or to $K_4$.
However, there are exactly $6$ non-trivial harmonic morphisms $\Gamma\to St_4$
and exactly $8$ non-trivial harmonic morphisms $\Gamma\to P_3$.

\end{enumerate}



\subsubsection{Examples}

\begin{example}% d4reg7a
 {\tt d4reg7a}\newline
{\rm  
  Let $\Gamma_2$ denote the graph {\tt d4reg7a} as given above.
There are no harmonic quotient morphisms
$\pi: \Gamma_2\to \Gamma_1=\Gamma_2/G$,
for some $G$.
}
\end{example}
  

\begin{example}% d4reg7b
 {\tt d4reg7b}\newline
{\rm  
  Let $\Gamma_2$ denote the graph {\tt d4reg7b} as given above.
There are $27$ (not necessarily distinct) harmonic quotient morphisms
$\pi: \Gamma_2\to \Gamma_1=\Gamma_2/G$,
for some $G$.

Two of these $27$ are distinct harmonic quotient morphisms
$\pi: \Gamma_2\to \Gamma_1=\Gamma_2/G$,
for some $G$, where $\Gamma_1=K_{2,3}$.
An example for
\[
G = \langle (1,6)(3,4) \rangle
\]
is tabulated below and the uncolored graphs are depicted in
Figure \ref{fig:harmonic-quotient_d7reg7b+K23}.

\begin{center}
\begin{tabular}{c|c|c|c}  
  $x$ & $\phi(x)$ & $m_\phi(x)$ & $\nu_\phi(x)$ \\ \hline
  0   &    3   &        2         &  0   \\ 
  1   &    2   &        1         &  1  \\ 
  2   &    1   &        2         &  0  \\ 
  3   &    0   &        1         &  1  \\ 
  4   &    0   &        1         &  1   \\
  5   &    4   &        2         &  0   \\
  6   &    2   &        1         &  1   \\   
\end{tabular}
\end{center}

\[
  {\rm genus}(\Gamma_2) =8,\ \ \ \ \
  {\rm genus}(\Gamma_1) =2,
\]

\[
deg(\phi) =2.
\]

\begin{figure}[H]
\begin{minipage}{\textwidth}
\begin{center}
%\vspace{8.0 cm}
\includegraphics[height=5cm,width=5cm]{graphics/graph_d4reg7b.png}
\end{center}
\begin{center}
%\vspace{8.0 cm}
\includegraphics[height=4cm,width=4cm]{graphics/K23-quotient-of-d4reg7b.png}
\end{center}
\end{minipage}
\caption{The quartic graph $\Gamma_2$ ({\tt d4reg7b}), with a degree two harmonic quotient morphism to the $K_{2,3}$ graph, $\Gamma_1=\Gamma_2/G$.}
\label{fig:harmonic-quotient_d7reg7b+K23}
\end{figure}

Two of the remaining $25$ are distinct harmonic quotient morphisms
$\pi: \Gamma_2\to \Gamma_1=\Gamma_2/G$,
for some $G$, where $\Gamma_1$ is the graph \verb+DF{+,
a $5$-graph\footnote{This is the graph6 notation. In the ISGCI notation, it is $\overline{K_3\cup 2K_1}$.}.
An example for
\[
G = \langle (1,3)(4,6) \rangle,
\]
which is in fact a homomorphism,
is tabulated below and the uncolored graphs are depicted in
Figure \ref{fig:harmonic-quotient_d4reg7b+fan5}.


\begin{center}
\begin{tabular}{c|c|c|c}  
  $x$ & $\phi(x)$ & $m_\phi(x)$ & $\nu_\phi(x)$ \\ \hline
  0   &    4   &        2         &  0   \\ 
  1   &    3   &        1         &  0  \\ 
  2   &    0   &        2         &  0  \\ 
  3   &    3   &        1         &  0  \\ 
  4   &    1   &        1         &  0   \\
  5   &    2   &        2         &  0   \\
  6   &    1   &        1         &  0   \\   
\end{tabular}
\end{center}

\[
  {\rm genus}(\Gamma_2) =8,\ \ \ \ \
  {\rm genus}(\Gamma_1) =3,
\]

\[
deg(\phi) =2.
\]

\begin{figure}[H]
\begin{minipage}{\textwidth}
\begin{center}
%\vspace{8.0 cm}
\includegraphics[height=5cm,width=5cm]{graphics/graph_d4reg7b.png}
\end{center}
\begin{center}
%\vspace{8.0 cm}
\includegraphics[height=4cm,width=4cm]{graphics/fan5-quotient-of-d4reg7b.png}
\end{center}
\end{minipage}
\cprotect\caption{The quartic graph  $\Gamma_2$ ({\tt d4reg7b}), with a degree two harmonic quotient morphism to the the graph $\Gamma_1=\Gamma_2/G$.}
\label{fig:harmonic-quotient_d4reg7b+fan5}
\end{figure}

Three of the remaining $23$ are (not necessarily distinct)
harmonic quotient morphisms
$\pi: \Gamma_2\to \Gamma_1=\Gamma_2/G$,
for some $G$, where $\Gamma_1=St_4$ is the star graph.
\index{star graph}

Six of the remaining $20$ are (not necessarily distinct)
harmonic quotient morphisms
$\pi: \Gamma_2\to \Gamma_1=\Gamma_2/G$,
for some $G$, where $\Gamma_1=P_3$ is the path graph of length $3$.

The remaining $14$ are (not necessarily distinct)
harmonic quotient morphisms
$\pi: \Gamma_2\to \Gamma_1=\Gamma_2/G$,
for some $G$, where $\Gamma_1=P_2$ is the path graph of length $2$.

}
\end{example}
  

  
\subsection{$8$ vertices}
\label{sec:d4reg8}

There are (up to isomorphism) exactly six $4$-regular connected
graphs on $8$ vertices, $V=\{0,1,2,3,4,5,6,7\}$.
By Ore's Theorem, all of these graphs are Hamiltonian.

\begin{enumerate}
\item
{\tt d4reg8a}
\newline
Define the graph $\Gamma = (V,E)$, where


\[
E = \{
(0, 1), (0, 5), (0, 6), (0, 7), (1, 3), (1, 6), (1, 7), (2, 3), (2, 4),
\]
\[
(2, 5), (2, 7), (3, 4), (3, 6), (4, 5), (4, 6), (5, 7)
\}.
\]
It has diameter $2$,
girth $3$, circumference $8$,
chromatic number $4$, 
and automorphism group of size $16$.
This graph is vertex transitive but not edge transitive.
%aut gp size:  16 
%aut gp gens:  [(1,7)(2,3)(5,6), (0,1)(2,4)(3,5)(6,7)]
The size of the critical group of this graph is $3528=2^3 3^2 7^2$.
\index{vertex transitive graph}

There are no non-trivial harmonic morphisms from $\Gamma$ to
$St_4$, $P_3$, $P_4$, $C_4$, or to $K_4$. However, there are $12$
non-trivial harmonic morphisms from $\Gamma$ to $C_3$.

\item
{\tt d4reg8b}
\newline  
Define the graph $\Gamma = (V,E)$, where

\[
E =
\{
(0, 1), (0, 2), (0, 3), (0, 6), (1, 2), (1, 3), (1, 4), (2, 3),
\]
\[
(2, 7), (3, 5), (4, 5), (4, 6), (4, 7), (5, 6), (5, 7), (6, 7)
\}.
\]
It has diameter $2$, girth $3$, circumference $8$,
chromatic number $4$, and automorphism group of size $48$.
This graph is vertex transitive but not edge transitive.
%aut gp size:  48
%aut gp gens:  [(2,3)(5,7), (1,3)(4,5), (0,1,3)(4,5,6), (0,4)(1,6)(2,5)(3,7)] 
The size of the critical group of this graph is $3456=2^7 3^3$.
\index{vertex transitive graph}

There are $24$ non-trivial harmonic morphisms from $\Gamma$ to $K_4$,
$36$ non-trivial harmonic morphisms from $\Gamma$ to $C_3$,
but no non-trivial harmonic morphisms to $St_4$, $P_3$, $P_4$, or to $C_4$.

\item
{\tt d4reg8c}
\newline  
Define the graph $\Gamma = (V,E)$, where

\[
E = \{
(0, 1), (0, 2), (0, 5), (0, 6), (1, 3), (1, 4), (1, 7), (2, 3),
\]
\[
(2, 4), (2, 7), (3, 5), (3, 6), (4, 5), (4, 6), (5, 7), (6, 7)
\}.
\]
This graph is isomorphic to the complete bipartite graph
$K_{4,4}$.
\index{complete bipartite graph}
It has diameter $2$,
girth $3$, circumference $8$,
chromatic number $2$, 
critical group of size $4096 = 2^{12}$, 
and automorphism group of size $1152$.
It is also a circulant graph with parameters
$(8, [1, 3, 5, 7])$.
\index{circulant graph}
This graph is both vertex transitive and edge transitive,
and is strongly regular with ``trivial'' paramemters $(8,4,0,4)$.
%aut gp size:  1152
%aut gp gens:  [(5,6), (4,7), (3,4), (2,5), (1,2), (0,1)(2,3)(4,5)(6,7)] 
\index{vertex transitive graph}
\index{edge transitive graph}

There are $36$ non-trivial harmonic morphisms from $\Gamma$ to $P_3$,
$72$ non-trivial harmonic morphisms\footnote{Moreover, it appears that
  all of these are multicovers.} from $\Gamma$ to $C_4$,
$576$ non-trivial harmonic
morphisms\footnote{Moreover, it appears that all of these
  are quasicovers.}
from $\Gamma$ to $K_4$, and
$48$ non-trivial harmonic morphisms from $\Gamma$ to $St_4$.
However, it has
no non-trivial harmonic morphisms from $\Gamma$ to $C_3$ or to $P_4$.

\item
{\tt d4reg8d}
\newline  
Define the graph $\Gamma = (V,E)$, where

\[
E = \{
(0, 1), (0, 2), (0, 4), (0, 6), (1, 3), (1, 5), (1, 6), (2, 3),
\]
\[
(2, 5), (2, 7), (3, 4), (3, 7), (4, 5), (4, 6), (5, 7), (6, 7)
\}.
\]
It has diameter $2$,
girth $3$, circumference $8$,
chromatic number $3$, 
critical group of size $3840 = 2^8 3 \cdot 5$, 
and automorphism group of size $16$.
This graph is neither vertex transitive nor edge transitive.
%aut gp size:  16
%aut gp gens:  [(3,5), (1,4), (0,2)(1,3)(4,5)(6,7), (0,6)(2,7)] 

There are no non-trivial harmonic morphisms from $\Gamma$ to $St_4$,
$P_4$ or to $C_4$. However,
there are $4$ non-trivial harmonic morphisms to $P_3$,
$6$ non-trivial harmonic morphisms to $C_3$,
and $48$ non-trivial harmonic morphisms\footnote{Moreover,
  it appears that all of these are quasicovers.}
to $K_4$.

\item
{\tt d4reg8e}
\newline  
Define the graph $\Gamma = (V,E)$, where

\[
E = \{
(0, 1), (0, 2), (0, 6), (0, 7), (1, 3), (1, 4), (1, 5), (2, 3), 
\]
\[
(2, 4), (2, 7), (3, 5), (3, 7), (4, 5), (4, 6), (5, 6), (6, 7)
\}.
\]
It has diameter $2$,
girth $3$, circumference $8$,
chromatic number $3$, 
critical group of size $3696 = 2^4 3\cdot 7\cdot 11$, 
and automorphism group of size $4$.
%aut gp size:  4 
%aut gp gens:  [(0,1)(2,4)(3,6)(5,7), (0,2)(1,4)(3,6)]

There are no non-trivial harmonic morphisms from $\Gamma$ to $St_4$,
$P_3$, $P_4$, $C_4$ or to $K_4$, but there are $6$ non-trivial harmonic
morphisms to $C_3$.

\item
{\tt d4reg8f}
\newline  
Define the graph $\Gamma = (V,E)$, where

\[
E = \{
(0, 1), (0, 2), (0, 6), (0, 7), (1, 2), (1, 3), (1, 4), (2, 3),
\]
\[
(2, 4), (3, 5), (3, 7), (4, 5), (4, 6), (5, 6), (5, 7), (6, 7)
\}.
\]
It has diameter $2$,
girth $3$, circumference $8$,
chromatic number $4$, 
critical group of size $3610 = 2\cdot 5\cdot 19^2$, 
and automorphism group of size $12$.
%aut gp size:  4 
%aut gp gens:  [(0,1)(2,4)(3,6)(5,7), (0,2)(1,4)(3,6)]

There are no non-trivial harmonic morphisms from $\Gamma$ to $St_4$, $P_4$, $C_3$, $C_4$ or to $K_4$. However, there are $2$ non-trivial harmonic
morphisms to $P_3$.

\end{enumerate}



\subsubsection{Examples}

\begin{example}% d4reg8a
 {\tt d4reg8a}\newline
{\rm  
  Let $\Gamma_2$ denote the graph {\tt d4reg8a} as given above.
There are $10$ (not necessarily distinct) harmonic quotient morphisms
$\pi: \Gamma_2\to \Gamma_1=\Gamma_2/G$,
for some $G$.

Four of these $10$ are distinct harmonic quotient morphisms
$\pi: \Gamma_2\to \Gamma_1=\Gamma_2/G$,
for some $G$, where $\Gamma_1$ is the \verb+Dh{+ graph\footnote{This is the graph6 notation. In the ISGCI notation, it is called a gem graph or a $3$-fan.}.
\index{gem graph}
An example for
\[
G = \langle  (1,7)(2,3)(5,6)\rangle,
\]
is tabulated below and the uncolored graphs are depicted in
Figure \ref{fig:harmonic-quotient_d4reg8a+gem}.

\begin{center}
\begin{tabular}{c|c|c|c}  
  $x$ & $\phi(x)$ & $m_\phi(x)$ & $\nu_\phi(x)$ \\ \hline
  0   &    4   &        2         &  0   \\ 
  1   &    1   &        1         &  1  \\ 
  2   &    2   &        1         &  1  \\ 
  3   &    2   &        1         &  1  \\ 
  4   &    3   &        2         &  0   \\
  5   &    0   &        1         &  0   \\
  6   &    0   &        1         &  0   \\
  7   &    1   &        1         &  1   \\     
\end{tabular}
\end{center}

\[
  {\rm genus}(\Gamma_2) =9,\ \ \ \ \
  {\rm genus}(\Gamma_1) =3,
\]

\[
deg(\phi) =2.
\]

\begin{figure}[H]
\begin{minipage}{\textwidth}
\begin{center}
%\vspace{8.0 cm}
\includegraphics[height=5cm,width=5cm]{graphics/graph_d4reg8a.png}
\end{center}
\begin{center}
%\vspace{8.0 cm}
\includegraphics[height=4cm,width=4cm]{graphics/gem-quotient-of-d4reg8a.png}
\end{center}
\end{minipage}
\cprotect\caption{The quartic graph  $\Gamma_2$ ({\tt d4reg8a}), with a degree two harmonic quotient morphism to the graph $\Gamma_1=\Gamma_2/G$.}
\label{fig:harmonic-quotient_d4reg8a+gem}
\end{figure}

Two of the remaining $6$ are distinct harmonic quotient morphisms
$\pi: \Gamma_2\to \Gamma_1=\Gamma_2/G$,
for some $G$, where $\Gamma_1=C_3$. The rest are (not necessarily distinct) 
harmonic quotient morphisms
$\pi: \Gamma_2\to \Gamma_1=\Gamma_2/G$,
for some $G$, where $\Gamma_1=P_2$.

}
\end{example}




\begin{example}% d4reg8b
 {\tt d4reg8b}\newline
{\rm  
  Let $\Gamma_2$ denote the graph {\tt d4reg8b} as given above.
There are $46$ (not necessarily distinct) harmonic quotient morphisms
$\pi: \Gamma_2\to \Gamma_1=\Gamma_2/G$,
for some $G$.

Exactly one of these $46$ is a distinct harmonic quotient morphism
$\pi: \Gamma_2\to \Gamma_1=\Gamma_2/G$,
for some $G$, where $\Gamma_1=K_4$. In fact, it is a quasicover.
This example, for

\[
G = \langle (0,6)(1,4)(2,7)(3,5)  \rangle,
\]
is tabulated below and depicted in Figure \ref{fig:harmonic-quotient_d4reg8b+K4}.

\begin{center}
\begin{tabular}{c|c|c|c}  
  $x$ & $\phi(x)$ & $m_\phi(x)$ & $\nu_\phi(x)$ \\ \hline
  0   &    2   &        1         &  1   \\ 
  1   &    0   &        1         &  1  \\ 
  2   &    1   &        1         &  1  \\ 
  3   &    3   &        1         &  1  \\ 
  4   &    0   &        1         &  1   \\
  5   &    3   &        1         &  1   \\
  6   &    2   &        1         &  1   \\
  7   &    1   &        1         &  1   \\     
\end{tabular}
\end{center}

\[
  {\rm genus}(\Gamma_2) =9,\ \ \ \ \
  {\rm genus}(\Gamma_1) =3,
\]

\[
deg(\phi) =2.
\]
\begin{figure}[H]
\begin{minipage}{\textwidth}
\begin{center}
%\vspace{8.0 cm}
\includegraphics[height=5cm,width=5cm]{graphics/harmonic-quotient_d4reg8b-K4_20130321.png}
\end{center}
\begin{center}
%\vspace{8.0 cm}
\includegraphics[height=4cm,width=4cm]{graphics/K4-quotient-of-d4reg8b_0123.png}
\end{center}
\end{minipage}
\caption{The quartic graph  $\Gamma_2$, with a degree two harmonic quotient morphism to the complete graph, $\Gamma_1=K_4$.}
\label{fig:harmonic-quotient_d4reg8b+K4}
\end{figure}


Six of the remaining $45$ are distinct harmonic quotient morphisms
$\pi: \Gamma_2\to \Gamma_1=\Gamma_2/G$,
for some $G$, where $\Gamma_1=C_3$.

The remaining $39$ are (not necessarily distinct)
harmonic quotient morphisms
$\pi: \Gamma_2\to \Gamma_1=\Gamma_2/G$,
for some $G$, where $\Gamma_1=P_2$.

}
\end{example}




\begin{example}% d4reg8c
 {\tt d4reg8c}\newline
{\rm  
  Let $\Gamma_2$ denote the graph $K_{4,4}$ as given above.
There are $2221$ (not necessarily distinct) harmonic quotient morphisms
$\pi: \Gamma_2\to \Gamma_1=\Gamma_2/G$,
for some $G$.

Of these $2221$, $24$ are (not necessarily distinct) harmonic quotient morphisms
$\pi: \Gamma_2\to \Gamma_2/G = K_4$, for some $G$. In fact, each of the $24$ is a quasicover.
Of the remaining $2191$, $234$ are (not necessarily distinct) harmonic quotient morphisms
$\pi: \Gamma_2\to \Gamma_2/G = C_4$, for some $G$.
Of the remaining $1963$, $636$ are (not necessarily distinct) harmonic quotient morphisms
$\pi: \Gamma_2\to \Gamma_2/G = P_3$, for some $G$.
}
\end{example}




\begin{example}% d4reg8d
 {\tt d4reg8d}\newline
{\rm  
  Let $\Gamma_2$ denote the graph {\tt d4reg8d} as given above.
There are $11$ (not necessarily distinct) harmonic quotient morphisms
$\pi: \Gamma_2\to \Gamma_1=\Gamma_2/G$,
for some $G$.

Of these $11$, $2$ are distinct harmonic quotient morphisms
$\pi: \Gamma_2\to \Gamma_2/G = K_4$, for some $G$. In fact, they are each a quasicover.
Of the remaining $9$, $3$ are (not necessarily distinct) harmonic quotient morphisms
$\pi: \Gamma_2\to \Gamma_2/G = C_3$, for some $G$.
Of the remaining $6$, $636$ are (not necessarily distinct) harmonic quotient morphisms
$\pi: \Gamma_2\to \Gamma_2/G = P_2$, for some $G$.
}
\end{example}





\begin{example}% d4reg8e
 {\tt d4reg8e}\newline
{\rm  
  Let $\Gamma_2$ denote the graph {\tt d4reg8e} as given above.
There is exactly one harmonic quotient morphism
$\pi: \Gamma_2\to \Gamma_1=\Gamma_2/G$,
for some subgroup $G$ of $Aut(\Gamma_2)$.
Tabulated below, this is a quotient morphism
$\phi: \Gamma_2\to \Gamma_1=\Gamma_2/G$,
where $\Gamma_1$ is the gem graph and

\[
G = \langle (0,2)(1,4)(3,6) \rangle,
\]
is tabulated below.

\begin{center}
\begin{tabular}{c|c|c|c}  
  $x$ & $\phi(x)$ & $m_\phi(x)$ & $\nu_\phi(x)$ \\ \hline
  0   &    2   &        1         &  1   \\ 
  1   &    0   &        1         &  1  \\ 
  2   &    2   &        1         &  1  \\ 
  3   &    3   &        1         &  0  \\ 
  4   &    0   &        1         &  1   \\
  5   &    1   &        2         &  0   \\
  6   &    3   &        1         &  0   \\
  7   &    4   &        2         &  0   \\     
\end{tabular}
\end{center}

\[
  {\rm genus}(\Gamma_2) =9,\ \ \ \ \
  {\rm genus}(\Gamma_1) =3,
\]

\[
deg(\phi) =2.
\]

}
\end{example}



\begin{example}% d4reg8f
 {\tt d4reg8f}\newline
{\rm  
  Let $\Gamma_2$ denote the graph {\tt d4reg8f} as given above.
There are two (identical) harmonic quotient morphisms
$\pi: \Gamma_2\to \Gamma_1=\Gamma_2/G$,
for some subgroup $G$ of $Aut(\Gamma_2)$.
Each of these is a quotient morphism
$\phi: \Gamma_2\to \Gamma_1=\Gamma_2/G$,
where $\Gamma_1=P_3$ (in each case, the vertices of $P_3$
are labeled in such a way that $0$ has degree $2$).
One such subgroup is

\[
G_1 = \langle (1,2), (0,4,3)(5,7,6) \rangle,
\]
and the other is

\[
G_2 = \langle (1,2)(3,4)(6,7), (0,4,3)(5,7,6) \rangle.
\]
Each has cardinality $6$. We tabulate
$\phi_1: \Gamma_2\to \Gamma_1=\Gamma_2/G_1=P_3$
below.

\begin{center}
\begin{tabular}{c|c|c|c}  
  $x$ & $\phi(x)$ & $m_\phi(x)$ & $\nu_\phi(x)$ \\ \hline
  0   &    0   &        2         &  0   \\ 
  1   &    2   &        3         &  1  \\ 
  2   &    2   &        3         &  1  \\ 
  3   &    0   &        2         &  0  \\ 
  4   &    0   &        2         &  0   \\
  5   &    1   &        2         &  2   \\
  6   &    1   &        2         &  2   \\
  7   &    1   &        2         &  2   \\     
\end{tabular}
\end{center}

\[
  {\rm genus}(\Gamma_2) =9,\ \ \ \ \
  {\rm genus}(\Gamma_1) =0,
\]

\[
deg(\phi) =6.
\]

}
\end{example}


\subsection{$9$ vertices}
\label{sec:d4reg9}

There are (up to isomorphism) exactly $16$ $4$-regular connected graphs on $9$ vertices.
%For the majority of these, it's currently unknown how many (if any)
%non-trivial harmonic morphisms map from these graphs to $C_4$ or $K_4$.
Below, we include $8$ of these $16$ as examples.

\begin{enumerate}
\item
  {\tt d4reg9a}\newline
Define the graph  $\Gamma$ to have edge set:

\[
E = \{
(0, 1), (0, 2), (0, 7), (0, 8), (1, 2), (1, 3), (1, 7), (2, 3), (2, 8),
\]
\[
 (3, 4), (3, 5), (4, 5), (4, 6), (4, 8), (5, 6), (5, 7), (6, 7), (6, 8)
\}.
\]
It has diameter $2$,
girth $3$, circumference $9$ (i.e., it's Hamiltonian),
critical group of size $10800 = 2^43^35^2$, 
and automorphism group of size $12$.

There are no non-trivial harmonic morphisms from $\Gamma$ to $St_4$,
$P_4$, $C_4$, or to $K_4$.  However, there are $2$ non-trivial
harmonic morphisms from $\Gamma$ to $P_3$ and $24$ non-trivial
harmonic morphisms from $\Gamma$ to $C_3$.

\item
  {\tt d4reg9b}\newline  
Define the graph $\Gamma$ to have edge set:

\[
E = \{
(0, 1), (0, 3), (0, 7), (0, 8), (1, 2), (1, 3), (1, 7), (2, 3), (2, 5), (2, 8),
\]
\[
(3, 4), (4, 5), (4, 6), (4, 8), (5, 6), (5, 7), (6, 7), (6, 8)
\}.
\]
It has diameter $2$,
girth $3$, circumference $9$ (i.e., it's Hamiltonian),
critical group of size $11360 = 2^5 5\cdot 71$, 
and automorphism group of size $2$.

There are no non-trivial harmonic morphisms from $\Gamma$ to
$St_4$, $P_3$, $P_4$, $C_3$, $C_4$, or to $K_4$.  

\item
{\tt d4reg9c}\newline  
Define the graph $\Gamma$ to have edge set:

\[
E = \{
(0, 1), (0, 2), (0, 7), (0, 8), (1, 2), (1, 3), (1, 4), (2, 3), (2, 8), 
\]
\[
(3, 4), (3, 5), (4, 5), (4, 6), (5, 6), (5, 7), (6, 7), (6, 8), (7, 8)
\}.
\]
It has diameter $2$,
girth $3$, circumference $9$ (i.e., it's Hamiltonian),
critical group of size $10404 = 2^2 3^2 17^2$, 
and automorphism group of size $18$.
This is a vertex transitive but not edge transitive graph.
\index{vertex transitive graph}

There are no non-trivial harmonic morphisms from $\Gamma$ to $St_4$,
$P_3$, $P_4$, $C_4$, or to $K_4$.  However, there
are $6$ non-trivial harmonic morphisms from $\Gamma$ to $C_3$. 

\item
{\tt d4reg9d}\newline
Define the graph $\Gamma$ to have edge set:

\[
E = \{
(0, 1), (0, 5), (0, 7), (0, 8), (1, 2), (1, 4), (1, 7), (2, 3), (2, 4), (2, 5),
\]
\[
(3, 4), (3, 6), (3, 8), (4, 5), (5, 6), (6, 7), (6, 8), (7, 8)
\}.
\]
It has diameter $2$,
girth $3$, circumference $9$ (i.e., it's Hamiltonian),
critical group of size $11005 = 5\cdot 31\cdot 71$, 
and automorphism group of size $4$.

There are no non-trivial harmonic morphisms from $\Gamma$ to
$P_4$, $C_3$, $C_4$, or to $K_4$. However, there are $2$ non-trivial harmonic
morphisms from $\Gamma$ to $P_3$. 


\item
{\tt d4reg9e}\newline  
Define the graph $\Gamma$ to have edge set:

\[
E = \{
(0, 1), (0, 3), (0, 5), (0, 8), (1, 2), (1, 4), (1, 7), (2, 3), (2, 5), (2, 7),
\]
\[
 (3, 4), (3, 8), (4, 5), (4, 6), (5, 6), (6, 7), (6, 8), (7, 8)
\}.
\]
It has diameter $2$,
girth $3$, circumference $9$ (i.e., it's Hamiltonian),
critical group of size $12096 = 2^6 3^3 7$, 
and automorphism group of size $12$.

There are no non-trivial harmonic morphisms from $\Gamma$ to $St_4$,
$P_4$, $C_4$, or to $K_4$. However,
there are $6$ non-trivial harmonic
morphisms\footnote{Moreover,
  it appears that all of these are quasicovers
  (with constant horizontal multiplicity $1$ and constant vertical
  multiplicity $2$).}
from $\Gamma$ to $C_3$,
and $6$ non-trivial harmonic morphisms from $\Gamma$ to $P_3$.


\item
{\tt d4reg9f}\newline    
  Define the graph $\Gamma$
  %d4reg9l=d4reg9-12
  to have edge set:

\[
E = \{
(0, 1), (0, 3), (0, 6), (0, 8), (1, 2), (1, 4), (1, 6), (2, 3), (2, 5), (2, 7), (3, 4),
\]
\[
(3, 7), (4, 5), (4, 8), (5, 6), (5, 8), (6, 7), (7, 8)
\}.
\]
It has diameter $2$,
girth $3$, circumference $9$ (i.e., it's Hamiltonian),
critical group of size $12321 = 3^2 37^2$, 
and automorphism group of size $18$.
This is a vertex transitive but not edge transitive graph.
\index{vertex transitive graph}

There are no non-trivial harmonic morphisms from $\Gamma$ to $St_4$,
$P_3$, $P_4$, $C_4$, or to $K_4$. However,
there are $6$ non-trivial harmonic morphisms from $\Gamma$ to $C_3$.

\item
{\tt d4reg9g}\newline    
  Define the graph $\Gamma$
  %d4reg9n=d4reg9-14
  to have edge set:

\[
E = \{
(0, 1), (0, 3), (0, 4), (0, 8), (1, 2), (1, 5), (1, 8), (2, 3), (2, 5), (2, 7), (3, 4),
\]
\[
 (3, 7), (4, 5), (4, 6), (5, 6), (6, 7), (6, 8), (7, 8)
\}.
\]
It has diameter $2$,
girth $3$, circumference $9$ (i.e., it's Hamiltonian),
critical group of size $11664 = 2^4 3^6$, 
and automorphism group of size $72$.
This is a vertex transitive, edge transitive, arc transitive,
and distance regular graph.
Morover, it's $(9,4,1,2)$-strongly regular, isomorphic to the Paley graph
on $9$ vertices.
\index{vertex transitive graph}
\index{edge transitive graph}
\index{arc transitive graph}
\index{distance regular graph}
\index{strongly regular graph}
\index{Paley graph}

There are no non-trivial harmonic morphisms from $\Gamma$ to $St_4$,
$P_4$, $C_4$. or to $K_4$. However, there are $18$ non-trivial
harmonic morphisms from $\Gamma$ to $P_3$
and $78$ non-trivial harmonic
morphisms\footnote{Moreover,
  some, not not all, of these are quasicovers
  (with constant horizontal multiplicity $1$ and constant vertical
  multiplicity $2$).} from $\Gamma$ to $C_3$.
\label{page:paley9}

\item
{\tt d4reg9h}\newline    
\label{page:d4reg9-15}  
  Define the graph $\Gamma$ 
% d4reg9-15 ## integral spectrum [4, 2, 1, 0, 0, -1, -1, -2, -3]
  % ## denoted I_{9,3} in [CSS98]
to have edge set:

\[
E = \{
(0, 1), (0, 4), (0, 6), (0, 8), (1, 2), (1, 3), (1, 5), (2, 3),
\]
\[
(2, 4), (2, 7), (3, 4), (3, 7), (4, 5), (5, 6), (5, 8), (6, 7), (6, 8), (7, 8)
\}.
\]
It is an integral graph with diameter $2$,
girth $3$, circumference $9$ (i.e., it's Hamiltonian),
critical group of size $11200 = 2^6 5^27$, 
and automorphism group of size $32$.
It has no non-trivial harmonic morphisms to 
$C_3$, $C_4$, $P_4$, $K_4$, but $8$ non-trivial harmonic morphisms to
$P_3$.

See also page \pageref{page:d4reg9-15} below.
  
\item
{\tt d4reg9i}\newline  
  Define the graph $\Gamma$
  % d4reg9p
  to have edge set:

\[
E = \{
(0, 1), (0, 3), (0, 7), (0, 8), (1, 2), (1, 4), (1, 5), (2, 3), (2, 4), (2, 5),
\]
\[
(3, 7), (3, 8), (4, 5), (4, 6), (5, 6), (6, 7), (6, 8), (7, 8)
\}.
\]
It has diameter $2$,
girth $3$, circumference $9$ (i.e., it's Hamiltonian),
critical group of size $9600 = 2^7 3 \cdot 5^2$, 
and automorphism group of size $16$.

There are no non-trivial harmonic morphisms from $\Gamma$ to $St_4$,
$P_3$, $P_4$, $C_4$, or to $K_4$. However, there are
exactly $6$ non-trivial harmonic morphisms from $\Gamma$ to $C_3$.


\end{enumerate}

\subsubsection{Examples}


\begin{example}% d4reg9a
 {\tt d4reg9a}\newline
{\rm  
  Let $\Gamma_2$ denote the graph {\tt d4reg9a} as given above.
There are $11$ harmonic quotient morphisms
$\pi: \Gamma_2\to \Gamma_1=\Gamma_2/G$,
for some subgroup $G$ of $Aut(\Gamma_2)$.

Three of these harmonic quotient morphisms are to the co-antenna graph.
For example, if 
\index{coantenna graph}

\[
G = \langle (1,2)(4,5)(7,8) \rangle.
\]
We tabulate $\phi_1: \Gamma_2\to \Gamma_1=\Gamma_2/G$
below and give the uncolored graphs in Figure \ref{fig:harmonic-quotient_d4reg9a+coantenna}.

\begin{center}
\begin{tabular}{c|c|c|c}  
  $x$ & $\phi(x)$ & $m_\phi(x)$ & $\nu_\phi(x)$ \\ \hline
  0   &    3   &        2         &  0   \\ 
  1   &    2   &        1         &  1  \\ 
  2   &    1   &        1         &  1  \\ 
  3   &    0   &        2         &  0  \\ 
  4   &    1   &        1         &  1   \\
  5   &    2   &        1         &  1   \\
  6   &    3   &        2         &  0   \\
  7   &    5   &        1         &  0   \\
  8   &    4   &        1         &  0   \\ 
\end{tabular}
\end{center}

\[
  {\rm genus}(\Gamma_2) =10,\ \ \ \ \
  {\rm genus}(\Gamma_1) =3,
\]

\[
deg(\phi) =2.
\]

\begin{figure}[H]
\begin{minipage}{\textwidth}
\begin{center}
%\vspace{8.0 cm}
\includegraphics[height=5cm,width=5cm]{graphics/graph_d4reg9a.png}
\end{center}
\begin{center}
%\vspace{8.0 cm}
\includegraphics[height=4cm,width=4cm]{graphics/coantenna-quotient-of-d4reg9a.png}
\end{center}
\end{minipage}
\caption{The quartic graph  $\Gamma_2$ ({\tt d4reg9a}), with a degree two harmonic quotient morphism to the coantenna graph, $\Gamma_1$.}
\label{fig:harmonic-quotient_d4reg9a+coantenna}
\end{figure}


One of the remaining $8$ is a harmonic quotient homomorphism
$\phi: \Gamma_2\to \Gamma_1=\Gamma_2/G$,
where $\Gamma_1=S_3$, where
\index{$S_3$ graph}

\[
G = \langle (0,6)(1,5)(2,4) \rangle.
\]
We tabulate $\phi_1: \Gamma_2\to \Gamma_1=\Gamma_2/G=S_3$
below and give the uncolored graphs in Figure \ref{fig:harmonic-quotient_d4reg9a+S3}.

\begin{center}
\begin{tabular}{c|c|c|c}  
  $x$ & $\phi(x)$ & $m_\phi(x)$ & $\nu_\phi(x)$ \\ \hline
  0   &    3   &        1         &  0   \\ 
  1   &    2   &        1         &  0  \\ 
  2   &    1   &        1         &  0  \\ 
  3   &    0   &        2         &  0  \\ 
  4   &    1   &        1         &  0   \\
  5   &    2   &        1         &  0   \\
  6   &    3   &        2         &  0   \\
  7   &    5   &        2         &  0   \\
  8   &    4   &        2         &  0   \\ 
\end{tabular}
\end{center}

\[
  {\rm genus}(\Gamma_2) =10,\ \ \ \ \
  {\rm genus}(\Gamma_1) =6,
\]

\[
deg(\phi) =2.
\]

\begin{figure}[H]
\begin{minipage}{\textwidth}
\begin{center}
%\vspace{8.0 cm}
\includegraphics[height=5cm,width=5cm]{graphics/graph_d4reg9a.png}
\end{center}
\begin{center}
%\vspace{8.0 cm}
\includegraphics[height=4cm,width=4cm]{graphics/S3-quotient-of-d4reg9a.png}
\end{center}
\end{minipage}
\caption{The quartic graph  $\Gamma_2$ ({\tt d4reg9a}), with a degree two harmonic quotient morphism to the matchstick graph, $\Gamma_1=S_3$.}
\label{fig:harmonic-quotient_d4reg9a+S3}
\end{figure}

All the $7$ remaining harmonic quotients map to $P_2$ or $P_3$ or the
paw graph (which is $C_3$ plus a leaf ``tail').
\index{paw graph}
One of these, a harmonic quotient morphisms
$\phi: \Gamma_2\to \Gamma_1=\Gamma_2/G$,
where $\Gamma_1$ is the paw graph and where

\[
G = \langle (1,2)(4,5)(7,8), (0,6)(1,4)(2,5)(7,8) \rangle,
\]
is tabulated below.

\begin{center}
\begin{tabular}{c|c|c|c}  
  $x$ & $\phi(x)$ & $m_\phi(x)$ & $\nu_\phi(x)$ \\ \hline
  0   &    1   &        2         &  0   \\ 
  1   &    2   &        1         &  0  \\ 
  2   &    2   &        1         &  0  \\ 
  3   &    0   &        4         &  0  \\ 
  4   &    2   &        1         &  0   \\
  5   &    2   &        1         &  0   \\
  6   &    1   &        2         &  0   \\
  7   &    3   &        2         &  0   \\
  8   &    3   &        2         &  0   \\ 
\end{tabular}
\end{center}

\[
  {\rm genus}(\Gamma_2) =10,\ \ \ \ \
  {\rm genus}(\Gamma_1) =1,
\]

\[
deg(\phi) =4.
\]


}
\end{example}


\begin{example}% d4reg9b
 {\tt d4reg9b}\newline
{\rm  
  Let $\Gamma_2$ denote the graph {\tt d4reg9b} as given above.
There are no harmonic quotient morphisms
$\pi: \Gamma_2\to \Gamma_1=\Gamma_2/G$,
for some subgroup $G$ of $Aut(\Gamma_2)$.

}
\end{example}


\begin{example}% d4reg9c
 {\tt d4reg9c}\newline
{\rm  
Let $\Gamma_2$ denote the graph {\tt d4reg9c} as given above.
There are $4$ harmonic quotient morphisms
$\pi: \Gamma_2\to \Gamma_1=\Gamma_2/G$,
for some subgroup $G$ of $Aut(\Gamma_2)$.

There is only one harmonic quotient morphisms
$\pi: \Gamma_2\to \Gamma_1=\Gamma_2/G$,
for some subgroup $G$ of $Aut(\Gamma_2)$, to
$\Gamma_1=C_3$. In fact, this is a multicover
of degree $3$ and multiplicity $2$.
In this case,

\[
G = \langle (0,3,6)(1,5,8)(2,4,7) \rangle.
\]
We tabulate $\phi_1: \Gamma_2\to \Gamma_1=\Gamma_2/G=C_3$
below and depict it in Figure \ref{fig:harmonic-quotient_d4reg9c-C3_021012012}.
%harmonic-quotient_d4reg9c-C3_021012012

\begin{center}
\begin{tabular}{c|c|c|c}  
  $x$ & $\phi(x)$ & $m_\phi(x)$ & $\nu_\phi(x)$ \\ \hline
  0   &    0   &        2         &  0   \\ 
  1   &    2   &        2         &  0  \\ 
  2   &    1   &        2         &  0  \\ 
  3   &    0   &        2         &  0  \\ 
  4   &    1   &        2         &  0   \\
  5   &    2   &        2         &  0   \\
  6   &    0   &        2         &  0   \\
  7   &    1   &        2         &  0   \\
  8   &    2   &        2         &  0   \\ 
\end{tabular}
\end{center}

\[
  {\rm genus}(\Gamma_2) =10,\ \ \ \ \
  {\rm genus}(\Gamma_1) =1,
\]

\[
deg(\phi) =3.
\]

\begin{figure}[H]
\begin{minipage}{\textwidth}
\begin{center}
%\vspace{8.0 cm}
\includegraphics[height=5cm,width=5cm]{graphics/harmonic-quotient_d4reg9c-C3_021012012.png}
\end{center}
\begin{center}
%\vspace{8.0 cm}
\includegraphics[height=4cm,width=4cm]{graphics/cycle3-012.png}
\end{center}
\end{minipage}
\caption{The quartic graph  $\Gamma_2$ ({\tt d4reg9c}), with a degree three harmonic quotient homomorphism (in fact, a multicover) to the cycke graph, $\Gamma_1=C_3$.}
\label{fig:harmonic-quotient_d4reg9c-C3_021012012}
\end{figure}


}
\end{example}



\begin{example}% d4reg9d
 {\tt d4reg9d}\newline
{\rm  
  Let $\Gamma_2$ denote the graph {\tt d4reg9d} as given above.
There are no harmonic quotient morphisms
$\pi: \Gamma_2\to \Gamma_1=\Gamma_2/G$,
for some subgroup $G$ of $Aut(\Gamma_2)$.

}
\end{example}




\begin{example}% d4reg9e
 {\tt d4reg9e}\newline
{\rm  
Let $\Gamma_2$ denote the graph {\tt d4reg9e} as given above.
There are $5$ harmonic quotient morphisms
$\pi: \Gamma_2\to \Gamma_1=\Gamma_2/G$,
for some subgroup $G$ of $Aut(\Gamma_2)$.
Three of these $5$ are quotients to the \verb+D]w+
graph\footnote{This is the graph6 notation. However, the ISGCI notation is
$\overline{P_2\cup P_3}$.}. The remaining two are
(identical morphisms but with different subgroups)
harmonic quotients to $P_2$.

One of the three distinct harmonic quotient morphisms
$\pi: \Gamma_2\to \Gamma_1=\Gamma_2/G=\Gamma_1=C_4+2C_3$,
for the subgroup

\[
G = \langle (0,1)(2,3)(4,5)(7,8) \rangle.
\]
of $Aut(\Gamma_2)$, is tabulated below.
In addition, the uncolored graphs are depicted
in Figure \ref{fig:harmonic-quotient_d4reg9e-C4+2C3}.

\begin{center}
\begin{tabular}{c|c|c|c}  
  $x$ & $\phi(x)$ & $m_\phi(x)$ & $\nu_\phi(x)$ \\ \hline
  0   &    3   &        1         &  1   \\ 
  1   &    3   &        1         &  1  \\ 
  2   &    1   &        1         &  1  \\ 
  3   &    1   &        1         &  1  \\ 
  4   &    2   &        1         &  1   \\
  5   &    2   &        1         &  1   \\
  6   &    0   &        2         &  0   \\
  7   &    4   &        1         &  1   \\
  8   &    4   &        1         &  1   \\ 
\end{tabular}
\end{center}

\[
  {\rm genus}(\Gamma_2) =10,\ \ \ \ \
  {\rm genus}(\Gamma_1) =3,
\]

\[
deg(\phi) =2.
\]

\begin{figure}[H]
\begin{minipage}{\textwidth}
\begin{center}
%\vspace{8.0 cm}
\includegraphics[height=5cm,width=5cm]{graphics/graph_d4reg9e.png}
\end{center}
\begin{center}
%\vspace{8.0 cm}
\includegraphics[height=4cm,width=4cm]{graphics/C4+2C3-quotient-of-d4reg9e.png}
\end{center}
\end{minipage}
\cprotect\caption{The quartic graph  $\Gamma_2$ ({\tt d4reg9e}), with a degree two harmonic quotient morphism to the graph \verb+D]w=+$\Gamma_1=\overline{P_2\cup P_3}$.}
\label{fig:harmonic-quotient_d4reg9e-C4+2C3}
\end{figure}

}
\end{example}



\begin{example}% d4reg9f
 {\tt d4reg9f}\newline
{\rm  
  Let $\Gamma_2$ denote the graph {\tt d4reg9f} as given above.
There are $4$ harmonic quotient morphisms
$\pi: \Gamma_2\to \Gamma_1=\Gamma_2/G$,
for some subgroup $G$ of $Aut(\Gamma_2)$.

Three of these harmonic quotient morphisms are to the path graph,
$\Gamma_1=P_2$.
One of these harmonic quotient morphisms is to the cycle graph,
$\Gamma_1=C_3$.
If 

\[
G = \langle (0,6,1)(2,3,7)(4,8,5) \rangle,
\]
then the quotient $\phi: \Gamma_2\to \Gamma_1=\Gamma_2/G$,
tabulated below, is in fact a quasicover.

\begin{center}
\begin{tabular}{c|c|c|c}  
  $x$ & $\phi(x)$ & $m_\phi(x)$ & $\nu_\phi(x)$ \\ \hline
  0   &    0   &        1         &  2   \\ 
  1   &    0   &        1         &  2  \\ 
  2   &    1   &        1         &  2  \\ 
  3   &    1   &        1         &  2  \\ 
  4   &    2   &        1         &  2   \\
  5   &    2   &        1         &  2   \\
  6   &    0   &        1         &  2   \\
  7   &    1   &        1         &  2   \\
  8   &    2   &        1         &  2   \\ 
\end{tabular}
\end{center}

\[
  {\rm genus}(\Gamma_2) =10,\ \ \ \ \
  {\rm genus}(\Gamma_1) =1,
\]
and $\deg(\phi)=3$.

}
\end{example}



\begin{example}% d4reg9g
 {\tt d4reg9g}\newline
{\rm  
  Let $\Gamma_2$ denote the graph {\tt d4reg9g} as given above.
There are $59$ harmonic quotient morphisms
$\pi: \Gamma_2\to \Gamma_1=\Gamma_2/G$,
for some subgroup $G$ of $Aut(\Gamma_2)$.

Six of these harmonic quotient morphisms are to the matchstick graph,
$\Gamma_1=S_3$.
\index{$S_3$ graph}
For example, if 

\[
G = \langle (1,3)(4,8)(5,7) \rangle,
\]
then the quotient $\phi: \Gamma_2\to \Gamma_1=\Gamma_2/G$,
tabulated below, is in fact a homomorphism to $\Gamma_1=S_3$.

\begin{center}
\begin{tabular}{c|c|c|c}  
  $x$ & $\phi(x)$ & $m_\phi(x)$ & $\nu_\phi(x)$ \\ \hline
  0   &    5   &        2         &  0   \\ 
  1   &    3   &        1         &  0  \\ 
  2   &    0   &        2         &  0  \\ 
  3   &    3   &        1         &  0  \\ 
  4   &    4   &        1         &  0   \\
  5   &    1   &        1         &  0   \\
  6   &    2   &        2         &  0   \\
  7   &    1   &        1         &  0   \\
  8   &    4   &        1         &  0   \\ 
\end{tabular}
\end{center}

\[
  {\rm genus}(\Gamma_2) =10,\ \ \ \ \
  {\rm genus}(\Gamma_1) =6,
\]
and $\deg(\phi)=2$.

Of the remaining $56$, $9$ are harmonic quotient morphisms to
the paw graph (i.e., $C_3$ plus a leaf edge).
\index{paw graph}
Of the remaining $47$, $8$ are (not necessarily distinct) harmonic
quotient morphisms to the graph $C_3$.
All of the remaining $39$ are (not necessarily distinct) harmonic
quotient morphisms to the graph $P_2$.

}
\end{example}



\begin{example}% d4reg9h
 {\tt d4reg9h}\newline
{\rm  
Let $\Gamma_2$ denote the graph {\tt d4reg9h} as given above.
There are $4$ harmonic quotient morphisms
$\pi: \Gamma_2\to \Gamma_1=\Gamma_2/G$,
for some subgroup $G$ of $Aut(\Gamma_2)$.
All of these are quotients to a graph we shall call
the {\it complete house graph}, however the ISGCI notation is
$\overline{P_3\cup 2K_1}$. 
\index{house graph!complete}
\index{complete house graph}
One of the these harmonic quotient morphisms,
$\pi: \Gamma_2\to \Gamma_1=\Gamma_2/G=\Gamma_1$,
for the subgroup

\[
G = \langle (0,1)(2,6)(3,8)(4,5) \rangle,
\]
is tabulated below.
In addition, the uncolored graphs are depicted
in Figure \ref{fig:harmonic-quotient_d4reg9h-K4+C3}.

\begin{center}
\begin{tabular}{c|c|c|c}  
  $x$ & $\phi(x)$ & $m_\phi(x)$ & $\nu_\phi(x)$ \\ \hline
  0   &    3   &        1         &  1   \\ 
  1   &    3   &        1         &  1  \\ 
  2   &    1   &        1         &  0  \\ 
  3   &    0   &        1         &  0  \\ 
  4   &    2   &        1         &  1   \\
  5   &    2   &        1         &  1   \\
  6   &    1   &        1         &  0   \\
  7   &    4   &        2         &  0   \\
  8   &    0   &        1         &  0   \\ 
\end{tabular}
\end{center}

\[
  {\rm genus}(\Gamma_2) =10,\ \ \ \ \
  {\rm genus}(\Gamma_1) =4,
\]

\[
deg(\phi) =2.
\]

\begin{figure}[H]
\begin{minipage}{\textwidth}
\begin{center}
%\vspace{8.0 cm}
\includegraphics[height=5cm,width=5cm]{graphics/graph_d4reg9h.png}
\end{center}
\begin{center}
%\vspace{8.0 cm}
\includegraphics[height=4cm,width=4cm]{graphics/K4+C3-quotient-of-d4reg9h.png}
\end{center}
\end{minipage}
\caption{The quartic graph  $\Gamma_2$ ({\tt d4reg9h}), with a degree two harmonic quotient morphism to the complete house graph, $\Gamma_1=\overline{P_3\cup 2K_1}$.}
\label{fig:harmonic-quotient_d4reg9h-K4+C3}
\end{figure}

}
\end{example}


\begin{example}% d4reg9i
 {\tt d4reg9i}\newline
{\rm  
Let $\Gamma_2$ denote the graph {\tt d4reg9i} as given above.
There are $5$ harmonic quotient morphisms
$\pi: \Gamma_2\to \Gamma_1=\Gamma_2/G$,
for some subgroup $G$ of $Aut(\Gamma_2)$.
Two of these are harmonic quotient morphisms to the complete house graph.
The remaining $3$ are harmonic quotient morphisms to the paw graph.
\index{complete house graph}
\index{paw graph}

One of the these harmonic quotient morphisms to the paw graph,
$\pi: \Gamma_2\to \Gamma_1=\Gamma_2/G=\Gamma_1$,
for the subgroup

\[
G = \langle (4,5)(7,8), (0,2)(1,3)(4,7)(5,8) \rangle,
\]
is tabulated below and the morphism is depicted
in Figure \ref{fig:harmonic-quotient_d4reg9i-paw_232311011}.

\begin{center}
\begin{tabular}{c|c|c|c}  
  $x$ & $\phi(x)$ & $m_\phi(x)$ & $\nu_\phi(x)$ \\ \hline
  0   &    2   &        2         &  0   \\ 
  1   &    3   &        2         &  0  \\ 
  2   &    2   &        2         &  0  \\ 
  3   &    3   &        2         &  0  \\ 
  4   &    1   &        1         &  1   \\
  5   &    1   &        1         &  1   \\
  6   &    0   &        4         &  0   \\
  7   &    1   &        1         &  1   \\
  8   &    1   &        1         &  1   \\ 
\end{tabular}
\end{center}

\[
  {\rm genus}(\Gamma_2) =10,\ \ \ \ \
  {\rm genus}(\Gamma_1) =1,
\]

\[
deg(\phi) =4.
\]

\begin{figure}[H]
\begin{minipage}{\textwidth}
\begin{center}
%\vspace{8.0 cm}
\includegraphics[height=5cm,width=5cm]{graphics/harmonic-quotient_d4reg9i-paw_232311011.png}
\end{center}
\begin{center}
%\vspace{8.0 cm}
\includegraphics[height=4cm,width=4cm]{graphics/paw_0123.png}
\end{center}
\end{minipage}
\caption{The quartic graph  $\Gamma_2$ ({\tt d4reg9i}), with a degree four harmonic quotient morphism to $\Gamma_1={\rm paw}$.}
\label{fig:harmonic-quotient_d4reg9i-paw_232311011}
\end{figure}
}
\end{example}


\subsection{$10$ vertices}

There are (up to isomorphism) exactly $59$ $4$-regular connected graphs on $10$ vertices, all of which are Hamiltonian.
Three of these $59$ are vertex transitive,
two of those three are also edge transitive, and one of those two is also distance regular.
Of all these, $11$ have a non-trivial harmonic morphism to $C_3$.
Here, we include only a few examples (including the above-mentioned
vertex transitive graphs), some having a non-trivial harmonic
morphism to $C_3$, some not.


\begin{enumerate}
\item
{\tt d4reg10a}\newline  
Define the graph  $\Gamma$
  % d4reg10-1
  to have edge set:

\[
E = \{
(0, 1), (0, 4), (0, 8), (0, 9), (1, 3), (1, 6), (1, 7), (2, 3), (2, 5), (2, 7), (2, 9),
\]
\[
 (3, 6), (3, 7), (4, 5), (4, 8), (4, 9), (5, 6), (5, 8), (6, 7), (8, 9) 
\}.
\]
It has diameter $3$,
girth $3$, circumference $10$ (i.e., it's Hamiltonian),
critical group of size $27750 = 2\cdot 3\cdot 5^3 37$, 
and automorphism group of size $8$.

There are no non-trivial harmonic morphisms from $\Gamma$ to $P_4$,
$C_3$, $C_4$, or to $K_4$.
However, there are exactly $2$ non-trivial
harmonic morphisms from $\Gamma$ to $P_3$.

\item
  {\tt d4reg10b}\newline  
Define the graph  $\Gamma$
  % d4reg10-6
  to have edge set:

\[
E = \{
(0, 1), (0, 4), (0, 8), (0, 9), (1, 3), (1, 6), (1, 9), (2, 3), (2, 5), (2, 7), (2, 9),
\]
\[
 (3, 7), (3, 8), (4, 5), (4, 8), (4, 9), (5, 6), (5, 7), (6, 7), (6, 8)
\}.
\]
It has diameter $3$,
girth $3$, circumference $10$ (i.e., it's Hamiltonian),
critical group of size $35190 = 2\cdot 3^2 5 \cdot 17\cdot 23$, 
and automorphism group of size $4$.

There are no non-trivial harmonic morphisms from $\Gamma$ to $P_4$, $C_4$,
or to $K_4$.
However, there are exactly $4$ non-trivial
harmonic morphisms from $\Gamma$ to $P_3$
and $12$ non-trivial harmonic morphisms from $\Gamma$ to $C_3$.


\item
  {\tt d4reg10c}\newline  
  Define the graph  $\Gamma$
  % d4reg10-30
  to have edge set:

\[
E = \{
(0, 1), (0, 5), (0, 6), (0, 8), (1, 3), (1, 7), (1, 8), (2, 3), (2, 6), (2, 7), (2, 9),
\]
\[
(3, 4), (3, 9), (4, 5), (4, 8), (4, 9), (5, 6), (5, 9), (6, 7), (7, 8)
\}.
\]
It has diameter $2$,
girth $3$, circumference $10$ (i.e., it's Hamiltonian),
critical group of size $34680 = 2^33\cdot 5\cdot 17^2$, 
and automorphism group of size $2$.

There are no non-trivial harmonic morphisms from $\Gamma$ to $P_3$, $P_4$,
$C_4$, or to $K_4$.
However, there are exactly $6$ non-trivial harmonic morphisms
from $\Gamma$ to $C_3$. 

\item
  {\tt d4reg10d}\newline  
  Define the graph  $\Gamma$
  % d4reg10-34
  to have edge set:

\[
E = \{
(0, 1), (0, 5), (0, 6), (0, 8), (1, 3), (1, 6), (1, 7), (2, 3), (2, 4), (2, 7), (2, 9),
\]
\[
(3, 7), (3, 9), (4, 5), (4, 8), (4, 9), (5, 6), (5, 8), (6, 7), (8, 9)
\}.
\]
It has diameter $3$,
girth $3$, circumference $10$ (i.e., it's Hamiltonian),
critical group of size $30250 = 2\cdot 5^3 11^2$, 
and automorphism group of size $20$ (which acts transitively on the vertex set
but not the edge set).

There are no non-trivial harmonic morphisms from $\Gamma$ to $P_3$
(and therefore $P_k$, $k>3$), $C_3$, $C_4$, or to $K_4$.


\item
  {\tt d4reg10e}\newline  
  Define the graph  $\Gamma$
  % d4reg10-46
  to have edge set:

\[
E = \{
(0, 1), (0, 2), (0, 7), (0, 8), (1, 3), (1, 5), (1, 9), (2, 3), (2, 5), (2, 9),
\]
\[
(3, 4), (3, 6), (4, 5), (4, 7), (4, 8), (5, 6), (6, 7), (6, 8), (7, 9), (8, 9)
\}.
\]
It has diameter $2$,
girth $4$, circumference $10$ (i.e., it's Hamiltonian),
critical group of size $40960 =  2^{13} 5$, 
and automorphism group of size $320$ (which acts transitively on the vertex set {\it and} on the edge set). This is a symmetric (arc transitive) graph
with chromatic number $3$.
It is also a circulant graph with parameters
$(10, [1, 4, 6, 9])$ or with parameters $(10, [2, 3, 7, 8])$.
\index{circulant graph}

There are no non-trivial harmonic morphisms from $\Gamma$ to $C_3$,
$P_4$ (and therefore $P_k$, $k>4$), $K_4$.
There are $40$ non-trivial harmonic morphisms from $\Gamma$ to $P_3$,
$80$ non-trivial harmonic morphisms from $\Gamma$ to $C_4$.

\item
  {\tt d4reg10f}\newline  
  Define the graph  $\Gamma$
  % d4reg10-51
  to have edge set:

\[
E = \{
(0, 1), (0, 2), (0, 5), (0, 8), (1, 3), (1, 4), (1, 7), (2, 3), (2, 7), (2, 9),
\]
\[
(3, 5), (3, 6), (4, 5), (4, 6), (4, 8), (5, 9), (6, 7), (6, 9), (7, 8), (8, 9)
\}.
\]
It has diameter $3$,
girth $4$, circumference $10$ (i.e., it's Hamiltonian),
critical group of size $40500 = 2^2 3^4 5^3$, 
and automorphism group of size $240$ (which acts transitively on the vertex set {\it and} on the edge set). This is a bipartite, symmetric (= arc transitive), distance-regular graph.

There are no non-trivial harmonic morphisms from $\Gamma$ to $P_4$, $C_4$, or to $K_4$.
However, there are exactly $90$ non-trivial harmonic morphisms from $\Gamma$ to $C_3$ and $20$ non-trivial harmonic morphisms from $\Gamma$ to $P_3$.

\end{enumerate}



\subsubsection{Examples}



\begin{example}% d4reg10a
 {\tt d4reg10a}\newline
{\rm  
  Let $\Gamma_2$ denote the graph {\tt d4reg10a} as given above.
There are $2$ (not necessarily distinct) harmonic quotient morphisms
$\pi: \Gamma_2\to \Gamma_1=\Gamma_2/G$,
for some $G$. Both are to the $5$ vertex graph $\Gamma_1$
(denoted \verb+D~k+ in graph6 notation)
given by removing an edge from $K_5$.
\index{\verb+D~k+ graph}
An example for
\[
G = \langle (0,1)(2,5)(3,4)(6,9)(7,8) \rangle
\]
is tabulated below and the uncolored graphs are depicted in
Figure \ref{fig:harmonic-quotient_d4reg10a+K5me}.
It is a quasicover.

\begin{center}
\begin{tabular}{c|c|c|c}  
  $x$ & $\phi(x)$ & $m_\phi(x)$ & $\nu_\phi(x)$ \\ \hline
  0   &    2   &        1         &  1   \\ 
  1   &    2   &        1         &  1  \\ 
  2   &    3   &        1         &  1  \\ 
  3   &    0   &        1         &  0  \\ 
  4   &    0   &        1         &  0   \\
  5   &    3   &        1         &  1   \\
  6   &    1   &        1         &  0   \\
  7   &    4   &        1         &  0   \\
  8   &    4   &        1         &  0   \\
  9   &    1   &        1         &  0   \\     
\end{tabular}
\end{center}

\[
  {\rm genus}(\Gamma_2) =11,\ \ \ \ \
  {\rm genus}(\Gamma_1) =5,
\]

\[
deg(\phi) =2.
\]

\begin{figure}[H]
\begin{minipage}{\textwidth}
\begin{center}
%\vspace{8.0 cm}
\includegraphics[height=5cm,width=5cm]{graphics/graph_d4reg10a.png}
\end{center}
\begin{center}
%\vspace{8.0 cm}
\includegraphics[height=4cm,width=4cm]{graphics/K5me-quotient-of-d4reg10a.png}
\end{center}
\end{minipage}
\caption{The quartic graph $\Gamma_2$ ({\tt d4reg10a}), with a degree two harmonic quotient morphism to the {\tt D\~\,k}$=K_{5}-\{e\}$ graph, $\Gamma_1=\Gamma_2/G$.}
\label{fig:harmonic-quotient_d4reg10a+K5me}
\end{figure}

}
\end{example}



\begin{example}% d4reg10b
 {\tt d4reg10b}\newline
{\rm  
  Let $\Gamma_2$ denote the graph {\tt d4reg10b} as given above.
There is only $1$ harmonic quotient morphism
$\pi: \Gamma_2\to \Gamma_1=\Gamma_2/G$, for some $G$. 
It happens to be a quasicover of a wheel graph, as shown below.
\index{wheel graph}
When
\[
G = \langle (0,7)(1,3)(2,9)(4,5)(6,8) \rangle ,
\]
the quotient morphism
$\pi: \Gamma_2\to \Gamma_1=\Gamma_2/G = W_5$
is tabulated below and the uncolored graphs are depicted in
Figure \ref{fig:harmonic-quotient_d4reg10b+W5}.

\begin{center}
\begin{tabular}{c|c|c|c}  
  $x$ & $\phi(x)$ & $m_\phi(x)$ & $\nu_\phi(x)$ \\ \hline
  0   &    1   &        1         &  0   \\ 
  1   &    3   &        1         &  1  \\ 
  2   &    4   &        1         &  1  \\ 
  3   &    3   &        1         &  1  \\ 
  4   &    2   &        1         &  1   \\
  5   &    2   &        1         &  1   \\
  6   &    0   &        1         &  1   \\
  7   &    1   &        1         &  0   \\
  8   &    0   &        1         &  1   \\
  9   &    4   &        1         &  1   \\     
\end{tabular}
\end{center}

\[
  {\rm genus}(\Gamma_2) =11,\ \ \ \ \
  {\rm genus}(\Gamma_1) =4,
\]
\[
deg(\phi) =2.
\]

\begin{figure}[H]
\begin{minipage}{\textwidth}
\begin{center}
%\vspace{8.0 cm}
\includegraphics[height=5cm,width=5cm]{graphics/graph_d4reg10b.png}
\end{center}
\begin{center}
%\vspace{8.0 cm}
\includegraphics[height=4cm,width=4cm]{graphics/W5-quotient-of-d4reg10b.png}
\end{center}
\end{minipage}
\caption{The quartic graph $\Gamma_2$ ({\tt d4reg10b}), with a degree two harmonic quotient morphism to the wheel graph $W_{5}=\Gamma_1=\Gamma_2/G$.}
\label{fig:harmonic-quotient_d4reg10b+W5}
\end{figure}


}
\end{example}


\begin{example}% d4reg10c
 {\tt d4reg10c}\newline
{\rm  
  Let $\Gamma_2$ denote the graph {\tt d4reg10c} as given above.
There are no harmonic quotient morphisms
$\pi: \Gamma_2\to \Gamma_1=\Gamma_2/G$,
for some subgroup $G$ of $Aut(\Gamma_2)$.

}
\end{example}


\begin{example}% d4reg10d
 {\tt d4reg10d}\newline
{\rm  
  Let $\Gamma_2$ denote the graph {\tt d4reg10d} as given above.
There is $8$ (not necessarily distinct) harmonic quotient morphism
$\pi: \Gamma_2\to \Gamma_1=\Gamma_2/G$, for some $G$. 

One of these is a cover of a complete graph, as shown below.
\index{complete graph}
When
\[
G = \langle (0,2)(1,4)(3,5)(6,9)(7,8) \rangle ,
\]
the quotient morphism maps to the quartic graph
$\pi: \Gamma_2\to \Gamma_1=\Gamma_2/G = K_5$. This example
is tabulated below and the uncolored graphs are depicted in
Figure \ref{fig:harmonic-quotient_d4reg10d+K5}.

\begin{center}
\begin{tabular}{c|c|c|c}  
  $x$ & $\phi(x)$ & $m_\phi(x)$ & $\nu_\phi(x)$ \\ \hline
  0   &    3   &        1         &  0   \\ 
  1   &    1   &        1         &  0  \\ 
  2   &    3   &        1         &  0  \\ 
  3   &    2   &        1         &  0  \\ 
  4   &    1   &        1         &  0   \\
  5   &    2   &        1         &  0   \\
  6   &    0   &        1         &  0   \\
  7   &    4   &        1         &  0   \\
  8   &    4   &        1         &  0   \\
  9   &    0   &        1         &  0   \\     
\end{tabular}
\end{center}

\[
  {\rm genus}(\Gamma_2) =11,\ \ \ \ \
  {\rm genus}(\Gamma_1) =6,
\]
\[
deg(\phi) =2.
\]

\begin{figure}[H]
\begin{minipage}{\textwidth}
\begin{center}
%\vspace{8.0 cm}
\includegraphics[height=5cm,width=5cm]{graphics/graph_d4reg10d.png}
\end{center}
\begin{center}
%\vspace{8.0 cm}
\includegraphics[height=4cm,width=4cm]{graphics/K5-quotient-of-d4reg10d.png}
\end{center}
\end{minipage}
\caption{The quartic graph $\Gamma_2$ ({\tt d4reg10d}), with a degree two harmonic quotient morphism to the quartic graph $K_{5}=\Gamma_1=\Gamma_2/G$.}
\label{fig:harmonic-quotient_d4reg10d+K5}
\end{figure}


Six of the remaining $7$ morphisms are to the 
\verb+EUjo+ graph\footnote{In the graph6 notation.
In the ISGCI notation, it is $\overline{A}$.}.
When
\[
G = \langle (1,8)(3,9)(4,7)(5,6) \rangle ,
\]
the quotient morphism maps to the  graph
$\pi: \Gamma_2\to {\tt EUjo}=\Gamma_1=\Gamma_2/G $. This example
is tabulated below and the uncolored graphs are depicted in
Figure \ref{fig:harmonic-quotient_d4reg10d+barA}.

\begin{center}
\begin{tabular}{c|c|c|c}  
  $x$ & $\phi(x)$ & $m_\phi(x)$ & $\nu_\phi(x)$ \\ \hline
  0   &    4   &        2         &  0   \\ 
  1   &    1   &        1         &  0  \\ 
  2   &    0   &        2         &  0  \\ 
  3   &    3   &        1         &  1  \\ 
  4   &    5   &        1         &  0   \\
  5   &    2   &        1         &  1   \\
  6   &    2   &        1         &  1   \\
  7   &    5   &        1         &  0   \\
  8   &    1   &        1         &  0   \\
  9   &    3   &        1         &  1   \\     
\end{tabular}
\end{center}

\[
  {\rm genus}(\Gamma_2) =11,\ \ \ \ \
  {\rm genus}(\Gamma_1) =4,
\]
\[
deg(\phi) =2.
\]

\begin{figure}[H]
\begin{minipage}{\textwidth}
\begin{center}
%\vspace{8.0 cm}
\includegraphics[height=5cm,width=5cm]{graphics/graph_d4reg10d.png}
\end{center}
\begin{center}
%\vspace{8.0 cm}
\includegraphics[height=4cm,width=4cm]{graphics/barA-quotient-of-d4reg10d.png}
\end{center}
\end{minipage}
\caption{The quartic graph $\Gamma_2$ ({\tt d4reg10d}), with a degree two harmonic quotient morphism to the graph ${\tt EUjo}=\Gamma_1=\Gamma_2/G$.}
\label{fig:harmonic-quotient_d4reg10d+barA}
\end{figure}

The $8$th and last of these harmonic quotient morphisms is to $P_2$.

}
\end{example}


\begin{example}% d4reg10e
 {\tt d4reg10e}\newline
{\rm  
  Let $\Gamma_2$ denote the graph {\tt d4reg10e} as given above.
There are $850$ harmonic quotient morphisms
$\pi: \Gamma_2\to \Gamma_1=\Gamma_2/G$, for some $G$. 

All $20$ graphs $\Gamma_1$ arising as a quotient with
$6$ vertices are isomorphic to the graph \verb+E@~o+ (in graph6 notation). 

For instance, when
\[
G = \langle  (1,7)(2,8)(3,4)(5,6)\rangle ,
\]
the quotient morphism maps to the  graph
$\pi: \Gamma_2\to $\verb+E@~o+$=\Gamma_1=\Gamma_2/G $. This example
is tabulated below and the uncolored graphs are depicted in
Figure \ref{fig:harmonic-quotient_d4reg10e+E-AT-SIM-o}.

\begin{center}
\begin{tabular}{c|c|c|c}  
  $x$ & $\phi(x)$ & $m_\phi(x)$ & $\nu_\phi(x)$ \\ \hline
  0   &    5   &        2         &  0   \\ 
  1   &    3   &        1         &  0  \\ 
  2   &    1   &        1         &  0  \\ 
  3   &    0   &        1         &  1  \\ 
  4   &    0   &        1         &  1   \\
  5   &    2   &        1         &  1   \\
  6   &    2   &        1         &  1   \\
  7   &    3   &        1         &  0   \\
  8   &    1   &        1         &  0   \\
  9   &    4   &        2         &  0   \\     
\end{tabular}
\end{center}

\[
  {\rm genus}(\Gamma_2) =11,\ \ \ \ \
  {\rm genus}(\Gamma_1) =4,
\]
\[
deg(\phi) =2.
\]

\begin{figure}[H]
\begin{minipage}{\textwidth}
\begin{center}
%\vspace{8.0 cm}
\includegraphics[height=5cm,width=5cm]{graphics/graph_d4reg10e.png}
\end{center}
\begin{center}
%\vspace{8.0 cm}
\includegraphics[height=4cm,width=4cm]{graphics/E-AT-SIM-o-quotient-of-d4reg10e.png}
\end{center}
\end{minipage}
\caption{The quartic graph $\Gamma_2$ ({\tt d4reg10e}), with a degree two harmonic quotient morphism to the graph {\tt E@\,\~\,o}$=\Gamma_1=\Gamma_2/G$.}
\label{fig:harmonic-quotient_d4reg10e+E-AT-SIM-o}
\end{figure}


All $158$ graphs $\Gamma_1$ arising as a quotient with
$5$ vertices are isomorphic to the graph $C_5$.

For instance, when
\[
G = \langle  (0,9)(1,2)(3,5)(4,6)(7,8) \rangle ,
\]
the quotient morphism maps to the  graph
$\pi: \Gamma_2\to C_5=\Gamma_1=\Gamma_2/G $.
This multicover example is tabulated below.

\begin{center}
\begin{tabular}{c|c|c|c}  
  $x$ & $\phi(x)$ & $m_\phi(x)$ & $\nu_\phi(x)$ \\ \hline
  0   &    3   &        2         &  0   \\ 
  1   &    2   &        2         &  0  \\ 
  2   &    2   &        2         &  0  \\ 
  3   &    1   &        2         &  0  \\ 
  4   &    0   &        2         &  0   \\
  5   &    1   &        2         &  0   \\
  6   &    0   &        2         &  0   \\
  7   &    4   &        2         &  0   \\
  8   &    4   &        2         &  0   \\
  9   &    3   &        2         &  0   \\     
\end{tabular}
\end{center}

\[
  {\rm genus}(\Gamma_2) =11,\ \ \ \ \
  {\rm genus}(\Gamma_1) =1,
\]
\[
deg(\phi) =2.
\]


All $100$ graphs $\Gamma_1$ arising as a quotient with
$4$ vertices are isomorphic to the graph $C_4$.

All $540$ graphs $\Gamma_1$ arising as a quotient with
$3$ vertices are isomorphic to the graph $P_3$.

All $32$ graphs $\Gamma_1$ arising as a quotient with
$2$ vertices are isomorphic to the graph $P_2$.


}
\end{example}



\begin{example}% d4reg10f
 {\tt d4reg10f}\newline
{\rm  
  Let $\Gamma_2$ denote the graph {\tt d4reg10f} as given above.
There are $251$ harmonic quotient morphisms
$\pi: \Gamma_2\to \Gamma_1=\Gamma_2/G$, for some $G$. 

One of these $251$ harmonic quotient morphisms is a
cover of the quartic complete graph, $K_5$.
\index{complete graph}
Ten of the remaining $250$ are harmonic quotient morphisms
to a graph isomorphic to $K_5$ minus an edge - the
\verb+D~k+ graph in the graph6 notation.
\index{\verb+D~k+ graph}
Fifteen of the remaining $240$ are harmonic quotient morphisms
to a graph isomorphic to $W_5$, the wheel graph with $5$ vertices.
\index{wheel graph}

Of the remaining $225$, $45$ of the graphs $\Gamma_1$ arise as a quotient with
$3$ vertices isomorphic to the graph $C_3$.

Of the remaining $180$, $30$ of the graphs $\Gamma_1$ arise as a quotient with
$3$ vertices isomorphic to the graph $P_3$.

All of the $150$ remaining graphs $\Gamma_1$ arise as a quotient with
$3$ vertices isomorphic to the graph $P_2$.
}
\end{example}


\subsection{$11$ vertices}

There are (up to isomorphism) exactly $265$ $4$-regular connected graphs on $11$ vertices. All but one of these are Hamiltonian.
Only two of these are vertex transitive, each with an automorphism group of order $22$, one having critical
group of order $130691 = 11\cdot 109^2$, the other order $87131 = 11\cdot 89^2$.
% d4reg11-82, d4reg11-103

Of all $265$ graphs, only $48$ have a non-trivial harmonic morphism to $C_3$.
This includes an example of a graph having trivial automorphism group, yet $6$ non-trivial harmonic morphisms to $C_3$.
% d4reg11-172

If we counted correctly (``by hand''),
there are $100$ of these $265$ graphs with an odd critical group
(which therefore have no non-trivial harmonic morphisms to $C_4$ or $K_4$).
%Neither of the vertex transitive graphs have any non-trivial harmonic morphisms to $C_3$ or to $C_4$ or to $K_4$.

\begin{enumerate}
\item
  % d4reg11a
 {\tt d4reg11a}\newline
  Define the graph  $\Gamma$
  % d4reg11-3
  to have edge set:

\[
E = \{
(0, 1), (0, 5), (0, 8), (0, 10), (1, 2), (1, 3), (1, 9), (2, 3), (2, 8), (2, 9),
\]
\[
(3, 8), (3, 9), (4, 5), (4, 6), (4, 7), (4, 10), (5, 6), (5, 7), (6, 7),
\]
\[
(6, 10), (7, 10), (8, 9)
\}.
\]
It has diameter $4$,
girth $3$, circumference $6$ (thus not Hamiltonian),
% what is   sage: len(Gamma.longest_path())  ????
critical group of size $40000 = 2^6 5^4$, 
and automorphism group of size $288$. 

There are no non-trivial harmonic morphisms from $\Gamma$ to $P_4$, $C_3$, $C_4$, or to $K_4$. 
However, there are exactly $10$ non-trivial harmonic morphisms from $\Gamma$ to $P_3$. 

\item
  % d4reg11b
 {\tt d4reg11b}\newline
  Define the graph  $\Gamma$
  % d4reg11-8
  to have edge set:

\[
E = \{
(0, 1), (0, 6), (0, 9), (0, 10), (1, 2), (1, 3), (1, 5), (2, 3), (2, 7), (2, 8),
\]
\[
(3, 4), (3, 10), (4, 5), (4, 6), (4, 9), (5, 8), (5, 9), (6, 7), (6, 10),
\]
\[
(7, 8), (7, 10), (8, 9)
\}.
\]
It has diameter $3$,
girth $3$, circumference $11$ (i.e., Hamiltonian),
critical group of size $112404 = 2^2 3 \cdot 17 \cdot 19\cdot 29$, 
and automorphism group of size $2$. 

There are no non-trivial harmonic morphisms from $\Gamma$ to $P_3$, $P_4$, $C_4$, or to $K_4$.
However, there are exactly $6$ non-trivial harmonic morphisms from $\Gamma$ to $C_3$. 


\item
  % d4reg11c
 {\tt d4reg11c}\newline
  Define the graph  $\Gamma$
  % d4reg11-81 or -82
  to have edge set:

\[
E = \{
(0, 1), (0, 2), (0, 9), (0, 10), (1, 3), (1, 5), (1, 7), (2, 3), (2, 4), (2, 7),
\]
\[
(3, 6), (3, 9), (4, 5), (4, 9), (4, 10), (5, 6), (5, 8), (6, 7), (6, 10),
\]
\[
(7, 8), (8, 9), (8, 10)
\}.
\]
It has diameter $2$,
girth $4$, circumference $11$ (i.e., Hamiltonian),
critical group of size $130691 = 11\cdot 109^2$, 
and automorphism group of size $22$. This graph is vertex transitive but not
edge transitive.
\index{vertex transitive graph}

There are no non-trivial harmonic morphisms from $\Gamma$ to $P_3$, $P_4$, $C_3$, $C_4$, or to $K_4$.


\item
  % d4reg11d
 {\tt d4reg11d}\newline
  Define the graph  $\Gamma$
  % d4reg11-102
  to have edge set:

\[
E = \{
(0, 1), (0, 5), (0, 9), (0, 10), (1, 3), (1, 8), (1, 9), (2, 3), (2, 4), (2, 6),
\]
\[
(2, 7), (3, 7),  (3, 8), (4, 5), (4, 6), (4, 10), (5, 6), (5, 10), (6, 7),
\]
\[
(7, 8), (8, 9), (9, 10)
\}.
\]
It has diameter $3$,
girth $3$, circumference $11$ (i.e., Hamiltonian),
critical group of size $87131 = 11\cdot 89^2$, 
and automorphism group of size $22$. This graph is vertex transitive but not
edge transitive.
\index{vertex transitive graph}

There are no non-trivial harmonic morphisms from $\Gamma$ to $P_3$, $P_4$, $C_3$, $C_4$, or to $K_4$.


\item
  % d4reg11e
 {\tt d4reg11e}\newline
  Define the graph  $\Gamma$
  % d4reg11-171 or -172
  to have edge set:

\[
E = \{
(0, 1), (0, 4), (0, 9), (0, 10), (1, 3), (1, 5), (1, 6), (2, 3), (2, 4), (2, 7),
\]
\[
(2, 8), (3, 8), (3, 9), (4, 5), (4, 7), (5, 6), (5, 10), (6, 7), (6, 10),
\]
\[
(7, 8), (8, 9), (9, 10)
\}.
\]
It has diameter $3$,
girth $3$, circumference $11$ (i.e., Hamiltonian),
critical group of size $106056 = 2^3 3^3 491$, 
and automorphism group of size $1$. 

There are no non-trivial harmonic morphisms from $\Gamma$ to $P_3$, $P_4$, $C_4$, or to $K_4$.
However, there are exactly $6$ non-trivial harmonic morphisms from $\Gamma$ to $C_3$. 


\item
  % d4reg11f
 {\tt d4reg11f}\newline
  Define the graph  $\Gamma$
  % d4reg11-13 
  to have edge set:

\[
E = \{
(0, 1), (0, 2), (0, 4), (0, 10), (1, 2), (1, 3), (1, 9), (2, 3), (2, 9), (3, 4), 
\]
\[
(3, 10), (4, 5), (4, 6), (5, 7), (5, 8), (5, 10), (6, 7), (6, 8), (6, 10),
\]
\[
 (7, 8), (7, 9), (8, 9)
\}.
\]
It has diameter $3$,
girth $3$, circumference $11$ (i.e., Hamiltonian),
critical group of size $102400=2^{12} 5^2$, 
and automorphism group of size $64$. 

There are $9466$ non-trivial harmonic morphisms from $\Gamma$ to $P_3$ and
$16014$ non-trivial harmonic morphisms from $\Gamma$ to $P_4$.
Since the diameter is $3$, there are no non-trivial harmonic morphisms
from $\Gamma$ to $P_k$, for $k>4$, by \S 3.8 of \cite{MJ22}.
There are no non-trivial harmonic morphisms from $\Gamma$ to $C_3$ or
to $K_4$, however, there are $3488$ non-trivial harmonic morphisms
from $\Gamma$ to $C_4$.


\item
  % d4reg11g
 {\tt d4reg11g}\newline
  Define the graph  $\Gamma$
  % d4reg11-253 
  to have edge set:

\[
E = \{
(0, 1), (0, 8), (0, 9), (0, 10), (1, 3), (1, 4), (1, 6), (2, 3), (2, 4), (2, 6),
\]
\[
 (2, 10), (3, 5), (3, 7), (4, 5), (4, 7), (5, 6), (5, 9), (6, 7),
\]
\[
(7, 8), (8, 9), (8, 10), (9, 10)
\}.
\]
It has diameter $2$,
girth $3$, circumference $11$ (i.e., Hamiltonian),
critical group of size $109744=2^{4} 19^3$, 
and automorphism group of size $144$.

There are $8642$ non-trivial harmonic morphisms from $\Gamma$ to $P_3$.
Since the diameter is $2$, there are no non-trivial harmonic morphisms
from $\Gamma$ to $P_k$, for $k>3$, by \S 3.8 of \cite{MJ22}.
There are $2304$ non-trivial harmonic morphisms
from $\Gamma$ to $C_4$ but none to $K_4$.

\end{enumerate}

% 55, 80, 106, 122, 128, 219, 253, 263 have aut gp >22
% in file3-of-4reg11-graph-data.txt
% 253 has size 144


\subsubsection{Examples}



\begin{example}% d4reg11a
 {\tt d4reg11a}\newline
{\rm  
  Let $\Gamma_2$ denote the graph {\tt d4reg11a} as given above.

  There are $12$ harmonic quotient morphisms
$\pi: \Gamma_2\to \Gamma_1=\Gamma_2/G$, for some $G$. 
All of these $12$ (distinct) harmonic quotient morphisms are homomorphisms to the
graph, \verb+EJ~o+ (in the graph6 notation).
\index{\verb+EJ~o+ graph}

For instance, when
\[
G = \langle  (1,5)(2,4)(3,6)(7,9)(8,10)\rangle ,
\]
the quotient morphism maps to the  graph
$\pi: \Gamma_2\to$ {\tt EJ$\sim$o}$=\Gamma_1=\Gamma_2/G $. This example
is tabulated below and the uncolored graphs are depicted in
Figure \ref{fig:harmonic-quotient_d4reg11a+EJ-SIM-o}.

\begin{center}
\begin{tabular}{c|c|c|c}  
  $x$ & $\phi(x)$ & $m_\phi(x)$ & $\nu_\phi(x)$ \\ \hline
  0   &    3   &        2         &  0   \\ 
  1   &    2   &        1         &  0  \\ 
  2   &    1   &        1         &  0  \\ 
  3   &    4   &        1         &  0  \\ 
  4   &    1   &        1         &  0   \\
  5   &    2   &        1         &  0   \\
  6   &    4   &        1         &  0   \\
  7   &    5   &        1         &  0   \\
  8   &    0   &        1         &  0   \\
  9   &    5   &        1         &  0   \\
  10   &   0   &        1         &  0  \\   
\end{tabular}
\end{center}

\[
  {\rm genus}(\Gamma_2) =12,\ \ \ \ 
  {\rm genus}(\Gamma_1) =6,
\]
\[
deg(\phi) =2.
\]

\begin{figure}[H]
\begin{minipage}{\textwidth}
\begin{center}
%\vspace{8.0 cm}
\includegraphics[height=5cm,width=5cm]{graphics/graph_d4reg11a.png}
\end{center}
\begin{center}
%\vspace{8.0 cm}
\includegraphics[height=4cm,width=4cm]{graphics/EJ-SIM-o-quotient-of-d4reg11a.png}
\end{center}
\end{minipage}
\caption{The quartic graph $\Gamma_2$ ({\tt d4reg11a}), with a degree two harmonic quotient morphism to the graph {\tt EJ$\sim$o}$=\Gamma_1=\Gamma_2/G$.}
\label{fig:harmonic-quotient_d4reg11a+EJ-SIM-o}
\end{figure}

}
\end{example}


\begin{example}% d4reg11b
 {\tt d4reg11b}\newline
{\rm  
  Let $\Gamma_2$ denote the graph {\tt d4reg11b} as given above.
There are no harmonic quotient morphisms
$\pi: \Gamma_2\to \Gamma_1=\Gamma_2/G$,
for some subgroup $G$ of $Aut(\Gamma_2)$.

}
\end{example}



\begin{example}% d4reg11c
 {\tt d4reg11c}\newline
{\rm  
  Let $\Gamma_2$ denote the graph {\tt d4reg11c} as given above.
There are no harmonic quotient morphisms
$\pi: \Gamma_2\to \Gamma_1=\Gamma_2/G$,
for some subgroup $G$ of $Aut(\Gamma_2)$.

}
\end{example}


\begin{example}% d4reg11d
 {\tt d4reg11d}\newline
{\rm  
  Let $\Gamma_2$ denote the graph {\tt d4reg11d} as given above.
There are no harmonic quotient morphisms
$\pi: \Gamma_2\to \Gamma_1=\Gamma_2/G$,
for some subgroup $G$ of $Aut(\Gamma_2)$.

}
\end{example}


\begin{example}% d4reg11e
 {\tt d4reg11e}\newline
{\rm  
  Let $\Gamma_2$ denote the graph {\tt d4reg11e} as given above.
There are no harmonic quotient morphisms
$\pi: \Gamma_2\to \Gamma_1=\Gamma_2/G$,
for some subgroup $G$ of $Aut(\Gamma_2)$.

}
\end{example}




\begin{example}% d4reg11f
 {\tt d4reg11f}\newline
{\rm  
  Let $\Gamma_2$ denote the graph {\tt d4reg11f} as given above.
There are $30$ harmonic quotient morphisms
$\pi: \Gamma_2\to \Gamma_1=\Gamma_2/G$,
for some subgroup $G$ of $Aut(\Gamma_2)$.
In some cases, $\Gamma_1$ is a connected graph with $8$ vertices.

Of these $30$, $6$ are (not necessarily distinct)
non-trivial harmonic quotient morphisms
$\pi: \Gamma_2\to \Gamma_1=\Gamma_2/G$, for some $G$, where
$\Gamma_1$ is a graph (depicted below) which we shall
denote\footnote{I tried but could not locate it in the graph6 nor the ISGCI databases.}
$C_6+C_4+C_4$.
For instance, when
\[
G = \langle  (1,2)(4,10)(7,8)\rangle ,
\]
the quotient morphism maps to the  graph
$\pi: \Gamma_2\to \Gamma_1=\Gamma_2/G $. This example
is tabulated below and the uncolored graphs are depicted in
Figure \ref{fig:harmonic-quotient_d4reg11f+C6+C4+C4}.
\index{$C_6+C_4+C_4$ graph}

\begin{center}
\begin{tabular}{c|c|c|c}  
  $x$ & $\phi(x)$ & $m_\phi(x)$ & $\nu_\phi(x)$ \\ \hline
  0   &    7   &        2         &  0   \\ 
  1   &    1   &        1         &  1  \\ 
  2   &    1   &        1         &  1  \\ 
  3   &    4   &        2         &  0  \\ 
  4   &    0   &        2         &  0   \\
  5   &    6   &        2         &  0   \\
  6   &    5   &        2         &  0   \\
  7   &    3   &        1         &  1   \\
  8   &    3   &        1         &  1   \\
  9   &    2   &        4         &  0   \\
  10   &   0   &        2         &  0  \\   
\end{tabular}
\end{center}

\[
  {\rm genus}(\Gamma_2) =12,\ \ \ \ \
  {\rm genus}(\Gamma_1) =3,
\]
\[
deg(\phi) =2.
\]

\begin{figure}[H]
\begin{minipage}{\textwidth}
\begin{center}
%\vspace{8.0 cm}
\includegraphics[height=5cm,width=5cm]{graphics/graph_d4reg11f.png}
\end{center}
\begin{center}
%\vspace{8.0 cm}
\includegraphics[height=4cm,width=4cm]{graphics/C6+C4+C4-quotient-of-d4reg11f.png}
\end{center}
\end{minipage}
\caption{The quartic graph $\Gamma_2$ ({\tt d4reg11f}), with a degree two harmonic quotient morphism to the graph $C_6+C_4+C_4=\Gamma_1=\Gamma_2/G$.}
\label{fig:harmonic-quotient_d4reg11f+C6+C4+C4}
\end{figure}

Of the remaining $24$, $4$ are (distinct)
non-trivial harmonic quotient morphisms
$\pi: \Gamma_2\to \Gamma_1=\Gamma_2/G$, for some $G$, where
$\Gamma_1$ is the graph (depicted below) denoted
{\tt FCdzo} (in the graph6 database).
For instance, when
\[
G = \langle  (0,5)(1,7)(2,8)(3,6)\rangle ,
\]
the quotient morphism maps to the  graph
$\pi: \Gamma_2\to {\tt FCdzo}=\Gamma_1=\Gamma_2/G $. This example
is tabulated below and the uncolored graphs are depicted in
Figure \ref{fig:harmonic-quotient_d4reg11f+FCdzo}.
This is a homomorphism.
\index{FCdzo graph}

\begin{center}
\begin{tabular}{c|c|c|c}  
  $x$ & $\phi(x)$ & $m_\phi(x)$ & $\nu_\phi(x)$ \\ \hline
  0   &    6   &        1         &  0   \\ 
  1   &    2   &        1         &  0  \\ 
  2   &    1   &        1         &  0  \\ 
  3   &    4   &        1         &  0  \\ 
  4   &    5   &        2         &  0   \\
  5   &    6   &        1         &  0   \\
  6   &    4   &        1         &  0   \\
  7   &    2   &        1         &  0   \\
  8   &    1   &        1         &  0   \\
  9   &    3   &        2         &  0   \\
  10   &   0   &        2         &  0  \\   
\end{tabular}
\end{center}

\[
  {\rm genus}(\Gamma_2) =12,\ \ \ \ \
  {\rm genus}(\Gamma_1) =5,
\]
\[
deg(\phi) =2.
\]

\begin{figure}[H]
\begin{minipage}{\textwidth}
\begin{center}
%\vspace{8.0 cm}
\includegraphics[height=5cm,width=5cm]{graphics/graph_d4reg11f.png}
\end{center}
\begin{center}
%\vspace{8.0 cm}
\includegraphics[height=4cm,width=4cm]{graphics/FCdzo-quotient-of-d4reg11f.png}
\end{center}
\end{minipage}
\caption{The quartic graph $\Gamma_2$ ({\tt d4reg11f}), with a degree two harmonic quotient morphism to the graph ${\tt FCdzo}=\Gamma_1=\Gamma_2/G$.}
\label{fig:harmonic-quotient_d4reg11f+FCdzo}
\end{figure}

All of the remaining $20$ are (not necessarily distinct)
non-trivial harmonic quotient morphisms
$\pi: \Gamma_2\to \Gamma_1=\Gamma_2/G$, for some $G$, where
$\Gamma_1$ is the graph (depicted below) denoted
{\tt DrG} (in the graph6 database), also called the {\it banner graph}.
\index{DrG graph}
\index{banner graph}
For instance, when\footnote{Note the
close relationship with this example and the $C_6+C_4+C_4$ example above,
giving us another illustration of composition of quotients.}
\[
G = \langle  (1,2)(4,10)(7,8), (0,5)(1,7)(2,8)(3,6)\rangle ,
\]
the quotient morphism maps to the  graph
$\pi: \Gamma_2\to {\tt DrG}=\Gamma_1=\Gamma_2/G $. This example
is tabulated below and the uncolored graphs are depicted in
Figure \ref{fig:harmonic-quotient_d4reg11f+DrG}. 

\begin{center}
\begin{tabular}{c|c|c|c}  
  $x$ & $\phi(x)$ & $m_\phi(x)$ & $\nu_\phi(x)$ \\ \hline
  0   &    4   &        2         &  0   \\ 
  1   &    3   &        1         &  1  \\ 
  2   &    3   &        1         &  1  \\ 
  3   &    2   &        2         &  0  \\ 
  4   &    0   &        2         &  0   \\
  5   &    4   &        2         &  0   \\
  6   &    2   &        2         &  0   \\
  7   &    3   &        1         &  1   \\
  8   &    3   &        1         &  1   \\
  9   &    1   &        4         &  0   \\
  10   &   0   &        2         &  0  \\   
\end{tabular}
\end{center}

\[
  {\rm genus}(\Gamma_2) =12,\ \ \ \ \
  {\rm genus}(\Gamma_1) =1,
\]
\[
deg(\phi) =4.
\]

\begin{figure}[H]
\begin{minipage}{\textwidth}
\begin{center}
%\vspace{8.0 cm}
\includegraphics[height=5cm,width=5cm]{graphics/graph_d4reg11f.png}
\end{center}
\begin{center}
%\vspace{8.0 cm}
\includegraphics[height=4cm,width=4cm]{graphics/DrG-quotient-of-d4reg11f.png}
\end{center}
\end{minipage}
\caption{The quartic graph $\Gamma_2$ ({\tt d4reg11f}), with a degree two harmonic quotient morphism to the graph ${\tt DrG}=\Gamma_1=\Gamma_2/G$.}
\label{fig:harmonic-quotient_d4reg11f+DrG}
\end{figure}

}
\end{example}
  

\begin{example}% d4reg11g
 {\tt d4reg11g}\newline
{\rm  
  Let $\Gamma_2$ denote the graph {\tt d4reg11g} as given above.
There are $9$ harmonic quotient morphisms
$\pi: \Gamma_2\to \Gamma_1=\Gamma_2/G$,
for some subgroup $G$ of $Aut(\Gamma_2)$. In each case,
the quotient graph is the star graph with $5$ vertices:
$\Gamma_1=St_5$.
\index{star graph}

}
\end{example}

\subsection{$12$ vertices}
\label{sec:d4reg12}

There are (up to isomorphism) exactly $1544$ $4$-regular connected graphs on $12$ vertices. Ten of these are vertex transitive, but only two of those ten are also edge transitive. 


\begin{enumerate}
\item
  % d4reg12a
 {\tt d4reg12a}\newline
  Define the graph  $\Gamma$
  % d4reg12-14
  to have edge set:

\[
E = \{
(0, 1), (0, 2), (0, 7), (0, 10), (1, 3), (1, 5), (1, 6), (2, 3), (2, 9),
\]
\[
(2, 11), (3, 7), (3, 10), (4, 5), (4, 6), (4, 9), (4, 11), (5, 7), (5, 8),
\]
\[
(6, 7), (6, 8), (8, 9), (8, 11), (9, 10), (10, 11)
\}.
\]
It has diameter $3$,
girth $4$, circumference $12$ (i.e., Hamiltonian),
critical group of size $293216 = 2^5 7^2 11\cdot 17$, 
and automorphism group of size $768$. This is a vertex transitive and
edge transitive graph. It is also bipartite, symmetric
(= arc transitive), and isomorphic to the circulant graph with
parameters $(12,[1,5])$ or $(12, [1, 5, 7, 11])$.
\index{circulant graph}
\index{vertex transitive graph}
\index{edge transitive graph}
\index{arc transitive graph}
\index{bipartite transitive graph}

There are $6$ non-trivial harmonic
morphisms\footnote{Moreover, it appears all of these are multicovers
  (with constant horizontal multiplicity $2$ and constant vertical
  multiplicity $0$).}
from $\Gamma$ to $C_3$,
$272$ non-trivial harmonic
morphisms\footnote{Moreover, some, but not all, of these are multicovers
  (with constant horizontal multiplicity $2$ and constant vertical
  multiplicity $0$).} from $\Gamma$ to $C_4$,
$82$ non-trivial harmonic morphisms from $\Gamma$ to $P_3$,
$6$ non-trivial harmonic morphisms from $\Gamma$ to $P_4$,
but no non-trivial harmonic morphisms from $\Gamma$ to $K_4$.

\item
  % d4reg12b
 {\tt d4reg12b}\newline
  Define the graph  $\Gamma$
  % d4reg12-47 - 8 to C4, aut gp 24, vert trans, not edge trans
  to have edge set:

\[
E = \{
(0, 1), (0, 5), (0, 7), (0, 10), (1, 3), (1, 6), (1, 10), (2, 3), (2, 4),
\]
\[
(2, 9), (2, 11), (3, 5), (3, 8), (4, 5), (4, 6), (4, 9), (5, 8), (6, 7),
\]
\[
(6, 11), (7, 9), (7, 11), (8, 9), (8, 10), (10, 11)
\}.
\]
It has diameter $3$,
girth $3$, circumference $12$ (i.e., Hamiltonian),
critical group of size $371712 = 2^{10} 3 \cdot 11^2$, 
and automorphism group of size $24$. This is a vertex transitive graph,
but is not edge transitive. It is isomorphic to the circulant graph
having parameters $(12,[1,4])$ or $[(12, [4, 5, 7, 8])$ or
$(12, [1, 4, 8, 11])$.
\index{circulant graph}
\index{vertex transitive graph}

There are $6$ non-trivial harmonic
morphisms\footnote{This is another example where it appears
  all of these are multicovers
  (with constant horizontal multiplicity $2$ and constant vertical
  multiplicity $0$).} from $\Gamma$ to $C_3$,
$8$ non-trivial harmonic
morphisms\footnote{Moreover, it appears all of these are quasicovers
  (with constant horizontal multiplicity $1$ and constant vertical
  multiplicity $2$).}
from $\Gamma$ to $C_4$,
$10$ non-trivial harmonic morphisms from $\Gamma$ to $P_3$.
However, there
are no non-trivial harmonic morphisms from $\Gamma$ to $P_4$ or to $K_4$.
% just the test to K4 took several days!


\item
  
  % d4reg12c
 {\tt d4reg12c}\newline
  Define the graph  $\Gamma$
  % d4reg12-59 
  to have edge set:

\[
E = \{
(0, 1), (0, 2), (0, 7), (0, 10), (1, 3), (1, 5), (1, 11), (2, 3), (2, 9),
\]
\[
 (2, 11), (3, 4), (3, 10), (4, 5), (4, 6), (4, 9), (5, 7), (5, 8),
\]
\[
 (6, 7), (6, 8), (6, 10), (7, 9), (8, 9), (8, 11), (10, 11)
\}.
\]
It has diameter $3$,
girth $4$, circumference $12$ (i.e., Hamiltonian),
critical group of size $405000 = 2^3 3^4 5^4$, 
and automorphism group of size $144$. This is a vertex transitive graph,
but is not edge transitive. It is bipartite.
\index{vertex transitive graph}
\index{bipartite transitive graph}

There are $12$ non-trivial harmonic morphisms from $\Gamma$ to $P_3$
and $36$ non-trivial harmonic morphisms from $\Gamma$ to $C_3$.
However, there are no non-trivial harmonic morphisms
from $\Gamma$ to $P_4$, $C_4$, or to $K_4$.
  
\item
  
  % d4reg12d
 {\tt d4reg12d}\newline
  Define the graph  $\Gamma$
  % d4reg12-195
  to have edge set:

\[
E = \{
(0, 1), (0, 9), (0, 10), (0, 11), (1, 2), (1, 3), (1, 9), (2, 3), (2, 4), 
\]
\[
(2, 11), (3, 5), (3, 8), (4, 5), (4, 6), (4, 11), (5, 6), (5, 8), (6, 7),
\]
\[
(6, 10), (7, 8), (7, 9), (7, 10), (8, 9), (10, 11)
\}.
\]
It has diameter $3$,
girth $3$, circumference $12$,
critical group of size $331776 = 2^{12} 3^4$, 
and automorphism group of size $48$. This is a vertex transitive and
edge transitive graph. It is symmetric, has chromatic number $3$,
and is isomorphic to the cuboctahedron graph.
\index{vertex transitive graph}
\index{cuboctahedron graph}
\index{edge transitive graph}
\index{arc transitive graph}

There are exactly $240$ non-trivial harmonic morphisms from $\Gamma$ to $C_3$,
$20$ non-trivial harmonic morphisms from $\Gamma$ to $P_3$,
$24$ non-trivial harmonic morphisms from $\Gamma$ to $C_4$,
but no non-trivial harmonic morphisms from $\Gamma$ to $P_4$ or to $K_4$.

%It is unknown at this time if there are any non-trivial harmonic morphisms from $\Gamma$ to $K_4$.


\item
  
  % d4reg12e
 {\tt d4reg12e}\newline
  % d4reg12-1269 - 0 to C4, Hamiltonian, aut gp 24, vert trans, not edge trans, circ (12,[2,1])
Define the graph $\Gamma$ to have edge set

\[
E = \{
(0, 1), (0, 3), (0, 10), (0, 11), (1, 5), (1, 9), (1, 10), (2, 3), (2, 6),
\]
\[
 (2, 7), (2, 11), (3, 10), (3, 11), (4, 5), (4, 6), (4, 7), (4, 8),
\]
\[
 (5, 8), (5, 9), (6, 7), (6, 11), (7, 8), (8, 9), (9, 10)
\}.
\]
It has diameter $3$,
girth $3$, circumference $12$,
critical group of size $248832 = 2^{10} 3^5$ (which
equals\footnote{It happens, coincidentally, this formula is in wikipedia
under ``Circulant graph'' but general formulas
for the number of spanning trees in a circulant graph have
been treated by many authors. See for example \cite{Me17} and the
references cited there.} $nF_n^2$, where $n=12$), 
and automorphism group of size $24$. This is a vertex transitive but not
edge transitive graph. It has chromatic number $3$, and is isomorphic to the circulant graph with parameter $(12, [1, 2, 10, 11])$ and
$(12, [2, 5, 7, 10])$.
\index{circulant graph}
\index{vertex transitive graph}

There are $42$ non-trivial harmonic morphisms from $\Gamma$ to $C_3$
and $6$ non-trivial harmonic morphisms from $\Gamma$ to $P_3$.
However, there are no non-trivial harmonic morphisms from $\Gamma$ to $C_4$, $P_4$, or to $K_4$.  

%At this time, it's unknown if this graph has any non-trivial harmonic
%morphisms from $\Gamma$ to $K_4$.

\item
  % d4reg12f
 {\tt d4reg12f}\newline
  % d4reg12-1537 - 48 to C4, bipartite, aut gp 64,
Define the graph $\Gamma$ to have edge set

\[
E = \{
(0, 1), (0, 2), (0, 4), (0, 8), (1, 3), (1, 5), (1, 6), (2, 3), (2, 5), (2, 9),
\]
\[
(3, 7), (3, 11), (4, 5), (4, 6), (4, 10), (5, 8), (6, 7), (6, 11), 
\]
\[
(7, 9), (7, 10), (8, 9), (8, 10), (9, 11), (10, 11)
\}.
\]
It has diameter $3$,
girth $3$, circumference $12$,
critical group of size $401408 = 2^{13} 7^2$, 
and automorphism group of size $64$. This is a bipartite graph.

There are $48$ non-trivial harmonic morphisms from $\Gamma$ to $C_4$,
$32$ non-trivial harmonic morphisms from $\Gamma$ to $P_3$,
$2$ non-trivial harmonic morphisms from $\Gamma$ to $P_4$, but
there are no non-trivial harmonic morphisms from $\Gamma$ to $C_3$
or to $K_4$.

%At this time, it's unknown if this graph has any non-trivial harmonic
%morphisms from $\Gamma$ to $K_4$.

\end{enumerate}



\subsubsection{Examples}



\begin{example}% d4reg12a
\label{example:d4reg12a}  
 {\tt d4reg12a}\newline
{\rm  
Let $\Gamma_2$ denote the graph {\tt d4reg12a} as given above.
This graph, because of its large automorphism group, has a lot of
graph quotients. In fact, it has
$7195$  non-trivial harmonic quotient morphisms
$\pi: \Gamma_2\to \Gamma_1=\Gamma_2/G$,
for some subgroup $G$ of $Aut(\Gamma_2)$.

There are $12$ in the case $\Gamma_1$ has $9$ vertices,
$12$ in the case $\Gamma_1$ has $8$ vertices,
$1482$ in the case $\Gamma_1$ has $6$ vertices,
$2294$ in the case $\Gamma_1$ has $4$ vertices,
$1634$ in the case $\Gamma_1$ has $3$ vertices,
and  $1761$ in the case $\Gamma_1$ has $2$ vertices
(i.e., $\Gamma_1=P_2$).

For instance, when
\[
G = \langle  (1,7)(2,10)(4,8) \rangle ,
\]
the quotient morphism maps to the graph
$\pi: \Gamma_2\to $\verb+H??E\hw+$=\Gamma_1=\Gamma_2/G $ (in graph6 notation).
All of the $12$ quotients $\Gamma_1=\Gamma_2/G$ having $9$ vertices
are isomorphic to this graph, \verb+H??E\hw+.
This example
is tabulated below and the uncolored graphs are depicted in
Figure \ref{fig:harmonic-quotient_d4reg12a+v9e12ds444222222}. 
Note this is a graph homomorphism.

\begin{center}
\begin{tabular}{c|c|c|c}  
  $x$ & $\phi(x)$ & $m_\phi(x)$ & $\nu_\phi(x)$ \\ \hline
  0   &    8   &        2         &  0   \\ 
  1   &    2   &        1         &  0  \\ 
  2   &    1   &        1         &  0  \\ 
  3   &    4   &        2         &  0  \\ 
  4   &    7   &        1         &  0   \\
  5   &    6   &        2         &  0   \\
  6   &    5   &        2         &  0   \\
  7   &    2   &        1         &  0   \\
  8   &    7   &        1         &  0   \\
  9   &    3   &        2         &  0   \\
  10   &   1   &        1         &  0  \\   
  11   &   0   &        2         &  0  \\
\end{tabular}
\end{center}

\[
  {\rm genus}(\Gamma_2) =13,\ \ \ \ \
  {\rm genus}(\Gamma_1) =4,
\]
\[
deg(\phi) =2.
\]

\begin{figure}[H]
\begin{minipage}{\textwidth}
\begin{center}
%\vspace{8.0 cm}
\includegraphics[height=5cm,width=5cm]{graphics/graph_d4reg12a.png}
\end{center}
\begin{center}
%\vspace{8.0 cm}
\includegraphics[height=4cm,width=4cm]{graphics/v9e12ds444222222-quotient-of-d4reg12a.png}
\end{center}
\end{minipage}
\caption{The quartic graph $\Gamma_2$ ({\tt d4reg12a}), with a degree two harmonic quotient morphism to the $9$ vertex graph $\Gamma_1=\Gamma_2/G$ described above.}
\label{fig:harmonic-quotient_d4reg12a+v9e12ds444222222}
\end{figure}

% v8e12ds44442222 = G?r@xw
For instance, when
\[
G = \langle  (1,2)(5,9)(6,11)(7,10) \rangle ,
\]
the quotient morphism maps to the graph
$\pi: \Gamma_2\to $\verb+G?r@xw+$=\Gamma_1=\Gamma_2/G $ (in graph6 notation).
All of the $12$ quotients $\Gamma_1=\Gamma_2/G$ having $8$ vertices
are isomorphic to this graph, \verb+G?r@xw+.
This example
is tabulated below and the uncolored graphs are depicted in
Figure \ref{fig:harmonic-quotient_d4reg12a+v8e12ds44442222}. 
Note this is a graph homomorphism.

\begin{center}
\begin{tabular}{c|c|c|c}  
  $x$ & $\phi(x)$ & $m_\phi(x)$ & $\nu_\phi(x)$ \\ \hline
  0   &    7   &        2         &  0   \\ 
  1   &    0   &        1         &  0  \\ 
  2   &    0   &        1         &  0  \\ 
  3   &    4   &        2         &  0  \\ 
  4   &    6   &        2         &  0   \\
  5   &    5   &        1         &  0   \\
  6   &    3   &        1         &  0   \\
  7   &    1   &        1         &  0   \\
  8   &    2   &        2         &  0   \\
  9   &    5   &        1         &  0   \\
  10   &   1   &        1         &  0  \\   
  11   &   3   &        1         &  0  \\
\end{tabular}
\end{center}

\[
  {\rm genus}(\Gamma_2) =13,\ \ \ \ \
  {\rm genus}(\Gamma_1) =5,
\]
\[
deg(\phi) =2.
\]

\begin{figure}[H]
\begin{minipage}{\textwidth}
\begin{center}
%\vspace{8.0 cm}
\includegraphics[height=5cm,width=5cm]{graphics/graph_d4reg12a.png}
\end{center}
\begin{center}
%\vspace{8.0 cm}
\includegraphics[height=4cm,width=4cm]{graphics/v8e12ds44442222-quotient-of-d4reg12a.png}
\end{center}
\end{minipage}
\caption{The quartic graph $\Gamma_2$ ({\tt d4reg12a}), with a degree two harmonic quotient morphism to the $8$ vertex graph $\Gamma_1=\Gamma_2/G$ described above.}
\label{fig:harmonic-quotient_d4reg12a+v8e12ds44442222}
\end{figure}

For instance, when
\[
G = \langle (0,3)(1,7)(2,10)(4,8)(5,6)(9,11) \rangle ,
\]
the quotient morphism maps (is a $2$-multicover, in fact) to the cycle graph
$\pi: \Gamma_2\to C_6=\Gamma_1=\Gamma_2/G $.
All of the $1482$ quotients $\Gamma_1=\Gamma_2/G$ having $6$ vertices
are isomorphic to this cycle graph $C_6$.
This example
is tabulated below and the uncolored graphs are depicted in
Figure \ref{fig:harmonic-quotient_d4reg12a+C6}. 

\begin{center}
\begin{tabular}{c|c|c|c}  
  $x$ & $\phi(x)$ & $m_\phi(x)$ & $\nu_\phi(x)$ \\ \hline
  0   &    4   &        2         &  0   \\ 
  1   &    3   &        2         &  0  \\ 
  2   &    2   &        2         &  0  \\ 
  3   &    4   &        2         &  0  \\ 
  4   &    5   &        2         &  0   \\
  5   &    0   &        2         &  0   \\
  6   &    0   &        2         &  0   \\
  7   &    3   &        2         &  0   \\
  8   &    5   &        2         &  0   \\
  9   &    1   &        2         &  0   \\
  10   &   2   &        2         &  0  \\   
  11   &   1   &        2         &  0  \\
\end{tabular}
\end{center}

\[
  {\rm genus}(\Gamma_2) =13,\ \ \ \ \
  {\rm genus}(\Gamma_1) =1,
\]
\[
deg(\phi) =2.
\]

\begin{figure}[H]
\begin{minipage}{\textwidth}
\begin{center}
%\vspace{8.0 cm}
\includegraphics[height=5cm,width=5cm]{graphics/graph_d4reg12a.png}
\end{center}
\begin{center}
%\vspace{8.0 cm}
\includegraphics[height=4cm,width=4cm]{graphics/C6-quotient-of-d4reg12a.png}
\end{center}
\end{minipage}
\caption{The quartic graph $\Gamma_2$ ({\tt d4reg12a}), with a degree two harmonic quotient morphism to the graph $C_6=\Gamma_1=\Gamma_2/G$ described above.}
\label{fig:harmonic-quotient_d4reg12a+C6}
\end{figure}

For instance, when
\[
G = \langle (0,5,9)(1,4,2)(3,6,11)(7,8,10) \rangle ,
\]
the quotient morphism maps (is a $2$-multicover, in fact) to the cycle graph
$\pi: \Gamma_2\to C_4=\Gamma_1=\Gamma_2/G $.
All of the $1482$ quotients $\Gamma_1=\Gamma_2/G$ having $4$ vertices
are isomorphic to this cycle graph $C_4$.
This example
is tabulated below and the uncolored graphs are depicted in
Figure \ref{fig:harmonic-quotient_d4reg12a+C4}. 

\begin{center}
\begin{tabular}{c|c|c|c}  
  $x$ & $\phi(x)$ & $m_\phi(x)$ & $\nu_\phi(x)$ \\ \hline
  0   &    2   &        2         &  0   \\ 
  1   &    0   &        2         &  0  \\ 
  2   &    0   &        2         &  0  \\ 
  3   &    3   &        2         &  0  \\ 
  4   &    0   &        2         &  0   \\
  5   &    2   &        2         &  0   \\
  6   &    3   &        2         &  0   \\
  7   &    1   &        2         &  0   \\
  8   &    1   &        2         &  0   \\
  9   &    2   &        2         &  0   \\
  10   &   1   &        2         &  0  \\   
  11   &   3   &        2         &  0  \\
\end{tabular}
\end{center}

\[
  {\rm genus}(\Gamma_2) =13,\ \ \ \ \
  {\rm genus}(\Gamma_1) =1,
\]
\[
deg(\phi) =3.
\]

\begin{figure}[H]
\begin{minipage}{\textwidth}
\begin{center}
%\vspace{8.0 cm}
\includegraphics[height=5cm,width=5cm]{graphics/graph_d4reg12a.png}
\end{center}
\begin{center}
%\vspace{8.0 cm}
\includegraphics[height=4cm,width=4cm]{graphics/C4-quotient-of-d4reg12a.png}
\end{center}
\end{minipage}
\caption{The quartic graph $\Gamma_2$ ({\tt d4reg12a}), with a degree three harmonic quotient morphism to the graph $C_4=\Gamma_1=\Gamma_2/G$ described above.}
\label{fig:harmonic-quotient_d4reg12a+C4}
\end{figure}

All of the $1634$ quotients $\Gamma_1=\Gamma_2/G$ having $3$ vertices
are isomorphic to this cycle graph $P_3$.

All of the $1761$ quotients $\Gamma_1=\Gamma_2/G$ having $2$ vertices
are isomorphic to this cycle graph $P_2$.
}
\end{example}



\begin{example}% d4reg12b
 {\tt d4reg12b}\newline
{\rm  
  Let $\Gamma_2$ denote the graph {\tt d4reg12b} as given above.
It has $21$  non-trivial harmonic quotient morphisms
$\pi: \Gamma_2\to \Gamma_1=\Gamma_2/G$,
for some subgroup $G$ of $Aut(\Gamma_2)$.
There are $6$ in the case $\Gamma_1$ has $7$ vertices,
only $1$ in the case $\Gamma_1$ has $6$ vertices,
$4$ in the case $\Gamma_1$ has $4$ vertices,
$3$ in the case $\Gamma_1$ has $3$ vertices,
and  $7$ in the case $\Gamma_1$ has $2$ vertices
(i.e., $\Gamma_1=P_2$).


For instance, when
\[
G = \langle (1,10)(3,11)(4,9)(5,7)(6,8) \rangle ,
\]
the quotient morphism $\pi$ maps to the $7$ vertex graph
$\pi: \Gamma_2\to {\tt FKMiw}=\Gamma_1=\Gamma_2/G $ (in the graph6 notation).
In fact, every harmonic quotient having $7$ vertices is isomorphic to
this graph. This example
is tabulated below and the uncolored graphs are depicted in
Figure \ref{fig:harmonic-quotient_d4reg12b+FKMiw}. 

\begin{center}
\begin{tabular}{c|c|c|c}  
  $x$ & $\phi(x)$ & $m_\phi(x)$ & $\nu_\phi(x)$ \\ \hline
  0   &    2   &        2         &  0   \\ 
  1   &    5   &        1         &  1  \\ 
  2   &    0   &        2         &  0  \\ 
  3   &    4   &        1         &  0  \\ 
  4   &    3   &        1         &  1   \\
  5   &    6   &        1         &  0   \\
  6   &    1   &        1         &  0   \\
  7   &    6   &        1         &  0   \\
  8   &    1   &        1         &  0   \\
  9   &    3   &        1         &  1   \\
  10   &   5   &        1         &  1  \\   
  11   &   4   &        1         &  0  \\
\end{tabular}
\end{center}

\[
  {\rm genus}(\Gamma_2) =13,\ \ \ \ \
  {\rm genus}(\Gamma_1) =5,
\]
\[
deg(\phi) =2.
\]

\begin{figure}[H]
\begin{minipage}{\textwidth}
\begin{center}
%\vspace{8.0 cm}
\includegraphics[height=5cm,width=5cm]{graphics/graph_d4reg12b.png}
\end{center}
\begin{center}
%\vspace{8.0 cm}
\includegraphics[height=4cm,width=4cm]{graphics/FKMiw-quotient-of-d4reg12b.png}
\end{center}
\end{minipage}
\caption{The quartic graph $\Gamma_2$ ({\tt d4reg12b}), with a degree two harmonic quotient morphism to the graph {\tt FKMiw}$=\Gamma_1=\Gamma_2/G$ described above.}
\label{fig:harmonic-quotient_d4reg12b+FKMiw}
\end{figure}


For instance, when
\[
G = \langle (0,2)(1,9)(3,7)(4,10)(5,11)(6,8) \rangle ,
\]
the quotient morphism $\pi$ maps (as a cover, in fact) to the quartic graph
$\pi: \Gamma_2\to $\verb+E]~o+$=\Gamma_1=\Gamma_2/G$ (in the graph6 notation).
In fact, this graph \verb+E]~o+ is a circulant graph with parameters
$(6, [1, 2, 4, 5])$. This example
is tabulated below and the uncolored graphs are depicted in
Figure \ref{fig:harmonic-quotient_d4reg12b+circ6-1245}.
\index{circulant graph}
  
\begin{center}
\begin{tabular}{c|c|c|c}  
  $x$ & $\phi(x)$ & $m_\phi(x)$ & $\nu_\phi(x)$ \\ \hline
  0   &    5   &        1         &  0   \\ 
  1   &    2   &        1         &  0  \\ 
  2   &    5   &        1         &  0  \\ 
  3   &    4   &        1         &  0  \\ 
  4   &    0   &        1         &  0   \\
  5   &    1   &        1         &  0   \\
  6   &    3   &        1         &  0   \\
  7   &    4   &        1         &  0   \\
  8   &    3   &        1         &  0   \\
  9   &    2   &        1         &  0   \\
  10   &   0   &        1         &  0  \\   
  11   &   1   &        1         &  0  \\
\end{tabular}
\end{center}

\[
  {\rm genus}(\Gamma_2) = 13,\ \ \ \ \
  {\rm genus}(\Gamma_1) = 7,
\]
\[
deg(\phi) =2.
\]

\begin{figure}[H]
\begin{minipage}{\textwidth}
\begin{center}
%\vspace{8.0 cm}
\includegraphics[height=5cm,width=5cm]{graphics/graph_d4reg12b.png}
\end{center}
\begin{center}
%\vspace{8.0 cm}
\includegraphics[height=4cm,width=4cm]{graphics/circ6-1245-quotient-of-d4reg12b.png}
\end{center}
\end{minipage}
\caption{The quartic graph $\Gamma_2$ ({\tt d4reg12b}), with a degree two harmonic quotient morphism to the circulant graph $\Gamma_1=\Gamma_2/G$ described above.}
\label{fig:harmonic-quotient_d4reg12b+circ6-1245}
\end{figure}

For instance, when
\[
G = \langle (0,1,10)(2,9,4)(3,8,5)(6,11,7) \rangle ,
\]
the quotient morphism $\pi$ maps to the $4$ vertex graph
$\pi: \Gamma_2\to C_4=\Gamma_1=\Gamma_2/G $.
In fact, every harmonic quotient having $4$ vertices is isomorphic to
either this graph $C_4$ or to the diamond graph. This $C_4$ example
is tabulated below.
% and the uncolored graphs are depicted in
% Figure \ref{fig:harmonic-quotient_d4reg12b+C4}. 

\begin{center}
\begin{tabular}{c|c|c|c}  
  $x$ & $\phi(x)$ & $m_\phi(x)$ & $\nu_\phi(x)$ \\ \hline
  0   &    1   &        1         &  2   \\ 
  1   &    1   &        1         &  2  \\ 
  2   &    2   &        1         &  2  \\ 
  3   &    3   &        1         &  2  \\ 
  4   &    2   &        1         &  2   \\
  5   &    3   &        1         &  2   \\
  6   &    0   &        1         &  2   \\
  7   &    0   &        1         &  2   \\
  8   &    3   &        1         &  2   \\
  9   &    2   &        1         &  2   \\
  10   &   1   &        1         &  2  \\   
  11   &   0   &        1         &  2  \\
\end{tabular}
\end{center}

\[
  {\rm genus}(\Gamma_2) =13,\ \ \ \ \
  {\rm genus}(\Gamma_1) =1,
\]
\[
deg(\phi) =3.
\]

In fact, every harmonic quotient having $3$ vertices is isomorphic to
either $P_3$ (two of them) or $C_3$ (one of them).

Of course, every one of the $7$ harmonic quotients $\Gamma_1$ having $2$
vertices is isomorphic to $P_2$.

}
\end{example}



\begin{example}% d4reg12c
 {\tt d4reg12c}\newline
{\rm  
  Let $\Gamma_2$ denote the graph {\tt d4reg12c} as given above.
It has $142$  non-trivial harmonic quotient morphisms
$\pi: \Gamma_2\to \Gamma_1=\Gamma_2/G$,
for some subgroup $G$ of $Aut(\Gamma_2)$.
There are $22$ in the case $\Gamma_1$ has $6$ vertices,
$6$ in the case $\Gamma_1$ has $4$ vertices,
$6$ in the case $\Gamma_1$ has $3$ vertices,
and  $108$ in the case $\Gamma_1$ has $2$ vertices
(i.e., $\Gamma_1=P_2$).

For those quotients having $6$ vertices, the following graphs arise
(in graph6 notation): \verb+ELv_+ (= the prism graph),
\verb+EFz_+ (= $K_{3,3}$), \verb+Ejmw+, \verb+Es\w+. 
\index{prism graph}
\index{complete bipartite graph}

For instance, when
\[
G = \langle (0,1)(2,3)(4,9)(5,7)(6,8)(10,11) \rangle ,
\]
the quotient morphism $\pi$ (a quasicover) maps to the $6$ vertex prism graph
$\pi: \Gamma_2\to \Gamma_1=\Gamma_2/G $, as tabulated below.

\begin{center}
\begin{tabular}{c|c|c|c}  
  $x$ & $\phi(x)$ & $m_\phi(x)$ & $\nu_\phi(x)$ \\ \hline
  0   &    3   &        1         &  1   \\ 
  1   &    3   &        1         &  1  \\ 
  2   &    1   &        1         &  1  \\ 
  3   &    1   &        1         &  1  \\ 
  4   &    2   &        1         &  1   \\
  5   &    5   &        1         &  1   \\
  6   &    0   &        1         &  1   \\
  7   &    5   &        1         &  1   \\
  8   &    0   &        1         &  1   \\
  9   &    2   &        1         &  1   \\
  10   &   4   &        1         &  1  \\   
  11   &   4   &        1         &  1  \\
\end{tabular}
\end{center}

\[
  {\rm genus}(\Gamma_2) =13,\ \ \ \ \
  {\rm genus}(\Gamma_1) =4,
\]
\[
deg(\phi) =2.
\]

If you replace one of the rectangular faces (a $C_4$)
of the prism graph with a copy of $K_4$, then you get the
graph denoted {\tt Ejmw} (in the graph6 notation).
For instance, when
\[
G = \langle (0,4)(1,5)(2,6)(3,7)(8,11)(9,10) \rangle ,
\]
the quotient morphism $\pi$ (a quasicover) maps to the $6$ vertex prism graph
$\pi: \Gamma_2\to \Gamma_1=\Gamma_2/G $, as tabulated below.
The uncolored graphs are depicted in
Figure \ref{fig:harmonic-quotient_d4reg12c+Ejmw}. 

\begin{center}
\begin{tabular}{c|c|c|c}  
  $x$ & $\phi(x)$ & $m_\phi(x)$ & $\nu_\phi(x)$ \\ \hline
  0   &    5   &        1         &  0   \\ 
  1   &    1   &        1         &  1  \\ 
  2   &    4   &        1         &  0  \\ 
  3   &    3   &        1         &  0  \\ 
  4   &    5   &        1         &  0   \\
  5   &    1   &        1         &  1   \\
  6   &    4   &        1         &  0   \\
  7   &    3   &        1         &  0   \\
  8   &    0   &        1         &  1   \\
  9   &    2   &        1         &  0   \\
  10   &   2   &        1         &  0  \\   
  11   &   0   &        1         &  1  \\
\end{tabular}
\end{center}

\[
  {\rm genus}(\Gamma_2) =13,\ \ \ \ \
  {\rm genus}(\Gamma_1) =6,
\]
\[
deg(\phi) =2.
\]

\begin{figure}[H]
\begin{minipage}{\textwidth}
\begin{center}
%\vspace{8.0 cm}
\includegraphics[height=5cm,width=5cm]{graphics/graph_d4reg12c.png}
\end{center}
\begin{center}
%\vspace{8.0 cm}
\includegraphics[height=4cm,width=4cm]{graphics/Ejmw-quotient-of-d4reg12c.png}
\end{center}
\end{minipage}
\caption{The quartic graph $\Gamma_2$ ({\tt d4reg12c}), with a degree two harmonic quotient morphism (indeed, a quasicover) to the graph {\tt Ejmw}$=\Gamma_1=\Gamma_2/G$ described above.}
\label{fig:harmonic-quotient_d4reg12c+Ejmw}
\end{figure}

All $6$ of the non-trivial harmonic quotient morphisms
$\pi: \Gamma_2\to \Gamma_1=\Gamma_2/G$, in the case where
$\Gamma_1$ has $4$ vertices, satisfy $\Gamma_4=St_4$.
\index{star graph}

All $6$ of the non-trivial harmonic quotient morphisms
$\pi: \Gamma_2\to \Gamma_1=\Gamma_2/G$, in the case where
$\Gamma_1$ has $3$ vertices, satisfy $\Gamma_4=C_3$.
In fact, all these morphisms are quasicovers.
For instance, when
\[
G = \langle  (0,1)(2,3)(4,9)(5,7)(6,8)(10,11),
(0,7)(1,5)(2,9)(3,4)(6,10)(8,11)\rangle ,
\]
the quotient morphism $\pi$ (a quasicover) maps to the cycle graph
$\pi: \Gamma_2\to C_3=\Gamma_1=\Gamma_2/G $, as tabulated below
and depicted in
Figure \ref{fig:harmonic-quotient_d4reg12c+C3}. 

\begin{center}
\begin{tabular}{c|c|c|c}  
  $x$ & $\pi(x)$ & $m_\pi(x)$ & $\nu_\pi(x)$   \\
  0   &    1   &        1         &  2   \\   
  1   &    1   &        1         &  2  \\ 
  2   &    2   &        1         &  2  \\ 
  3   &    2   &        1         &  2  \\ 
  4   &    2   &        1         &  2   \\
  5   &    1   &        1         &  2   \\
  6   &    0   &        1         &  2   \\
  7   &    1   &        1         &  2   \\
  8   &    0   &        1         &  2   \\
  9   &    2   &        1         &  2   \\
  10   &   0   &        1         &  2  \\   
  11   &   0   &        1         &  2  \\
\end{tabular}
\end{center}

\[
  {\rm genus}(\Gamma_2) =13,\ \ \ \ \
  {\rm genus}(\Gamma_1) =1,
\]
\[
deg(\phi) =4.
\]

\begin{figure}[H]
\begin{minipage}{\textwidth}
\begin{center}
%\vspace{8.0 cm}
\includegraphics[height=5cm,width=5cm]{graphics/graph_d4reg12c_112221010200.png}
\end{center}
\begin{center}
%\vspace{8.0 cm}
\includegraphics[height=4cm,width=4cm]{graphics/cycle3-012.png}
\end{center}
\end{minipage}
\caption{The quartic graph $\Gamma_2$ ({\tt d4reg12c}), with a degree four harmonic quotient morphism (indeed, a quasicover) to the graph $C_3=\Gamma_1=\Gamma_2/G$.}
\label{fig:harmonic-quotient_d4reg12c+C3}
\end{figure}



The remaining $108$ non-trivial harmonic quotient morphisms,
in the case $\Gamma_1$ has $2$ vertices, all map to
$\Gamma_1=P_2$.
}
\end{example}



\begin{example}% d4reg12d
  {\tt d4reg12d}\newline
\label{example:d4reg12d}  
{\rm  
  Let $\Gamma_2$ denote the graph {\tt d4reg12d} as given above.
It has $79$  non-trivial harmonic quotient morphisms
$\pi: \Gamma_2\to \Gamma_1=\Gamma_2/G$,
for some subgroup $G$ of $Aut(\Gamma_2)$.
There are $3$ in the case $\Gamma_1$ has $8$ vertices,
$12$ in case $\Gamma_1$ has $7$ vertices,
only $1$ in case $\Gamma_1$ has $6$ vertices,
$9$ in case $\Gamma_1$ has $5$ vertices,
$9$ in case $\Gamma_1$ has $4$ vertices,
$25$ in the case $\Gamma_1$ has $3$ vertices,
and  $20$ in the case $\Gamma_1$ has $2$ vertices
(i.e., $\Gamma_1=P_2$).

All $3$ of the non-trivial harmonic quotient morphisms
$\pi: \Gamma_2\to \Gamma_1=\Gamma_2/G$, in the case where
$\Gamma_1$ has $8$ vertices, satisfy $\Gamma_1=$\verb+G?qipk+
(in the graph6 notation).
%\index{star graph}
For instance, when
\[
G = \langle  (1,11)(3,4)(6,8)(9,10) \rangle ,
\]
the quotient morphism $\pi$ (a homomorphism) maps to
$\Gamma_1=\Gamma_2/G =$\verb+G?qipk+
is tabulated below
and the uncolored graphs are depicted in
Figure \ref{fig:harmonic-quotient_d4reg12d+G_QU_qipk}. 

\begin{center}
\begin{tabular}{c|c|c|c}  
  $x$ & $\pi(x)$ & $m_\pi(x)$ & $\nu_\pi(x)$   \\ 
  0   &    7   &        2         &  0  \\ 
  1   &    3   &        1         &  0  \\ 
  2   &    1   &        2         &  0  \\ 
  3   &    0   &        1         &  0   \\
  4   &    0   &        1         &  0   \\
  5   &    6   &        2         &  0   \\
  6   &    5   &        1         &  0   \\
  7   &    4   &        2         &  0   \\
  8   &    5   &        1         &  0   \\
  9   &    2   &        1         &  0  \\   
  10  &    2  &        1         &  0  \\
  11  &    3   &        1         &  0  \\  
\end{tabular}
\end{center}

\[
  {\rm genus}(\Gamma_2) =13,\ \ \ \ \
  {\rm genus}(\Gamma_1) =5,
\]
\[
deg(\phi) =2.
\]

\begin{figure}[H]
\begin{minipage}{\textwidth}
\begin{center}
%\vspace{8.0 cm}
\includegraphics[height=5cm,width=5cm]{graphics/graph_d4reg12d2.png}
\end{center}
\begin{center}
%\vspace{8.0 cm}
\includegraphics[height=4cm,width=4cm]{graphics/G_QU_qipk-quotient-of-d4reg12d.png}
\end{center}
\end{minipage}
\caption{The quartic graph $\Gamma_2$ ({\tt d4reg12d}), with a degree two harmonic quotient morphism (indeed, a homomorphism) to the graph $\Gamma_1=\Gamma_2/G$ as described above.}
\label{fig:harmonic-quotient_d4reg12d+G_QU_qipk}
\end{figure}

Of the $12$ non-trivial harmonic quotient morphisms
$\pi: \Gamma_2\to \Gamma_1=\Gamma_2/G$, in the case where
$\Gamma_1$ has $7$ vertices, exactly $6$ of them
satisfy $\Gamma_1=$\verb+FIe`w+. (These quotients have
vertical edges.)
The remaining $6$ satisfy $\Gamma_1=$\verb+FB]lg+
(in the graph6 notation).
%\index{star graph}
For instance, when
\[
G = \langle (1,10)(2,7)(3,6)(4,8)(9,11) \rangle ,
\]
the quotient morphism $\pi$ (a homomorphism) maps to
$\Gamma_1=\Gamma_2/G =$\verb+FB]lg+
is tabulated below
and the uncolored graphs are depicted in
Figure \ref{fig:harmonic-quotient_d4reg12d+FB_BR_lg}. 

\begin{center}
\begin{tabular}{c|c|c|c}  
  $x$ & $\pi(x)$ & $m_\pi(x)$ & $\nu_\pi(x)$   \\ 
  0   &    6   &        2         &  0  \\ 
  1   &    4   &        1         &  0  \\ 
  2   &    0   &        1         &  0  \\ 
  3   &    3   &        1         &  0   \\
  4   &    5   &        1         &  0   \\
  5   &    2   &        2         &  0   \\
  6   &    3   &        1         &  0   \\
  7   &    0   &        1         &  0   \\
  8   &    5   &        1         &  0   \\
  9   &    1   &        1         &  0  \\   
  10  &    4  &         1         &  0  \\
  11  &    1   &        1         &  0  \\  
\end{tabular}
\end{center}

\[
  {\rm genus}(\Gamma_2) =13,\ \ \ \ \
  {\rm genus}(\Gamma_1) =6,
\]
\[
deg(\phi) =2.
\]

\begin{figure}[H]
\begin{minipage}{\textwidth}
\begin{center}
%\vspace{8.0 cm}
\includegraphics[height=5cm,width=5cm]{graphics/graph_d4reg12d2.png}
\end{center}
\begin{center}
%\vspace{8.0 cm}
\includegraphics[height=4cm,width=4cm]{graphics/FB_BR_lg-quotient-of-d4reg12d.png}
\end{center}
\end{minipage}
\caption{The quartic graph $\Gamma_2$ ({\tt d4reg12d}), with a degree two harmonic quotient morphism (indeed, a homomorphism) to the graph $\Gamma_1=\Gamma_2/G$ as described above.}
\label{fig:harmonic-quotient_d4reg12d+FB_BR_lg}
\end{figure}


The only non-trivial harmonic quotient morphism
$\pi: \Gamma_2\to \Gamma_1=\Gamma_2/G$, in the case where
$\Gamma_1$ has $6$ vertices, satisfies $\Gamma_1={\tt Oct}$,
the quartic octahedral graph. (This quotient is a cover.)
\index{octahedral graph}
When
\[
G = \langle (0,5)(1,6)(2,7)(3,10)(4,9)(8,11) \rangle ,
\]
the quotient morphism $\pi$ maps to
$\Gamma_1=\Gamma_2/G ={\tt Oct}$
is tabulated below
and the uncolored graphs are depicted in
Figure \ref{fig:harmonic-quotient_d4reg12d+octahedral}. 

\begin{center}
\begin{tabular}{c|c|c|c}  
  $x$ & $\pi(x)$ & $m_\pi(x)$ & $\nu_\pi(x)$   \\ 
  0   &    5   &        1         &  0  \\ 
  1   &    4   &        1         &  0  \\ 
  2   &    2   &        1         &  0  \\ 
  3   &    3   &        1         &  0   \\
  4   &    1   &        1         &  0   \\
  5   &    5   &        1         &  0   \\
  6   &    4   &        1         &  0   \\
  7   &    2   &        1         &  0   \\
  8   &    0   &        1         &  0   \\
  9   &    1   &        1         &  0  \\   
  10  &    3  &         1         &  0  \\
  11  &    0   &        1         &  0  \\  
\end{tabular}
\end{center}

\[
  {\rm genus}(\Gamma_2) =13,\ \ \ \ \
  {\rm genus}(\Gamma_1) =7,
\]
\[
deg(\phi) =2.
\]

\begin{figure}[H]
\begin{minipage}{\textwidth}
\begin{center}
%\vspace{8.0 cm}
\includegraphics[height=5cm,width=5cm]{graphics/graph_d4reg12d2.png}
\end{center}
\begin{center}
%\vspace{8.0 cm}
\includegraphics[height=4cm,width=4cm]{graphics/octahedral-quotient-of-d4reg12d.png}
\end{center}
\end{minipage}
\caption{The quartic graph $\Gamma_2$ ({\tt d4reg12d}), with a degree two harmonic quotient morphism (indeed, a cover) to the graph $\Gamma_1=\Gamma_2/G={\tt Oct}$.}
\label{fig:harmonic-quotient_d4reg12d+octahedral}
\end{figure}

Of the $9$ non-trivial harmonic quotient morphisms
$\pi: \Gamma_2\to \Gamma_1=\Gamma_2/G$, in the case where
$\Gamma_1$ has $5$ vertices, exactly $3$ of them
satisfy $\Gamma_1={\tt butterfly}$, the butterfly or bowtie graph.
(These quotients have no vertical edges.)
\index{butterfly graph}
\index{bowtie graph}
The remaining $6$ satisfy $\Gamma_1={\tt bull}$, the bull graph,
which can be invisioned as $C_3$ with two of the vertices having a
leaf\footnote{So if these leaves are pointing straight up
  and the degree 2 vertices is blwlow them, then the leafs look like
  horns and the degree 2 vertex looks like a chin.}.
\index{bull graph}
(These quotients do have vertical edges.)
When
\[
G = \langle (1,11)(3,4)(6,8)(9,10), (0,2)(1,11)(3,10)(4,9)(5,7)(6,8) \rangle ,
\]
the quotient morphism $\pi$ maps to
$\Gamma_1=\Gamma_2/G ={\tt butterfly}$
is tabulated below
and the uncolored graphs are depicted in
Figure \ref{fig:harmonic-quotient_d4reg12d+butterfly}. 

\begin{center}
\begin{tabular}{c|c|c|c}  
  $x$ & $\pi(x)$ & $m_\pi(x)$ & $\nu_\pi(x)$   \\ 
  0   &    5   &        2         &  0  \\ 
  1   &    4   &        2         &  0  \\ 
  2   &    2   &        2         &  0  \\ 
  3   &    3   &        1         &  0   \\
  4   &    1   &        1         &  0   \\
  5   &    5   &        2         &  0   \\
  6   &    4   &        2         &  0   \\
  7   &    2   &        2         &  0   \\
  8   &    0   &        2         &  0   \\
  9   &    1   &        1         &  0  \\   
  10  &    3  &         1         &  0  \\
  11  &    0   &        2         &  0  \\  
\end{tabular}
\end{center}

\[
  {\rm genus}(\Gamma_2) =13,\ \ \ \ \
  {\rm genus}(\Gamma_1) =2,
\]
\[
deg(\phi) =4.
\]

\begin{figure}[H]
\begin{minipage}{\textwidth}
\begin{center}
%\vspace{8.0 cm}
\includegraphics[height=5cm,width=5cm]{graphics/graph_d4reg12d2.png}
\end{center}
\begin{center}
%\vspace{8.0 cm}
\includegraphics[height=4cm,width=4cm]{graphics/butterfly-quotient-of-d4reg12d.png}
\end{center}
\end{minipage}
\caption{The quartic graph $\Gamma_2$ ({\tt d4reg12d}), with a degree four harmonic quotient morphism (indeed, a homomorphism) to the graph $\Gamma_1=\Gamma_2/G={\tt butterfly}$.}
\label{fig:harmonic-quotient_d4reg12d+butterfly}
\end{figure}

Of the $9$ non-trivial harmonic quotient morphisms
$\pi: \Gamma_2\to \Gamma_1=\Gamma_2/G$, in the case where
$\Gamma_1$ has $4$ vertices, exactly $3$ of them
satisfy $\Gamma_1=C_4$, the cycle graph. 
The remaining $6$ satisfy $\Gamma_1={\tt diamond}$, the diamond graph.
\index{diamond graph}
(In each case, these quotients have vertical edges.)


Of the $25$ non-trivial harmonic quotient morphisms
$\pi: \Gamma_2\to \Gamma_1=\Gamma_2/G$, in the case where
$\Gamma_1$ has $3$ vertices, exactly $13$ of them
satisfy $\Gamma_1=P_3$, the path graph.
(In each case, these quotients have vertical edges.)
The remaining $6$ satisfy $\Gamma_1=C_3$, the cycle graph.
The quotients in this case $\pi: \Gamma_2\to C_3=\Gamma_2/G$
are multicovers of degree 4 and multiplicity $m=2$.
When
\[
G = \langle 
(0,2)(1,11)(3,10)(4,9)(5,7)(6,8), (0,5)(1,8)(2,7)(3,9)(4,10)(6,11)\rangle ,
\]
the quotient morphism $\pi$ maps to
$\Gamma_1=\Gamma_2/G =C_3$
is tabulated below
and  depicted in Figure \ref{fig:harmonic-quotient_d4reg12d-C3_212002121001}.

\begin{center}
\begin{tabular}{c|c|c|c}  
  $x$ & $\pi(x)$ & $m_\pi(x)$ & $\nu_\pi(x)$   \\ 
  0   &    2   &        2         &  0  \\ 
  1   &    1   &        2         &  0  \\ 
  2   &    2   &        2         &  0  \\ 
  3   &    0   &        2         &  0   \\
  4   &    0   &        2         &  0   \\
  5   &    2   &        2         &  0   \\
  6   &    1   &        2         &  0   \\
  7   &    2   &        2         &  0   \\
  8   &    1   &        2         &  0   \\
  9   &    0   &        2         &  0  \\   
  10  &    0  &         2         &  0  \\
  11  &    1   &        2         &  0  \\  
\end{tabular}
\end{center}

\[
  {\rm genus}(\Gamma_2) =13,\ \ \ \ \
  {\rm genus}(\Gamma_1) =1,
\]
\[
deg(\phi) =4.
\]

\begin{figure}[H]
\begin{minipage}{\textwidth}
\begin{center}
%\vspace{8.0 cm}
\includegraphics[height=5cm,width=5cm]{graphics/harmonic-quotient_d4reg12d-C3_212002121001b.png}
\end{center}
\begin{center}
%\vspace{8.0 cm}
\includegraphics[height=4cm,width=4cm]{graphics/cycle3-012.png}
\end{center}
\end{minipage}
\caption{The quartic graph $\Gamma_2$ ({\tt d4reg12d}) is a degree four $2$-multicover of $\Gamma_1=C_3$.}
\label{fig:harmonic-quotient_d4reg12d-C3_212002121001}
\end{figure}

All the $20$ non-trivial harmonic quotient morphisms
$\pi: \Gamma_2\to \Gamma_1=\Gamma_2/G$, in the case where
$\Gamma_1$ has $2$ vertices, map to $P_2$.

}
\end{example}



\begin{example}% d4reg12e
 {\tt d4reg12e}\newline
{\rm  
  Let $\Gamma_2$ denote the graph {\tt d4reg12e} as given above.
It has $20$  non-trivial harmonic quotient morphisms
$\pi: \Gamma_2\to \Gamma_1=\Gamma_2/G$,
for some subgroup $G$ of $Aut(\Gamma_2)$.
There are $6$ in the case $\Gamma_1$ has $7$ vertices,
only $1$ in case $\Gamma_1$ has $6$ vertices,
none in case $\Gamma_1$ has $5$ vertices,
$3$ in case $\Gamma_1$ has $4$ vertices,
$3$ in the case $\Gamma_1$ has $3$ vertices,
and  $7$ in the case $\Gamma_1$ has $2$ vertices
(i.e., $\Gamma_1=P_2$).

For instance, when
\[
G = \langle  (1,11)(2,9)(3,10)(5,6)(7,8) \rangle ,
\]
the quotient morphism maps to the graph
$\pi: \Gamma_2\to {\tt FKLkw}=\Gamma_1=\Gamma_2/G $ (in graph6 notation).
All of the $6$ quotients $\Gamma_1=\Gamma_2/G$ having $7$ vertices
are isomorphic to this graph, {\tt FKLkw}.
This example, which is a matchstick graph,
is tabulated below and the uncolored graphs are depicted in
Figure \ref{fig:harmonic-quotient_d4reg12e+FKLkw}.
\index{{\tt FKLkw} graph}

\begin{center}
\begin{tabular}{c|c|c|c}  
  $x$ & $\pi(x)$ & $m_\pi(x)$ & $\nu_\pi(x)$   \\ 
  0   &    5   &        2         &  0  \\ 
  1   &    2   &        1         &  0  \\ 
  2   &    3   &        1         &  0  \\ 
  3   &    4   &        1         &  1   \\
  4   &    0   &        2         &  0   \\
  5   &    1   &        1         &  0   \\
  6   &    1   &        1         &  0   \\
  7   &    6   &        1         &  1   \\
  8   &    6   &        1         &  1   \\
  9   &    3   &        1         &  0  \\   
  10  &    4  &         1         &  1  \\
  11  &    2   &        1         &  0  \\  
\end{tabular}
\end{center}

\[
  {\rm genus}(\Gamma_2) =13,\ \ \ \ \
  {\rm genus}(\Gamma_1) =5,
\]
\[
deg(\phi) =2.
\]

\begin{figure}[H]
\begin{minipage}{\textwidth}
\begin{center}
%\vspace{8.0 cm}
\includegraphics[height=5cm,width=5cm]{graphics/graph_d4reg12e1.png}
\end{center}
\begin{center}
%\vspace{8.0 cm}
\includegraphics[height=4cm,width=4cm]{graphics/FKLkw-quotient-of-d4reg12e.png}
\end{center}
\end{minipage}
\caption{The quartic graph $\Gamma_2$, a degree two harmonic quotient to $\Gamma_1={\tt FKLkw}$.}
\label{fig:harmonic-quotient_d4reg12e+FKLkw}
\end{figure}

The only non-trivial harmonic quotient morphism
$\pi: \Gamma_2\to \Gamma_1=\Gamma_2/G$, in the case where
$\Gamma_1$ has $6$ vertices, satisfies $\Gamma_1={\tt Oct}$,
the quartic octahedral graph.
This quotient is also a cover\footnote{As in the above case of
{\tt d4reg12d}, see Example \ref{example:d4reg12d}.}
\index{octahedral graph}


The $3$ non-trivial harmonic quotient morphisms
$\pi: \Gamma_2\to \Gamma_1=\Gamma_2/G$, in the case where
$\Gamma_1$ has $4$ vertices, satisfy $\Gamma_1={\tt diamond}$,
the diamond graph.
\index{diamond graph}
When
\[
G = \langle 
(1,11)(2,9)(3,10)(5,6)(7,8),
(0,4)(1,5)(3,7)(6,11)(8,10)\rangle ,
\]
the quotient morphism $\pi$ mapping to
$\Gamma_1=\Gamma_2/G ={\tt diamond}$
is tabulated below
and  depicted in Figure \ref{fig:harmonic-quotient_d4reg12e-diamond_201320033130}.

\begin{center}
\begin{tabular}{c|c|c|c}  
  $x$ & $\pi(x)$ & $m_\pi(x)$ & $\nu_\pi(x)$   \\ 
  0   &    2   &        2         &  0  \\ 
  1   &    0   &        1         &  1  \\ 
  2   &    1   &        2         &  0  \\ 
  3   &    3   &        1         &  1   \\
  4   &    2   &        2         &  0   \\
  5   &    0   &        1         &  1   \\
  6   &    0   &        1         &  1   \\
  7   &    3   &        1         &  1   \\
  8   &    3   &        1         &  1   \\
  9   &    1   &        2         &  0  \\   
  10  &    3  &         1         &  1  \\
  11  &    0   &        1         &  1  \\  
\end{tabular}
\end{center}

\[
  {\rm genus}(\Gamma_2) =13,\ \ \ \ \
  {\rm genus}(\Gamma_1) =2,
\]
\[
deg(\phi) =4.
\]

\begin{figure}[H]
\begin{minipage}{\textwidth}
\begin{center}
%\vspace{8.0 cm}
\includegraphics[height=5cm,width=5cm]{graphics/harmonic-quotient_d4reg12e-diamond_201320033130.png}
\end{center}
\begin{center}
%\vspace{8.0 cm}
\includegraphics[height=4cm,width=4cm]{graphics/diamond-0123d.png}
\end{center}
\end{minipage}
\caption{The quartic graph $\Gamma_2$ ({\tt d4reg12e}) is a degree four quotient morphism to $\Gamma_1={\tt diamond}$.}
\label{fig:harmonic-quotient_d4reg12e-diamond_201320033130}
\end{figure}

The $3$ non-trivial harmonic quotient morphisms
$\pi: \Gamma_2\to \Gamma_1=\Gamma_2/G$, in the case where
$\Gamma_1$ has $3$ vertices,
satisfy $\Gamma_1=C_3$, the cycle graph.
These quotients $\pi: \Gamma_2\to C_3=\Gamma_2/G$
are multicovers\footnote{This too was seen in the above case of
{\tt d4reg12d}, see Example \ref{example:d4reg12d}.}
of degree 4 and multiplicity $m=2$.

All the $7$ non-trivial harmonic quotient morphisms
$\pi: \Gamma_2\to \Gamma_1=\Gamma_2/G$, in the case where
$\Gamma_1$ has $2$ vertices, map to $P_2$.
}
\end{example}



\begin{example}% d4reg12f
 {\tt d4reg12f}\newline
{\rm  
  Let $\Gamma_2$ denote the graph {\tt d4reg12f} as given above.
It has $75$  non-trivial harmonic quotient morphisms
$\pi: \Gamma_2\to \Gamma_1=\Gamma_2/G$,
for some subgroup $G$ of $Aut(\Gamma_2)$.
There is exactly $1$ in the case $\Gamma_1$ has $8$ vertices,
none in case $\Gamma_1$ has $7$ vertices,
$16$ in case $\Gamma_1$ has $6$ vertices,
$30$ in case $\Gamma_1$ has $4$ vertices,
$6$ in the case $\Gamma_1$ has $3$ vertices,
and $22$ in the case $\Gamma_1$ has $2$ vertices
(i.e., $\Gamma_1=P_2$).

For instance, when
\[
G = \langle (1,8)(2,4)(3,10)(6,9)  \rangle ,
\]
the quotient morphism maps to the graph
$\pi: \Gamma_2\to $\verb+G?r@xw+$=\Gamma_1=\Gamma_2/G $ (in graph6 notation)
with degree sequence $[4,4,4,4,2,2,2,2]$. A very similar example occured in
Example \ref{example:d4reg12a}, where the quotient morphism
was also a graph homomorphism, so won't be tabulated here.

Of the $16$ non-trivial harmonic quotient morphisms
$\pi: \Gamma_2\to \Gamma_1=\Gamma_2/G$, in the case where
$\Gamma_1$ has $6$ vertices, exactly $8$ of them
satisfy $\Gamma_1=$\verb+EL~o+ (in graph6 notation).
(In each case, these quotients all seem to be quasicovers.)
The remaining $8$ satisfy $\Gamma_1=$\verb+E@JW+.

For instance, when
\[
G = \langle (0,7)(1,3)(2,6)(4,9)(5,11)(8,10) \rangle ,
\]
the quotient morphism maps to the graph
$\pi: \Gamma_2\to $\verb+EL~o+$=\Gamma_1=\Gamma_2/G $ (in graph6 notation).
This is tabulated below and the uncolored graphs are depicted in
Figure \ref{fig:harmonic-quotient_d4reg12f+EL_tilde_o}.
\index{\verb+EL~o+ graph}

\begin{center}
\begin{tabular}{c|c|c|c}  
  $x$ & $\pi(x)$ & $m_\pi(x)$ & $\nu_\pi(x)$   \\ 
  0   &    2   &        1         &  0  \\ 
  1   &    5   &        1         &  1  \\ 
  2   &    3   &        1         &  0  \\ 
  3   &    5   &        1         &  1   \\
  4   &    4   &        1         &  0   \\
  5   &    1   &        1         &  0   \\
  6   &    3   &        1         &  0   \\
  7   &    2   &        1         &  0   \\
  8   &    0   &        1         &  1   \\
  9   &    4   &        1         &  0  \\   
  10  &    0  &         1         &  1  \\
  11  &    1   &        1         &  0  \\  
\end{tabular}
\end{center}

\[
  {\rm genus}(\Gamma_2) =13,\ \ \ \ \
  {\rm genus}(\Gamma_1) =6,
\]
\[
deg(\phi) =2.
\]

\begin{figure}[H]
\begin{minipage}{\textwidth}
\begin{center}
%\vspace{8.0 cm}
\includegraphics[height=5cm,width=5cm]{graphics/graph_d4reg12f.png}
\end{center}
\begin{center}
%\vspace{8.0 cm}
\includegraphics[height=4cm,width=4cm]{graphics/EL_tilde_o-quotient-of-d4reg12f.png}
\end{center}
\end{minipage}
\caption{The quartic graph $\Gamma_2$ ({\tt d4reg12f}), a degree two harmonic quotient to $\Gamma_1={\tt EL~o}$.}
\label{fig:harmonic-quotient_d4reg12f+EL_tilde_o}
\end{figure}

For instance, when
\[
G = \langle (1,2)(4,8)(6,9)(7,11), (1,8)(2,4)(3,10)(6,9) \rangle ,
\]
the quotient morphism maps to the graph
$\pi: \Gamma_2\to $\verb+E@JW+$=\Gamma_1=\Gamma_2/G $ (in graph6 notation). This graph is sometimes called the {\it R graph},
because it can be drawn to look like the letter R.
Further, this quotient morphism is in fact a graph homomorphism
(without vertical edges).
This is tabulated below and the uncolored graphs are depicted in
Figure \ref{fig:harmonic-quotient_d4reg12f+E_at_JW}.
\index{\verb+E@JW+ graph}
\index{``R'' graph}

\begin{center}
\begin{tabular}{c|c|c|c}  
  $x$ & $\pi(x)$ & $m_\pi(x)$ & $\nu_\pi(x)$   \\ 
  0   &    5   &        4         &  0  \\ 
  1   &    0   &        1         &  0  \\ 
  2   &    0   &        1         &  0  \\ 
  3   &    4   &        2         &  0   \\
  4   &    0   &        1         &  0   \\
  5   &    2   &        4         &  0   \\
  6   &    1   &        2         &  0   \\
  7   &    3   &        2         &  0   \\
  8   &    0   &        1         &  0   \\
  9   &    1   &        2         &  0  \\   
  10  &    4  &         2         &  0  \\
  11  &    3   &        2         &  0  \\  
\end{tabular}
\end{center}

\[
  {\rm genus}(\Gamma_2) =13,\ \ \ \ \
  {\rm genus}(\Gamma_1) =1,
\]
\[
deg(\phi) =4.
\]

\begin{figure}[H]
\begin{minipage}{\textwidth}
\begin{center}
%\vspace{8.0 cm}
\includegraphics[height=5cm,width=5cm]{graphics/graph_d4reg12f.png}
\end{center}
\begin{center}
%\vspace{8.0 cm}
\includegraphics[height=4cm,width=4cm]{graphics/E_at_JW-quotient-of-d4reg12f.png}
\end{center}
\end{minipage}
\caption{The quartic graph $\Gamma_2$ ({\tt d4reg12f}), a degree 4 harmonic quotient to the R graph, $\Gamma_1$.}
\label{fig:harmonic-quotient_d4reg12f+E_at_JW}
\end{figure}


Of the $30$ non-trivial harmonic quotient morphisms
$\pi: \Gamma_2\to \Gamma_1=\Gamma_2/G$, in the case where
$\Gamma_1$ has $6$ vertices, exactly $4$ of them
satisfy $\Gamma_1={\tt diamond}$.
The remaining $26$ satisfy $\Gamma_1=P_4$.
In each case, these quotients to the ath graph $P_4$ all seem
to be graph homomorphisms.
\index{diamond graph}


For instance, when
\[
G = \langle (1,8)(2,4)(3,10)(6,9), (0,5)(1,4,8,2)(3,6,10,9)(7,11) \rangle ,
\]
the quotient morphism maps to the graph
$\pi: \Gamma_2\to P_4=\Gamma_1=\Gamma_2/G $.
Further, this quotient morphism is in fact a graph homomorphism
(without vertical edges).
This is tabulated below and depicted in
Figure \ref{fig:harmonic-quotient_d4reg12f-P4_311213201220}.

\begin{center}
\begin{tabular}{c|c|c|c}  
  $x$ & $\pi(x)$ & $m_\pi(x)$ & $\nu_\pi(x)$   \\ 
  0   &    3   &        4         &  0  \\ 
  1   &    1   &        2         &  0  \\ 
  2   &    1   &        2         &  0  \\ 
  3   &    2   &        2         &  0   \\
  4   &    1   &        2         &  0   \\
  5   &    3   &        4         &  0   \\
  6   &    2   &        2         &  0   \\
  7   &    0   &        4         &  0   \\
  8   &    1   &        2         &  0   \\
  9   &    2   &        2         &  0  \\   
  10  &    2  &         2         &  0  \\
  11  &    0   &        4         &  0  \\  
\end{tabular}
\end{center}

\[
  {\rm genus}(\Gamma_2) =13,\ \ \ \ \
  {\rm genus}(\Gamma_1) =0,
\]
\[
deg(\phi) =4.
\]

\begin{figure}[H]
\begin{minipage}{\textwidth}
\begin{center}
%\vspace{8.0 cm}
\includegraphics[height=5cm,width=5cm]{graphics/harmonic-quotient_d4reg12f-P4_311213201220.png}
\end{center}
\begin{center}
%\vspace{8.0 cm}
\includegraphics[height=4cm,width=4cm]{graphics/path4-3120.png}
\end{center}
\end{minipage}
\caption{The quartic graph $\Gamma_2$ ({\tt d4reg12f}), a degree 4 harmonic quotient to the path graph, $\Gamma_1=P_4$.}
\label{fig:harmonic-quotient_d4reg12f-P4_311213201220}
\end{figure}

All the $6$ non-trivial harmonic quotient morphisms
$\pi: \Gamma_2\to \Gamma_1=\Gamma_2/G$, in the case where
$\Gamma_1$ has $3$ vertices, map to $P_3$.

All $22$ of the non-trivial harmonic quotient morphisms
$\pi: \Gamma_2\to \Gamma_1=\Gamma_2/G$, in the case where
$\Gamma_1$ has $2$ vertices, map to $P_2$.
}
\end{example}


  
\section{Quintic cases}

\subsection{$6$ vertices}

Only the complete graph $K_6$ fits into this category.

There are exactly $90$ non-trivial harmonic morphisms from
$K_6$ to $K_3=C_3$. All of these morphisms are quasi-multicovers
of degree $4$, constant horizontal multiplicity $2$ and
constant vertical mutiplicity $1$.
%See Example \ref{example:K6-K3} for an example.
\index{quasi-multicover}


\subsubsection{Examples}

\begin{example}% d5reg6a
  {\tt d5reg6a}\newline
{\rm  
  Let $\Gamma_2=K_6$. All $90$ of the quasi-multicovers
of $C_3$ (with degree $4$, constant horizontal multiplicity $2$,
constant vertical mutiplicity $1$) are obtained as a harmonic
quotient morphism $\pi:\Gamma_2 \to \Gamma_1=\Gamma_2/G$, for some
subgroup $G\subset Aut(\Gamma_2)$.

}
\end{example}

\subsection{$8$ vertices}
\label{sec:d5reg8}

There are exactly three $5$-regular graphs with $8$ vertices, two of which are vertex transitive.
They are Hamiltonian by Ore's Theorem.


\begin{enumerate}
\item
  {\tt d5reg8a}\newline
  Define the graph  $\Gamma$
  % d5reg8-1 = d5reg8a
  to have edge set:

\[
E = \{
(0, 1), (0, 2), (0, 4), (0, 6), (0, 7), (1, 2), (1, 3), (1, 4), (1, 5), (2, 3),
\]
\[
(2, 5), (2, 7), (3, 4), (3, 6), (3, 7), (4, 5), (4, 6), (5, 6), (5, 7), (6, 7)
\}.
\]
It has diameter $2$,
girth $3$, circumference $8$, 
critical group of size $21025 = 5^2 29^2$, 
and automorphism group of size $60$. This has chromatic number $4$.

There are no non-trivial harmonic morphisms from $\Gamma$ to $P_3$, $P_4$,
$C_3$, $C_4$, or to $K_4$.

\item
  {\tt d5reg8b}\newline
Define the graph  $\Gamma$ to have edge set:
% d5reg8-2 = d5reg8b
  

\[
E = \{
(0, 1), (0, 2), (0, 3), (0, 5), (0, 6), (1, 2), (1, 3), (1, 4), (1, 7), (2, 3), 
\]
\[
(2, 4), (2, 7), (3, 5), (3, 6), (4, 5), (4, 6), (4, 7), (5, 6), (5, 7), (6, 7)
\}.
\]
This graph is isomorphic to the strong graph product
$C_4\times_{Str} P_2$. It has diameter $2$,
girth $3$, circumference $8$, 
critical group of size $20736 = 2^8 3^4$, 
and automorphism group of size $128$. This is a vertex transitive graph but not
edge transitive.
This has chromatic number $4$.
\index{vertex transitive graph}

There are $4$ non-trivial harmonic morphisms from $\Gamma$ to $P_3$,
$48$ non-trivial harmonic morphisms from $\Gamma$ to $C_3$, and
$8$ non-trivial harmonic
morphisms\footnote{Moreover, it appears all of these are quasi-multicovers
  (with constant horizontal multiplicity $2$ and constant vertical
  multiplicity $1$). This is consistent with the values predicted by
the known values for the projection morphism $pr_1:C_4\times_{Str}P_2 \to C_4$.}
\index{quasi-multicover}
from $\Gamma$ to $C_4$. However, there are
no non-trivial harmonic morphisms from $\Gamma$ to $P_4$ or to $K_4$.


\item
  {\tt d5reg8c}\newline
  Define the graph  $\Gamma$
  % d5reg8-3 = d5reg8c
  to have edge set:

\[
E = \{
(0, 1), (0, 2), (0, 3), (0, 4), (0, 6), (1, 2), (1, 3), (1, 5), (1, 6), (2, 3),
\]
\[
(2, 5), (2, 7), (3, 4), (3, 7), (4, 5), (4, 6), (4, 7), (5, 6), (5, 7), (6, 7)
\}.
\]
It has diameter $2$,
girth $3$, circumference $8$, 
critical group of size $20808 = 2^3 3^2 17^2$, 
and automorphism group of size $16$. This is a vertex transitive graph but not
edge transitive.
This has chromatic number $4$.
\index{vertex transitive graph}

There are $12$ non-trivial harmonic morphisms from $\Gamma$ to $C_3$.
However, there are no non-trivial harmonic morphisms from $\Gamma$ to $P_3$,
$P_4$, $C_4$, or to $K_4$.


\end{enumerate}



\subsubsection{Examples}

\begin{example}% d5reg8a
  {\tt d5reg8a}\newline
{\rm  
  Let $\Gamma_2$ denote the graph {\tt d5reg8a}, as above.
It has $6$  non-trivial harmonic quotient morphisms
$\pi: \Gamma_2\to \Gamma_1=\Gamma_2/G$,
for some subgroup $G$ of $Aut(\Gamma_2)$.
There are exactly $2$ in the case $\Gamma_1$ has $4$ vertices,
and $4$ in the case $\Gamma_1$ has $2$ vertices
(i.e., $\Gamma_1=P_2$).
Both of the quotient morphisms in case $\Gamma_1$ has $4$ vertices
are to $\Gamma_1=St_4$. 

For instance, when
\[
G = \langle (1,2,7,6,4) \rangle ,
\]
the quotient morphism maps to the graph
$\pi: \Gamma_2\to St_4=\Gamma_1=\Gamma_2/G $
This is tabulated below.

\begin{center}
\begin{tabular}{c|c|c|c}  
  $x$ & $\pi(x)$ & $m_\pi(x)$ & $\nu_\pi(x)$   \\ 
  0   &    2   &        5         &  0  \\ 
  1   &    1   &        1         &  2  \\ 
  2   &    1   &        1         &  2  \\ 
  3   &    0   &        5         &  0   \\
  4   &    1   &        1         &  2   \\
  5   &    3   &        5         &  0   \\
  6   &    1   &        1         &  2   \\
  7   &    1   &        1         &  2   \\
\end{tabular}
\end{center}

\[
  {\rm genus}(\Gamma_2) =13,\ \ \ \ \
  {\rm genus}(\Gamma_1) =0,
\]
\[
deg(\phi) =5.
\]
  
}
\end{example}


\begin{example}% d5reg8b
  \label{example:d5reg8b}
  {\tt d5reg8b}\newline
{\rm  
  Let $\Gamma_2$ denote the graph {\tt d5reg8b}, as above.
It has $304$  non-trivial harmonic quotient morphisms
$\pi: \Gamma_2\to \Gamma_1=\Gamma_2/G$,
for some subgroup $G$ of $Aut(\Gamma_2)$.
There are exactly $26$ in the case $\Gamma_1$ has $4$ vertices,
$116$ in the case $\Gamma_1$ has $3$ vertices
and $162$ in the case $\Gamma_1$ has $2$ vertices
(i.e., $\Gamma_1=P_2$).
All of the quotient morphisms in case $\Gamma_1$ has $4$ vertices
are to $\Gamma_1=C_4$. 
  When
\[
G = \langle (0,3)(1,2)(4,7)(5,6) \rangle ,
\]
the quotient morphism $\pi$ maps to
$\Gamma_1=\Gamma_2/G =C_4$
is tabulated below
and  depicted in Figure \ref{fig:harmonic-quotient_d5reg8b-C4_32231001}.

\begin{center}
\begin{tabular}{c|c|c|c}  
  $x$ & $\pi(x)$ & $m_\pi(x)$ & $\nu_\pi(x)$   \\ 
  0   &    3   &        2         &  1  \\ 
  1   &    2   &        2         &  1  \\ 
  2   &    2   &        2         &  1  \\ 
  3   &    3   &        2         &  1   \\
  4   &    1   &        2         &  1   \\
  5   &    0   &        2         &  1   \\
  6   &    0   &        2         &  1   \\
  7   &    1   &        2         &  1   \\
\end{tabular}
\end{center}

\[
  {\rm genus}(\Gamma_2) =13,\ \ \ \ \
  {\rm genus}(\Gamma_1) =1,
\]
\[
deg(\phi) =2.
\]

\begin{figure}[H]
\begin{minipage}{\textwidth}
\begin{center}
%\vspace{8.0 cm}
\includegraphics[height=5cm,width=5cm]{graphics/harmonic-quotient_d5reg8b-C4_32231001.png}
\end{center}
\begin{center}
%\vspace{8.0 cm}
\includegraphics[height=4cm,width=4cm]{graphics/Cycle4-0123c.png}
\end{center}
\end{minipage}
\caption{The quartic graph $\Gamma_2$ ({\tt d5reg8b}) is a degree two quasi-multicover of $\Gamma_1=C_4$.}
\label{fig:harmonic-quotient_d5reg8b-C4_32231001}
\end{figure}

All of the quotient morphisms in case
$\Gamma_1$ has $3$ vertices are to $\Gamma_1=C_3$. 
However, these are not quasi-multicovers.
For example, when
\[
G = \langle (4,7)(5,6), (0,1)(2,3)(4,6,7,5) \rangle ,
\]
the quotient morphism $\pi$ maps to
$\Gamma_1=\Gamma_2/G =C_3$
is tabulated below
and  depicted in Figure \ref{fig:harmonic-quotient_d5reg8b-C3_00112222}.

\begin{center}
\begin{tabular}{c|c|c|c}  
  $x$ & $\pi(x)$ & $m_\pi(x)$ & $\nu_\pi(x)$   \\ 
  0   &    0   &        2         &  1  \\ 
  1   &    0   &        2         &  1  \\ 
  2   &    1   &        2         &  1  \\ 
  3   &    1   &        2         &  1   \\
  4   &    2   &        1         &  3   \\
  5   &    2   &        1         &  3   \\
  6   &    2   &        1         &  3   \\
  7   &    2   &        1         &  3   \\
\end{tabular}
\end{center}

\[
  {\rm genus}(\Gamma_2) =13,\ \ \ \ \
  {\rm genus}(\Gamma_1) =1,
\]
\[
deg(\phi) =4.
\]

\begin{figure}[H]
\begin{minipage}{\textwidth}
\begin{center}
%\vspace{8.0 cm}
\includegraphics[height=5cm,width=5cm]{graphics/harmonic-quotient_d5reg8b-C3_00112222.png}
\end{center}
\begin{center}
%\vspace{8.0 cm}
\includegraphics[height=4cm,width=4cm]{graphics/cycle3-012.png}
\end{center}
\end{minipage}
\caption{The quartic graph $\Gamma_2$ ({\tt d5reg8b}) has a degree four harmonic quotient morphism to $\Gamma_1=C_3$.}
\label{fig:harmonic-quotient_d5reg8b-C3_00112222}
\end{figure}

}
\end{example}


\begin{example}% d5reg8c
  {\tt d5reg8c}\newline
{\rm  
  Let $\Gamma_2$ denote the graph {\tt d5reg8c}, as above.
It has $6$  non-trivial harmonic quotient morphisms
$\pi: \Gamma_2\to \Gamma_1=\Gamma_2/G$,
for some subgroup $G$ of $Aut(\Gamma_2)$.
There are exactly $2$ in the case $\Gamma_1$ has $3$ vertices,
and $4$ in the case $\Gamma_1$ has $2$ vertices
(i.e., $\Gamma_1=P_2$).
The morphisms to $C_3$ are very similar to those of the
previous Example \ref{example:d5reg8b} and will not be repeated.
  
}
\end{example}




\subsection{$10$ vertices}

Up to isomorphism, there are exactly $60$ $5$-regular graphs with $10$ vertices, four of which are vertex transitive,
% d5reg10-12, -24, -39
and exactly one
% d5reg10-53 
of those which is also edge transitive (and that one is also symmetric, distance regular, and bipartite). By Ore's Theorem, these are all Hamiltonian.



\begin{enumerate}
\item
  {\tt d5reg10a}\newline
  Define the graph  $\Gamma$
  % d5reg10-12 = d5reg10a
to have edge set:

\[
E = \{
(0, 1), (0, 2), (0, 3), (0, 8), (0, 9), (1, 2), (1, 3), (1, 5), (1, 9), (2, 3),
\]
\[
(2, 7), (2, 9), (3, 6), (3, 9), (4, 5), (4, 6), (4, 7), (4, 8), (4, 9), (5, 6),
\]
\[
(5, 7), (5, 8), (6, 7), (6, 8), (7, 8)
\}.
\]
It has diameter $2$,
girth $3$, circumference $10$, 
critical group of size $300125 = 5^3 7^4$, 
and automorphism group of size $240$. This is a vertex transitive graph but not
edge transitive.
This has chromatic number $5$.
\index{vertex transitive graph}

There are no non-trivial harmonic morphisms from $\Gamma$ to $P_3$, $P_4$, $C_3$, $C_4$, or to $K_4$.


\item
  {\tt d5reg10b}\newline
  Define the graph  $\Gamma$
  % d5reg10-24 = d5reg10b
  to have edge set:

\[
E = \{
(0, 1), (0, 6), (0, 7), (0, 8), (0, 9), (1, 3), (1, 4), (1, 5),(1, 7), (2, 3),
\]
\[
(2, 5), (2, 7), (2, 8), (2, 9), (3, 6), (3, 8), (3, 9), (4, 5), (4, 7), (4, 8),
\]
\[
(4, 9), (5, 6), (5, 9), (6, 7), (6, 8)
\}.
\]
It has diameter $2$,
girth $3$, circumference $10$, 
critical group of size $360375 = 3\cdot 5^3 31^2$, 
and automorphism group of size $20$. This is a vertex transitive graph but not
edge transitive.
This has chromatic number $4$.
\index{vertex transitive graph}

There are no non-trivial harmonic morphisms from $\Gamma$ to $P_3$, $P_4$, $C_3$, $C_4$, or to $K_4$.

\item
  {\tt d5reg10c}\newline
  Define the graph $\Gamma$
  % d5reg10-39 = d5reg10c
to have edge set:

\[
E = \{
(0, 1), (0, 2), (0, 7), (0, 8), (0, 9), (1, 2), (1, 3), (1, 5), (1, 9), (2, 3),
\]
\[
(2, 5), (2, 9), (3, 4), (3, 5), (3, 6), (4, 5), (4, 6), (4, 7), (4, 8), (5, 6), \]
\[
(6, 7), (6, 8), (7, 8), (7, 9), (8, 9)
\}.
\]
It has diameter $2$,
girth $3$, circumference $10$, 
critical group of size $311040 = 2^8 3^5 5$, 
and automorphism group of size $320$. This is a vertex transitive graph but not
edge transitive.
This has chromatic number $5$.
\index{vertex transitive graph}

There are $10$ non-trivial harmonic morphisms from $\Gamma$ to $P_3$
and $60$ non-trivial harmonic morphisms from $\Gamma$ to $C_3$.
However, there are
no non-trivial harmonic morphisms from $\Gamma$ to $P_4$, $C_4$, or to $K_4$.


\item
  {\tt d5reg10d}\newline
  Define the graph  $\Gamma$
  % d5reg10-53 = d5reg10d
to have edge set:

\[
E = \{
(0, 1), (0, 2), (0, 5), (0, 6), (0, 8), (1, 3), (1, 4), (1, 7), (1, 9), (2, 3),
\]
\[
(2, 4), (2, 7), (2, 9), (3, 5), (3, 6), (3, 8), (4, 5), (4, 6), (4, 8), (5, 7),
\]
\[
(5, 9), (6, 7), (6, 9), (7, 8), (8, 9)
\}.
\]
It has diameter $2$,
girth $4$, circumference $10$, 
critical group of size $390625 = 5^8$, 
and automorphism group of size $28800$.
This is a vertex transitive, edge transitive,
symmetric (=arc transitive), and distance regular.
This has chromatic number $2$ (bipartite).
It is a circulant graph with parameters $(10, [1, 3, 5, 7, 9])$
and is isomorphic to the complete bipartite graph $K_{5,5}$.
\index{complete bipartite graph}
\index{vertex transitive graph}
\index{edge transitive graph}
\index{arc transitive graph}
\index{bipartite graph}
\index{distance regular graph}

There are no non-trivial harmonic morphisms from $\Gamma$ to
$P_4$, $C_3$, $C_4$, or to $K_4$. However, there
are\footnote{One of these has CL $(0, 1, 1, 1, 1, 1, 1, 1, 1, 2)$,
yielding a weed with three roots.}
$100$ non-trivial harmonic morphisms from $\Gamma$ to $P_3$.

\end{enumerate}


\subsubsection{Examples}

\begin{example}% d5reg10a
\label{example:d5reg10a}
  {\tt d5reg10a}\newline
{\rm  
  Let $\Gamma_2$ denote the graph {\tt d5reg10a}, as above.
It has $151$  non-trivial harmonic quotient morphisms
$\pi: \Gamma_2\to \Gamma_1=\Gamma_2/G$,
for some subgroup $G$ of $Aut(\Gamma_2)$.
There is exactly $1$ in the case $\Gamma_1$ has $5$ vertices
(in which case $\Gamma_1=K_5$),
and $150$ in the case $\Gamma_1$ has $2$ vertices
(i.e., $\Gamma_1=P_2$).
\index{complete graph}

The quotient morphism for the
subgroup\footnote{In this example, the subgroup $G$
is actually the center of $Aut(\Gamma_2)$.}

\[
G=\langle (0,8)(1,5)(2,7)(3,6)(4,9) \rangle,
\]
is tabulated below. The uncolored graphs are depicted in
Figure \ref{fig:harmonic-quotient_d5reg10a+K5}.

\begin{center}
\begin{tabular}{c|c|c|c}  
  $x$ & $\phi(x)$ & $m_\phi(x)$ & $\nu_\phi(x)$ \\ \hline
  0   &    4   &        1         &  1   \\ 
  1   &    0   &        1         &  1  \\ 
  2   &    3   &        1         &  1  \\ 
  3   &    2   &        1         &  1  \\ 
  4   &    1   &        1         &  1   \\ 
  5   &    0   &        1         &  1    \\ 
  6   &    2   &        1         &  1   \\ 
  7   &    3   &        1         &  1   \\
  8   &    4   &        1         &  1   \\
  9   &    1   &        1         &  1   \\
\end{tabular}
\end{center}

\[
  {\rm genus}(\Gamma_2) =16,\ \ \ \ \
  {\rm genus}(\Gamma_1) =6,
\]

\[
deg(\phi) =2.
\]

\begin{figure}[H]
\begin{minipage}{\textwidth}
\begin{center}
%\vspace{8.0 cm}
\includegraphics[height=5cm,width=5cm]{graphics/graph_d5reg10a.png}
\end{center}
\begin{center}
%\vspace{8.0 cm}
\includegraphics[height=4cm,width=4cm]{graphics/K5-quotient-of-d5reg10a.png}
\end{center}
\end{minipage}
\caption{The quintic graph  $\Gamma_2$ ({\tt d5reg10a}) is a degree two quasicover of the complete graph, $\Gamma_1=\Gamma_2/G_2=K_5$.}
\label{fig:harmonic-quotient_d5reg10a+K5}
\end{figure}

In this example, $G$ is the center of $Aut(\Gamma_2)$ and
$Aut(\Gamma_2)/G=Aut(\Gamma_1)=Aut(K_5)=S_5$.

\noindent
{\bf Question}:
Is the quotient

\[
\pi: \Gamma_2 \to K_5=\Gamma_1=\Gamma_2/G
\]
an $S_5$-equivariant morphism (that is to say, the morphism commutes with the
group actions)?


}
\end{example}


\begin{example}% d5reg10b
  {\tt d5reg10b}\newline
{\rm  
  Let $\Gamma_2$ denote the graph {\tt d5reg10b}, as above.
It has $3$  non-trivial harmonic quotient morphisms
$\pi: \Gamma_2\to \Gamma_1=\Gamma_2/G$,
for some subgroup $G$ of $Aut(\Gamma_2)$.
There is exactly $1$ in the case $\Gamma_1$ has $5$ vertices
(in which case $\Gamma_1=K_5$),
and $2$ in the case $\Gamma_1$ has $2$ vertices
(i.e., $\Gamma_1=P_2$).
\index{complete graph}

This example is similar in some ways to {\tt d5reg10a}, except that
it's automorphism group is much smaller. For example, in this
example, $K_5$ is also the quotient of $\Gamma_2$ by the
center of $Aut(\Gamma_2)$. Moreover, the quotient morphism
in this case is also a degree two quasicover
where each edge has vertical multiplicity $1$.


}
\end{example}


\begin{example}% d5reg10c
  {\tt d5reg10c}\newline
{\rm  
  Let $\Gamma_2$ denote the graph {\tt d5reg10c}, as above.
It has $730$  non-trivial harmonic quotient morphisms
$\pi: \Gamma_2\to \Gamma_1=\Gamma_2/G$,
for some subgroup $G$ of $Aut(\Gamma_2)$.
There are $158$ in the case $\Gamma_1$ has $5$ vertices
(all of these have $\Gamma_1=C_5$),
$540$ in the case $\Gamma_1$ has $3$ vertices
(all of these have $\Gamma_1=P_3$),
and $32$ in the case $\Gamma_1$ has $2$ vertices
(i.e., $\Gamma_1=P_2$).

The quotient morphism for the
subgroup\footnote{In this example, the subgroup $G$
is actually the center of $Aut(\Gamma_2)$.}

\[
G=\langle (0,9)(1,2)(3,5)(4,6)(7,8) \rangle,
\]
is tabulated below. The uncolored graphs are depicted in
Figure \ref{fig:harmonic-quotient_d5reg10c+C5}.

\begin{center}
\begin{tabular}{c|c|c|c}  
  $x$ & $\phi(x)$ & $m_\phi(x)$ & $\nu_\phi(x)$ \\ \hline
  0   &    3   &        2         &  1   \\ 
  1   &    2   &        2         &  1  \\ 
  2   &    2   &        2         &  1  \\ 
  3   &    1   &        2         &  1  \\ 
  4   &    0   &        2         &  1   \\ 
  5   &    1   &        2         &  1    \\ 
  6   &    0   &        2         &  1   \\ 
  7   &    4   &        2         &  1   \\
  8   &    4   &        2         &  1   \\
  9   &    3   &        2         &  1   \\
\end{tabular}
\end{center}

\[
  {\rm genus}(\Gamma_2) =16,\ \ \ \ \
  {\rm genus}(\Gamma_1) =1,
\]

\[
deg(\phi) =2.
\]

\begin{figure}[H]
\begin{minipage}{\textwidth}
\begin{center}
%\vspace{8.0 cm}
\includegraphics[height=5cm,width=5cm]{graphics/graph_d5reg10c.png}
\end{center}
\begin{center}
%\vspace{8.0 cm}
\includegraphics[height=4cm,width=4cm]{graphics/C5-quotient-of-d5reg10c.png}
\end{center}
\end{minipage}
\caption{The quintic graph  $\Gamma_2$ ({\tt d5reg10c}) is a degree two quasi-$2$-multicover of the cycle graph, $\Gamma_1=\Gamma_2/G_2=C_5$.}
\label{fig:harmonic-quotient_d5reg10c+C5}
\end{figure}

The quotient morphism for the
subgroup\footnote{In this example, the subgroup $G$
contains the center of $Aut(\Gamma_2)$ as a subgroup of index $2$.}

\[
G=\langle (1,7)(2,8)(3,4)(5,6), (0,9)(1,2)(3,5)(4,6)(7,8) \rangle,
\]
is tabulated below. 

\begin{center}
\begin{tabular}{c|c|c|c}  
  $x$ & $\phi(x)$ & $m_\phi(x)$ & $\nu_\phi(x)$ \\ \hline
  0   &    2   &        4         &  1   \\ 
  1   &    1   &        2         &  1  \\ 
  2   &    1   &        2         &  1  \\ 
  3   &    0   &        2         &  3  \\ 
  4   &    0   &        2         &  3   \\ 
  5   &    0   &        2         &  3    \\ 
  6   &    0   &        2         &  3   \\ 
  7   &    1   &        2         &  1   \\
  8   &    1   &        2         &  1   \\
  9   &    2   &        4         &  1   \\
\end{tabular}
\end{center}

\[
  {\rm genus}(\Gamma_2) =16,\ \ \ \ \
  {\rm genus}(\Gamma_1) =0,
\]

\[
deg(\phi) =4.
\]


}
\end{example}


\begin{example}% d5reg10d
  {\tt d5reg10d}\newline
{\rm  
  Let $\Gamma_2$ denote the graph {\tt d5reg10d}, as above.
It has $28425$  non-trivial harmonic quotient morphisms
$\pi: \Gamma_2\to \Gamma_1=\Gamma_2/G$,
for some subgroup $G$ of $Aut(\Gamma_2)$.
There are $40$ in the case $\Gamma_1$ has $6$ vertices
(all of these have $\Gamma_1=St_6$),
$120$ in the case $\Gamma_1$ has $5$ vertices
(all of these have $\Gamma_1=K_5$),
and $28265$ in the case $\Gamma_1$ has $2$ vertices
(i.e., $\Gamma_1=P_2$).

The quotient morphism for the
subgroup\footnote{In this example, the subgroup $G$
is actually the center of $Aut(\Gamma_2)$.}

\[
G=\langle (1,5,8,2,6) \rangle,
\]
is tabulated below. The uncolored graphs are depicted in
Figure \ref{fig:harmonic-quotient_d5reg10d+St6}.

\begin{center}
\begin{tabular}{c|c|c|c}  
  $x$ & $\phi(x)$ & $m_\phi(x)$ & $\nu_\phi(x)$ \\ \hline
  0   &    4   &        5         &  0   \\ 
  1   &    1   &        1         &  0  \\ 
  2   &    1   &        1         &  0  \\ 
  3   &    0   &        5         &  0  \\ 
  4   &    3   &        5         &  0   \\ 
  5   &    1   &        1         &  0    \\ 
  6   &    1   &        1         &  0   \\ 
  7   &    5   &        5         &  0   \\
  8   &    1   &        1         &  0   \\
  9   &    2   &        5         &  0   \\
\end{tabular}
\end{center}

\[
  {\rm genus}(\Gamma_2) =16,\ \ \ \ \
  {\rm genus}(\Gamma_1) =0,
\]

\[
deg(\phi) =5.
\]

\begin{figure}[H]
\begin{minipage}{\textwidth}
\begin{center}
%\vspace{8.0 cm}
\includegraphics[height=5cm,width=5cm]{graphics/graph_d5reg10d.png}
\end{center}
\begin{center}
%\vspace{8.0 cm}
\includegraphics[height=4cm,width=4cm]{graphics/St6-quotient-of-d5reg10d.png}
\end{center}
\end{minipage}
\caption{The quintic graph  $\Gamma_2$ ({\tt d5reg10d}) is a degree five quotient homomorphism to the star graph, $\Gamma_1=\Gamma_2/G_2=St_6$.}
\label{fig:harmonic-quotient_d5reg10d+St6}
\end{figure}

This example is also similar in some ways to {\tt d5reg10a}
in Example \ref{example:d5reg10a}, except that
it's automorphism group is much larger. In this
example, $K_5$ is also the quotient of $\Gamma_2$ by
subgroup of $Aut(\Gamma_2)$ or order $2$. Moreover, the quotient morphism
in this case is also a degree two quasicover
where each edge has vertical multiplicity $1$.
}
\end{example}



\begin{thebibliography}{99}

\bibitem[BC76]{BC76}
  F. Bussemaker, D. Cvetkovi\'c,
  {\it There are exactly 13 connected,
    cubic integral graphs},
  Univ. Beograd. Publ. Elektrotehn Fak. (1976)43-48.
  
\bibitem[BN09]{BN09}  
M. Baker and S. Norine, 
{\it Harmonic morphisms and hyperelliptic graphs},
Int. Math. Res. Not. IMRN (2009) 2914-2955.

\bibitem[graph6]{graph6} B. McKay, graph formats,
  \url{https://users.cecs.anu.edu.au/~bdm/data/formats.html}
  
\bibitem[ISGCI]{ISGCI} Information System on Graph Classes and their Inclusions,
  \url{https://www.graphclasses.org/faq.html}
  
\bibitem[JM17]{JM17} Joyner, D. and  Melles, C. 
{\bf Adventures in graph theory}, Birkha\"user, Basel 2017.

\bibitem[MJ22]{MJ22} C. Melles and D. Joyner, 
  {\bf Hitchhiker's handbook to harmonic morphisms},
  in preparation, 2022.
  
\bibitem[Me17]{Me17} 
A.D. Mednykh and I.A. Mednykh,
{\it The number of spanning trees in circulant graphs, its arithmetic properties and asymptotic}, preprint, 2017. Available at:
\newline
\url{https://arxiv.org/abs/1711.00175}


  
   %\bibitem[]{}

   %\bibitem[]{}

   %\bibitem[]{} 
\end{thebibliography}

\printindex

\end{document}





sage: V2 = [0,1,2,3,4,5,6,7,8,9,10]
sage: E2 = [(0, 1), (0, 2), (0, 4), (0, 10), (1, 2), (1, 3), (1, 9), (2, 3), (2, 9), (3
....: , 4),(3, 10), (4, 5), (4, 6), (5, 7), (5, 8), (5, 10), (6, 7), (6, 8), (6, 10),(7
....: , 8), (7, 9), (8, 9)]
sage: Gamma2 = Graph([V2, E2]) # d4reg11f=d4reg11-13 
sage: G2 = Gamma2.automorphism_group(); G2.cardinality()
64
sage: L = has_harmonic_quotient(Gamma2, verbose=True)
30  non-trivial harmonic quotient morphism(s) found.

sage: for i in range(30):
....:     print(i, L[i][0], L[i][1].cardinality(), len(L[i][2].vertices()))
....: 
0 [6, 2, 1, 4, 5, 6, 4, 2, 1, 3, 0] 2 7
1 [6, 1, 4, 3, 5, 6, 3, 4, 1, 2, 0] 2 7
2 [3, 2, 1, 4, 6, 4, 3, 2, 1, 5, 0] 2 7
3 [2, 1, 5, 3, 6, 3, 2, 5, 1, 4, 0] 2 7
4 [7, 1, 1, 4, 0, 6, 5, 3, 3, 2, 0] 2 8
5 [4, 3, 3, 2, 0, 4, 2, 3, 3, 1, 0] 4 5
6 [4, 3, 3, 2, 0, 4, 2, 3, 3, 1, 0] 4 5
7 [1, 4, 4, 2, 0, 2, 1, 4, 4, 3, 0] 4 5
8 [1, 4, 4, 2, 0, 2, 1, 4, 4, 3, 0] 4 5
9 [4, 3, 3, 2, 0, 4, 2, 3, 3, 1, 0] 4 5
10 [1, 4, 4, 2, 0, 2, 1, 4, 4, 3, 0] 4 5
11 [4, 3, 3, 2, 0, 4, 2, 3, 3, 1, 0] 4 5
12 [1, 4, 4, 2, 0, 2, 1, 4, 4, 3, 0] 4 5
13 [7, 1, 1, 4, 0, 6, 5, 3, 3, 2, 0] 4 8
14 [7, 1, 1, 4, 0, 6, 5, 3, 3, 2, 0] 4 8
15 [7, 1, 1, 4, 0, 6, 5, 3, 3, 2, 0] 4 8
16 [7, 1, 1, 4, 0, 6, 5, 3, 3, 2, 0] 4 8
17 [7, 1, 1, 4, 0, 6, 5, 3, 3, 2, 0] 8 8
18 [4, 3, 3, 2, 0, 4, 2, 3, 3, 1, 0] 8 5
19 [1, 4, 4, 2, 0, 2, 1, 4, 4, 3, 0] 8 5
20 [4, 3, 3, 2, 0, 4, 2, 3, 3, 1, 0] 8 5
21 [1, 4, 4, 2, 0, 2, 1, 4, 4, 3, 0] 8 5
22 [4, 3, 3, 2, 0, 4, 2, 3, 3, 1, 0] 8 5
23 [1, 4, 4, 2, 0, 2, 1, 4, 4, 3, 0] 8 5
24 [4, 3, 3, 2, 0, 4, 2, 3, 3, 1, 0] 8 5
25 [1, 4, 4, 2, 0, 2, 1, 4, 4, 3, 0] 8 5
26 [4, 3, 3, 2, 0, 4, 2, 3, 3, 1, 0] 8 5
27 [1, 4, 4, 2, 0, 2, 1, 4, 4, 3, 0] 8 5
28 [4, 3, 3, 2, 0, 4, 2, 3, 3, 1, 0] 16 5
29 [1, 4, 4, 2, 0, 2, 1, 4, 4, 3, 0] 16 5
sage: f0 = L[0][0]; f0; Gamma1 = L[0][2]
[6, 2, 1, 4, 5, 6, 4, 2, 1, 3, 0]
sage: Gamma1.show(layout="circular")
Launched png viewer for Graphics object consisting of 19 graphics primitives
sage: Gamma1.show()
Launched png viewer for Graphics object consisting of 19 graphics primitives
sage: graph_id_graph6(Gamma1)
'FCdzo'
sage: Gamma = Graph('FCdzo')
sage: Gamma.is_isomorphic(Gamma1)
True
sage: FCdzo = Gamma
sage: graph_id_graph6(FCdzo)
'FCdzo'
sage: for x in L:
....:     f0 = x[0]; Gamma1 = x[2]
....:     if Gamma1.is_isomorphic(FCdzo):
....:         print(L.index(x), f0," is isom to FCdzo")
....: 
....: 
0 [6, 2, 1, 4, 5, 6, 4, 2, 1, 3, 0]  is isom to FCdzo
1 [6, 1, 4, 3, 5, 6, 3, 4, 1, 2, 0]  is isom to FCdzo
2 [3, 2, 1, 4, 6, 4, 3, 2, 1, 5, 0]  is isom to FCdzo
3 [2, 1, 5, 3, 6, 3, 2, 5, 1, 4, 0]  is isom to FCdzo
sage: 
sage: G0 = L[0][1]; G0; G0.cardinality()
Subgroup generated by [(0,5)(1,7)(2,8)(3,6)] of (Permutation Group with generators [(7,8), (5,6), (4,10), (1,2), (0,3), (0,5)(1,7)(2,8)(3,6)])
2
sage: 
sage: f0 = L[0][0]; f0; Gamma1 = L[0][2]
[6, 2, 1, 4, 5, 6, 4, 2, 1, 3, 0]
sage: f = [V2,f0]
sage: [horizontal_multiplicity_of_harmonic_map(Gamma1, Gamma2, f, v) for v in V2]
[1, 1, 1, 1, 2, 1, 1, 1, 1, 2, 2]
sage: [vertical_multiplicity_of_harmonic_map(Gamma1, Gamma2, f, v) for v in V2]
[0, 0, 0, 0, 0, 0, 0, 0, 0, 0, 0]
sage: genus(Gamma1)
5
sage: genus(Gamma2)
12
sage: 
sage: f0 = L[4][0]; f0; Gamma1 = L[4][2]
[7, 1, 1, 4, 0, 6, 5, 3, 3, 2, 0]
sage: graph_id_graph6(Gamma1, verbose=True)
None found
sage: # Gamma1 = C6+C4+C4
sage: C6_C4_C4 = Gamma1
sage: for x in L:
....:     f0 = x[0]; Gamma1 = x[2]
....:     if Gamma1.is_isomorphic(C6_C4_C4):
....:         print(L.index(x), f0," is isom to C6+C4+C4")
....: 
....: 
4 [7, 1, 1, 4, 0, 6, 5, 3, 3, 2, 0]  is isom to C6+C4+C4
13 [7, 1, 1, 4, 0, 6, 5, 3, 3, 2, 0]  is isom to C6+C4+C4
14 [7, 1, 1, 4, 0, 6, 5, 3, 3, 2, 0]  is isom to C6+C4+C4
15 [7, 1, 1, 4, 0, 6, 5, 3, 3, 2, 0]  is isom to C6+C4+C4
16 [7, 1, 1, 4, 0, 6, 5, 3, 3, 2, 0]  is isom to C6+C4+C4
17 [7, 1, 1, 4, 0, 6, 5, 3, 3, 2, 0]  is isom to C6+C4+C4
sage: f = [V2,f0]
sage: G0 = L[4][1]; G0; G0.cardinality()
Subgroup generated by [(1,2)(4,10)(7,8)] of (Permutation Group with generators [(7,8), (5,6), (4,10), (1,2), (0,3), (0,5)(1,7)(2,8)(3,6)])
2
sage: [horizontal_multiplicity_of_harmonic_map(Gamma1, Gamma2, f, v) for v in V2]
[2, 1, 1, 2, 2, 2, 2, 1, 1, 4, 2]
sage: [vertical_multiplicity_of_harmonic_map(Gamma1, Gamma2, f, v) for v in V2]
[0, 1, 1, 0, 0, 0, 0, 1, 1, 0, 0]
sage: 
sage: f0 = L[5][0]; f0; Gamma1 = L[5][2]
[4, 3, 3, 2, 0, 4, 2, 3, 3, 1, 0]
sage: G0 = L[5][1]; G0; G0.cardinality()
Subgroup generated by [(1,2)(4,10)(7,8), (0,5)(1,7)(2,8)(3,6)] of (Permutation Group with generators [(7,8), (5,6), (4,10), (1,2), (0,3), (0,5)(1,7)(2,8)(3,6)])
4
sage: Gamma1.show()
Launched png viewer for Graphics object consisting of 11 graphics primitives
sage: # Gamma1 = banner
sage: Gamma = Graph("DrG")
sage: Gamma1.is_isomorphic(Gamma)
True
sage: # Gamma1 = DrG
sage: f = [V2,f0]
sage: [horizontal_multiplicity_of_harmonic_map(Gamma1, Gamma2, f, v) for v in V2]
[2, 1, 1, 2, 2, 2, 2, 1, 1, 4, 2]
sage: [vertical_multiplicity_of_harmonic_map(Gamma1, Gamma2, f, v) for v in V2]
[0, 1, 1, 0, 0, 0, 0, 1, 1, 0, 0]
sage: ################# same as for L[4]!!
sage: DrG = Gamma
sage: for x in L:
....:     f0 = x[0]; Gamma1 = x[2]
....:     if Gamma1.is_isomorphic(DrG):
....:         print(L.index(x), f0," is isom to DrG=banner")
....: 
....: 
5 [4, 3, 3, 2, 0, 4, 2, 3, 3, 1, 0]  is isom to DrG=banner
6 [4, 3, 3, 2, 0, 4, 2, 3, 3, 1, 0]  is isom to DrG=banner
7 [1, 4, 4, 2, 0, 2, 1, 4, 4, 3, 0]  is isom to DrG=banner
8 [1, 4, 4, 2, 0, 2, 1, 4, 4, 3, 0]  is isom to DrG=banner
9 [4, 3, 3, 2, 0, 4, 2, 3, 3, 1, 0]  is isom to DrG=banner
10 [1, 4, 4, 2, 0, 2, 1, 4, 4, 3, 0]  is isom to DrG=banner
11 [4, 3, 3, 2, 0, 4, 2, 3, 3, 1, 0]  is isom to DrG=banner
12 [1, 4, 4, 2, 0, 2, 1, 4, 4, 3, 0]  is isom to DrG=banner
18 [4, 3, 3, 2, 0, 4, 2, 3, 3, 1, 0]  is isom to DrG=banner
19 [1, 4, 4, 2, 0, 2, 1, 4, 4, 3, 0]  is isom to DrG=banner
20 [4, 3, 3, 2, 0, 4, 2, 3, 3, 1, 0]  is isom to DrG=banner
21 [1, 4, 4, 2, 0, 2, 1, 4, 4, 3, 0]  is isom to DrG=banner
22 [4, 3, 3, 2, 0, 4, 2, 3, 3, 1, 0]  is isom to DrG=banner
23 [1, 4, 4, 2, 0, 2, 1, 4, 4, 3, 0]  is isom to DrG=banner
24 [4, 3, 3, 2, 0, 4, 2, 3, 3, 1, 0]  is isom to DrG=banner
25 [1, 4, 4, 2, 0, 2, 1, 4, 4, 3, 0]  is isom to DrG=banner
26 [4, 3, 3, 2, 0, 4, 2, 3, 3, 1, 0]  is isom to DrG=banner
27 [1, 4, 4, 2, 0, 2, 1, 4, 4, 3, 0]  is isom to DrG=banner
28 [4, 3, 3, 2, 0, 4, 2, 3, 3, 1, 0]  is isom to DrG=banner
29 [1, 4, 4, 2, 0, 2, 1, 4, 4, 3, 0]  is isom to DrG=banner
sage: 
